% Les catastrophes : classification, gestion et prévention — SVTEEHB Tle TI — Cours signé OnBuch+
% Programme officiel MINESEC. Compiler avec : tectonic N09-les-catastrophes.tex
\newcommand{\DOCMATIERE}{SVTEEHB}
\newcommand{\DOCNIVEAU}{Tle TI}
\newcommand{\DOCMODULE}{Module III : L'Éducation à l'Environnement et au Développement Durable}
\newcommand{\DOCLECON}{N09}
\newcommand{\DOCTITRE}{Les catastrophes : classification, gestion et prévention}
% OnBuch+ — préambule « cours signés » ENRICHI (compilé avec tectonic / XeLaTeX)
% Système visuel « Pop » d'OnBuch+ : fond crème (jamais blanc), bordures encre
% épaisses, ombres « dures » décalées, titres Archivo Black en capitales.
% Version MATHÉMATIQUES : graphes (pgfplots/tikz), boîtes pédagogiques
% (définition, propriété, méthode, attention, le savais-tu, exercice résolu,
% exercice, TP, bilan), page de garde et signature OnBuch+ ancrée en bas.
%
% Macros attendues dans main.tex AVANT \input{preamble} :
%   \DOCMATIERE  \DOCNIVEAU  \DOCMODULE  \DOCLECON (numéro)  \DOCTITRE
\documentclass[11pt]{article}

\usepackage{fontspec}
\usepackage[french]{babel}
\usepackage[a4paper,margin=2cm,top=3.4cm,bottom=2.8cm,headheight=1.4cm,footskip=1.3cm]{geometry}
\usepackage{amsmath,amssymb,amsfonts}
\usepackage{mathtools} % \xmapsto, \coloneqq... souvent utilisés par les modèles
\usepackage{cancel} % simplifications barrées (bilans de forces)
% \abs{z} (module, valeur absolue) et \norm{u} (norme), utilisés par les modèles
\providecommand{\abs}{}\let\abs\relax\DeclarePairedDelimiter{\abs}{\lvert}{\rvert}
\providecommand{\norm}{}\let\norm\relax\DeclarePairedDelimiter{\norm}{\lVert}{\rVert}
\usepackage[table]{xcolor} % [table] : fournit aussi \rowcolors (alternance de lignes)
\usepackage{pagecolor}
\usepackage{tikz}
\usetikzlibrary{arrows.meta,positioning,calc,shapes.geometric,shapes.misc,shapes.arrows,decorations.pathmorphing,decorations.pathreplacing,decorations.markings,patterns,shadows,mindmap,trees,fit,backgrounds,matrix,intersections,angles,quotes}
\usepackage{pgfplots}
\usepackage[version=4]{mhchem} % équations nucléaires \ce{^{235}_{92}U}
\usepackage[european,siunitx]{circuitikz} % schémas électriques
\pgfplotsset{compat=1.18}
\usepgfplotslibrary{fillbetween,groupplots} % aires entre courbes (calcul intégral)
\usepackage{enumitem}
\usepackage{fancyhdr}
% [most] charge toutes les bibliothèques tcolorbox (listings, minted, raster,
% documentation...) et gonfle inutilement la mémoire TeX sur un document long
% (~300 boîtes) : on ne charge que ce qui est réellement utilisé.
\usepackage[skins,breakable]{tcolorbox}
\usepackage{graphicx}
\usepackage{colortbl}
\usepackage{booktabs}
\usepackage{tabularx}
\usepackage{multirow}
\usepackage{diagbox} % case d'en-tête diagonale des tableaux à double entrée
% Colonnes extensibles centrée / gauche / droite (C, L, R), souvent utilisées par les modèles
\newcolumntype{C}{>{\centering\arraybackslash}X}
\newcolumntype{L}{>{\raggedright\arraybackslash}X}
\newcolumntype{R}{>{\raggedleft\arraybackslash}X}
\usepackage{array}
\usepackage{multicol}
\usepackage{siunitx}
% parse-numbers=false : accepte aussi \SI{2,5 \times 10^{-3}}{...}
\sisetup{locale=FR,output-decimal-marker={,},group-minimum-digits=4,parse-numbers=false}
\DeclareSIUnit\ohm{\ensuremath{\symup{\Omega}}} % U+2126 absent de la police
\DeclareSIUnit\division{div} % réglages d'oscilloscope (ms/div, V/div)
\DeclareSIUnit\tr{tr} % tours (vitesse de rotation, ex. \SI{1500}{\tr\per\minute})
\DeclareSIUnit\molar{M} % concentration molaire (ex. \SI{3}{\molar}, \SI{1}{\micro\molar})
\DeclareSIUnit\an{an} % année (le modèle confond parfois avec \year, un compteur TeX) — ex. \SI{5}{\kilo\meter\per\an}
\DeclareSIUnit\FCFA{FCFA} % franc CFA (ex. \SI{500}{\FCFA\per\kilo\gram})
\DeclareSIUnit\degreeBrix{\textdegree Brix} % teneur en sucre d'un sirop/jus (réfractométrie)
\DeclareSIUnit\month{mois} % le modèle utilise parfois \month, un compteur TeX, comme unité de durée
\DeclareSIUnit\microliter{\micro\liter} % le modèle écrit parfois \microliter en un seul token
\DeclareSIUnit\million{millions} % dénombrement (ex. \SI{15}{\million\per\mL} spermatozoïdes)
\usepackage{qrcode}
% \degree : commande fréquemment utilisée par les modèles pour un angle ou une
% température en dehors de \SI{}{} (non fournie par les paquets chargés).
\providecommand{\degree}{^{\circ}}
% \qed (fin de démonstration), utilisé par les modèles sans amsthm
\providecommand{\qed}{\hfill$\square$}
% \barre : double barre des valeurs interdites dans un tableau de variations (tkz-tab)
\providecommand{\barre}{\ensuremath{\|}}
% environnement proof (amsthm non chargé) utilisé par les modèles
\ifdefined\proof\else\newenvironment{proof}[1][Démonstration]{\par\noindent\textit{#1.}\ }{\qed\par}\fi
\usepackage{lastpage}
\usepackage{hyperref}
\hypersetup{hidelinks}

% ============================================================
% Palette « Pop » OnBuch+ — mêmes hex que la classe P (plus_ui.dart)
% ============================================================
\definecolor{poporange}{HTML}{F59321}
\definecolor{popdark}{HTML}{E07A0C}
\definecolor{popcream}{HTML}{FFF6E8}
\definecolor{popink}{HTML}{1C1714}
\definecolor{poppurple}{HTML}{7A5AE0}
\definecolor{popgreen}{HTML}{2E9E6A}
\definecolor{poppink}{HTML}{E0457B}
\definecolor{popblue}{HTML}{2F7DE1}
\definecolor{popgold}{HTML}{C98A12}
\definecolor{popmuted}{HTML}{8A7F76}
\definecolor{popline}{HTML}{EADFD2}
\definecolor{popgray}{HTML}{8A7F76}
\colorlet{popgrey}{popgray}
% Alias tolérants (le modèle nomme parfois les couleurs différemment)
\colorlet{popwhite}{white}
\colorlet{popblack}{popink}
\colorlet{popred}{poppink}
\colorlet{popyellow}{popgold}
% Teintes claires (fonds de tableaux/encadrés) — le modèle nomme parfois
% les couleurs avec un suffixe L (ex. popblueL) sans les avoir définies.
\colorlet{poporangeL}{poporange!20}
\colorlet{popdarkL}{popdark!20}
\colorlet{poppurpleL}{poppurple!20}
\colorlet{popgreenL}{popgreen!20}
\colorlet{poppinkL}{poppink!20}
\colorlet{popblueL}{popblue!20}
\colorlet{popgoldL}{popgold!20}
\colorlet{popredL}{popred!20}
\colorlet{popgrayL}{popgray!20}
% Teintes claires pour fonds de tableaux / zones
\colorlet{popblueL}{popblue!10!white}
\colorlet{poporangeL}{poporange!14!white}
\colorlet{popgreenL}{popgreen!12!white}
\colorlet{poppurpleL}{poppurple!12!white}
\colorlet{poppinkL}{poppink!12!white}
\colorlet{popgoldL}{popgold!14!white}

\pagecolor{popcream}

% ============================================================
% Typographie
% ============================================================
\defaultfontfeatures{Ligatures=TeX}
\setmainfont{PlusJakartaSans-Medium.ttf}[Path=../../../fonts/,BoldFont=PlusJakartaSans-Bold.ttf]
% Maths en sans-serif (Fira Math) avec chiffres et lettres droites de Plus
% Jakarta, pour que les formules se fondent dans le texte.
\usepackage[math-style=ISO,bold-style=ISO]{unicode-math}
\setmathfont{FiraMath-Regular.otf}[Path=../../../fonts/]
\setmathfont{PlusJakartaSans-Medium.ttf}[Path=../../../fonts/,range={up/num,up/{Latin,latin},\mathup}]
\usepackage{icomma} % virgule décimale sans espace en mode math
% Caractères mathématiques Unicode absents des polices de texte (Plus Jakarta,
% Archivo Black) : rendus par leur équivalent mathématique, en texte comme en
% formule — sinon ils s'affichent en carré vide (ex. « Arithmétique dans ℤ »).
\usepackage{newunicodechar}
\newunicodechar{ℝ}{\ensuremath{\mathbb{R}}}
\newunicodechar{ℤ}{\ensuremath{\mathbb{Z}}}
\newunicodechar{ℂ}{\ensuremath{\mathbb{C}}}
\newunicodechar{ℕ}{\ensuremath{\mathbb{N}}}
\newunicodechar{ℚ}{\ensuremath{\mathbb{Q}}}
\newunicodechar{𝔻}{\ensuremath{\mathbb{D}}}
\newunicodechar{α}{\ensuremath{\alpha}}
\newunicodechar{β}{\ensuremath{\beta}}
\newunicodechar{γ}{\ensuremath{\gamma}}
\newunicodechar{δ}{\ensuremath{\delta}}
\newunicodechar{ε}{\ensuremath{\varepsilon}}
\newunicodechar{θ}{\ensuremath{\theta}}
\newunicodechar{λ}{\ensuremath{\lambda}}
\newunicodechar{μ}{\ensuremath{\mu}}
\newunicodechar{σ}{\ensuremath{\sigma}}
\newunicodechar{τ}{\ensuremath{\tau}}
\newunicodechar{φ}{\ensuremath{\varphi}}
\newunicodechar{ψ}{\ensuremath{\psi}}
\newunicodechar{ω}{\ensuremath{\omega}}
\newunicodechar{Σ}{\ensuremath{\Sigma}}
% majuscules produites par \MakeUppercase dans les titres (α-aminés -> Α)
\newunicodechar{Α}{\ensuremath{\alpha}}
\newunicodechar{Β}{\ensuremath{\beta}}
\newunicodechar{Γ}{\ensuremath{\Gamma}}
\newunicodechar{Δ}{\ensuremath{\Delta}}
\newunicodechar{Ε}{\ensuremath{\varepsilon}}
\newunicodechar{Θ}{\ensuremath{\Theta}}
\newunicodechar{Λ}{\ensuremath{\Lambda}}
\newunicodechar{Μ}{\ensuremath{\mu}}
\newunicodechar{Τ}{\ensuremath{\tau}}
\newunicodechar{Φ}{\ensuremath{\Phi}}
\newunicodechar{Ψ}{\ensuremath{\Psi}}
\newunicodechar{Ω}{\ensuremath{\Omega}}
\newunicodechar{↦}{\ensuremath{\mapsto}}
\newunicodechar{∈}{\ensuremath{\in}}
\newunicodechar{∉}{\ensuremath{\notin}}
\newunicodechar{∧}{\ensuremath{\wedge}}
\newunicodechar{⇔}{\ensuremath{\Leftrightarrow}}
\newunicodechar{⟺}{\ensuremath{\Longleftrightarrow}}
\newunicodechar{⇒}{\ensuremath{\Rightarrow}}
\newunicodechar{⟹}{\ensuremath{\Longrightarrow}}
\newunicodechar{∘}{\ensuremath{\circ}}
\newunicodechar{∪}{\ensuremath{\cup}}
\newunicodechar{∩}{\ensuremath{\cap}}
\newunicodechar{≡}{\ensuremath{\equiv}}
\newunicodechar{⊂}{\ensuremath{\subset}}
\newunicodechar{⊥}{\ensuremath{\perp}}
\newunicodechar{∃}{\ensuremath{\exists}}
\newunicodechar{∀}{\ensuremath{\forall}}
\newunicodechar{∛}{\ensuremath{\sqrt[3]{\,}}}
\newunicodechar{∜}{\ensuremath{\sqrt[4]{\,}}}
\newunicodechar{⃗}{\ensuremath{{}^{\to}}}
\newunicodechar{ⁿ}{\ifmmode^{n}\else\textsuperscript{\ensuremath{n}}\fi}
\newunicodechar{ˣ}{\ifmmode^{x}\else\textsuperscript{\ensuremath{x}}\fi}
\newunicodechar{ᵇ}{\ifmmode^{b}\else\textsuperscript{\ensuremath{b}}\fi}
\newunicodechar{ʳ}{\ifmmode^{r}\else\textsuperscript{\ensuremath{r}}\fi}
\newunicodechar{ᵘ}{\ifmmode^{u}\else\textsuperscript{\ensuremath{u}}\fi}
\newunicodechar{ʸ}{\ifmmode^{y}\else\textsuperscript{\ensuremath{y}}\fi}
\newunicodechar{ᵃ}{\ifmmode^{a}\else\textsuperscript{\ensuremath{a}}\fi}
\newunicodechar{ᵉ}{\ifmmode^{e}\else\textsuperscript{\ensuremath{e}}\fi}
\newunicodechar{ᵏ}{\ifmmode^{k}\else\textsuperscript{\ensuremath{k}}\fi}
\newunicodechar{ᵖ}{\ifmmode^{p}\else\textsuperscript{\ensuremath{p}}\fi}
\newunicodechar{ᴺ}{\ifmmode^{N}\else\textsuperscript{\ensuremath{N}}\fi}
\newunicodechar{ᶜ}{\ifmmode^{c}\else\textsuperscript{\ensuremath{c}}\fi}
\newunicodechar{ʰ}{\ifmmode^{h}\else\textsuperscript{\ensuremath{h}}\fi}
\newunicodechar{ᵠ}{\ifmmode^{q}\else\textsuperscript{\ensuremath{q}}\fi}
\newunicodechar{ₙ}{\ifmmode_{n}\else\textsubscript{\ensuremath{n}}\fi}
\newunicodechar{ₐ}{\ifmmode_{a}\else\textsubscript{\ensuremath{a}}\fi}
\newunicodechar{ᵢ}{\ifmmode_{i}\else\textsubscript{\ensuremath{i}}\fi}
\newunicodechar{ᵣ}{\ifmmode_{r}\else\textsubscript{\ensuremath{r}}\fi}
\newunicodechar{ₖ}{\ifmmode_{k}\else\textsubscript{\ensuremath{k}}\fi}
\newunicodechar{ᵦ}{\ifmmode_{\beta}\else\textsubscript{\ensuremath{\beta}}\fi}
\newunicodechar{ₓ}{\ifmmode_{x}\else\textsubscript{\ensuremath{x}}\fi}
\newfontfamily\popdisplay{ArchivoBlack-Regular.ttf}[Path=../../../fonts/]
\newfontfamily\poplogo{SpaceGrotesk-Bold.ttf}[Path=../../../fonts/]
\newfontfamily\popsemi{PlusJakartaSans-SemiBold.ttf}[Path=../../../fonts/]
\newfontfamily\popxbold{PlusJakartaSans-ExtraBold.ttf}[Path=../../../fonts/]
\newfontfamily\popxboldls{PlusJakartaSans-ExtraBold.ttf}[Path=../../../fonts/,LetterSpace=3.0]

\setlist[enumerate]{leftmargin=1.7em,itemsep=5pt,topsep=4pt}
\setlist[itemize]{leftmargin=1.7em,itemsep=4pt,topsep=4pt,label={\color{poporange}\textbullet}}
\setlength{\parindent}{0pt}
\setlength{\parskip}{5pt}
\renewcommand{\arraystretch}{1.25}

% pgfplots : style Pop par défaut
\pgfplotsset{
  every axis/.append style={axis line style={popink,line width=1pt},tick style={popink},
    grid style={popline,line width=0.5pt},label style={font=\small},tick label style={font=\footnotesize}},
  popcurve/.style={poporange,line width=1.8pt,no markers,smooth},
}
% Même style disponible dans un \draw TikZ simple (hors pgfplots)
\tikzset{popcurve/.style={poporange,line width=1.8pt}}
\tikzset{popgrid/.style={popline,very thin}}
\colorlet{popcurve}{poporange} % le nom du style est parfois utilisé comme couleur

% ============================================================
% Pill / squiggle
% ============================================================
\newcommand{\poppill}[3]{%
  \tikz[baseline=-0.55ex]\node[draw=popink,line width=1.2pt,rounded corners=2.6pt,%
    fill=#2,inner xsep=6.5pt,inner ysep=3pt]{%
    \popxboldls\fontsize{7}{7}\selectfont\color{#3}\MakeUppercase{#1}};%
}
\newcommand{\popsquiggle}[1]{%
  \tikz[baseline=-0.4ex]{
    \draw[#1,line width=2.6pt,line cap=round]
      plot [smooth] coordinates {(0,0.09) (0.34,0.01) (0.68,0.21) (1.02,0.01) (1.36,0.10)};
  }%
}
% Mot-clé surligné (marqueur orange sous le texte)
\newcommand{\cle}[1]{{\popxbold\color{popdark}#1}}

% ============================================================
% En-tête / pied
% ============================================================
\pagestyle{fancy}
\fancyhf{}
\renewcommand{\headrulewidth}{0pt}
\renewcommand{\footrulewidth}{0pt}
\lhead{\raisebox{-0.15ex}{\poplogo\fontsize{13}{13}\selectfont\color{popink}OnBuch\color{poporange}+}}
\rhead{\poppill{\DOCMATIERE\ \textbullet\ \DOCNIVEAU\ \textbullet\ Leçon \DOCLECON}{white}{popink}}
\lfoot{\popxbold\fontsize{7.6}{7.6}\selectfont\color{popmuted}\MakeUppercase{Cours signé OnBuch+}\ \textbullet\ {\popsemi\DOCTITRE}}
\rfoot{\tikz[baseline=-0.6ex]\node[rounded corners=8.5pt,fill=popink,minimum height=17pt,minimum width=17pt,inner xsep=5pt,inner sep=0]{\popxbold\fontsize{8}{8}\selectfont\color{popcream}\thepage\,/\,\pageref*{LastPage}};}

% ============================================================
% Titres
% ============================================================
\newcommand{\chaptitre}[2]{%
  \begin{tcolorbox}[enhanced,colback=poporange,colframe=popink,boxrule=2.6pt,arc=11pt,
      drop shadow=popink,left=15pt,right=15pt,top=12pt,bottom=11pt,before skip=3pt,after skip=18pt]
    \poppill{#2}{popink}{popcream}\par\vspace{9pt}
    {\raggedright\hyphenpenalty=10000\exhyphenpenalty=10000\popdisplay\fontsize{22}{24}\selectfont\color{popcream}\MakeUppercase{#1}\par}
  \end{tcolorbox}
}
% Section numérotée : pastille orange avec le numéro + titre Archivo
\newcounter{popsec}
\newcommand{\coursec}[1]{%
  \stepcounter{popsec}\setcounter{popsub}{0}%
  \par\vspace{16pt}\noindent
  \begin{minipage}{\linewidth}
  \tikz[baseline=-0.7ex]\node[draw=popink,line width=1.6pt,fill=poporange,rounded corners=4pt,minimum size=22pt,inner sep=0]{\popdisplay\fontsize{12}{12}\selectfont\color{popcream}\thepopsec};%
  \hspace{8pt}{\popdisplay\fontsize{15}{17}\selectfont\color{popink}\MakeUppercase{#1}}\par
  \vspace{4pt}\noindent{\color{poporange}\rule{36pt}{3.6pt}}
  \end{minipage}\par\nopagebreak\vspace{8pt}
}
\newcounter{popsub}
\newcommand{\courssub}[1]{%
  \stepcounter{popsub}%
  \par\vspace{9pt}\noindent{\popxbold\fontsize{11.5}{13}\selectfont\color{popink}{\color{poporange}\thepopsec.\thepopsub}\hspace{6pt}#1}\par\nopagebreak\vspace{4pt}
}

% ============================================================
% Boîtes pédagogiques — même recette : fond blanc, bordure épaisse colorée,
% ombre dure encre, badge plein. Titre optionnel [#1].
% ============================================================
\tcbset{popbase/.style={breakable,enhanced,colback=white,boxrule=2.3pt,arc=9pt,
  shadow={3pt}{-3pt}{0pt}{popink},left=14pt,right=14pt,top=11pt,bottom=11pt,before skip=11pt,after skip=15pt,
  fontupper=\popsemi\fontsize{10.6}{14.5}\selectfont\color{popink},
  fontlower=\popsemi\fontsize{10.6}{14.5}\selectfont\color{popink}}}
% Coche dessinée (le glyphe ✓ n'existe pas dans les polices du document)
\newcommand{\popcheck}[1]{\tikz[baseline=-0.5ex]\draw[#1,line width=1.6pt,line cap=round,line join=round](-0.13,0.0)--(-0.04,-0.09)--(0.13,0.1);}
\providecommand{\coche}{\popcheck{popgreen}}
\providecommand{\resultat}[1]{\par\smallskip\noindent\fcolorbox{popgreen}{popgreenL}{\parbox{\dimexpr\linewidth-2\fboxsep-2\fboxrule\relax}{\textbf{Résultat.} #1}}\par\smallskip}
\newcommand{\popboxhead}[3]{\poppill{#1}{#2}{white}\ifx\relax#3\relax\else\hspace{6pt}{\popxbold\fontsize{10.6}{12}\selectfont\color{popink}#3}\fi\par\vspace{7pt}}

\newtcolorbox{aretenir}[1][]{popbase,colframe=popgreen,before upper={\popboxhead{À retenir}{popgreen}{#1}}}
\newtcolorbox{exemplebox}[1][]{popbase,colframe=poporange,before upper={\popboxhead{Exemple}{poporange}{#1}}}
\newtcolorbox{definition}[1][]{popbase,colframe=popblue,before upper={\popboxhead{Définition}{popblue}{#1}}}
\newtcolorbox{propriete}[1][]{popbase,colframe=poppurple,before upper={\popboxhead{Propriété}{poppurple}{#1}}}
\newtcolorbox{methode}[1][]{popbase,colframe=popgold,before upper={\popboxhead{Méthode}{popgold}{#1}}}
\newtcolorbox{attention}[1][]{popbase,colframe=poppink,before upper={\popboxhead{Attention}{poppink}{#1}}}
\newtcolorbox{experience}[1][]{popbase,colframe=popblue,colback=popblueL,before upper={\popboxhead{Expérience}{popblue}{#1}}}
% « Le savais-tu ? » : carte violette pleine (contexte, culture, Cameroun)
\newtcolorbox{savaistu}[1][]{popbase,colback=poppurple,colframe=popink,
  fontupper=\popsemi\fontsize{10.4}{14.2}\selectfont\color{white},
  before upper={\poppill{Le savais-tu\,?}{white}{poppurple}\ifx\relax#1\relax\else\hspace{6pt}{\popxbold\color{white}#1}\fi\par\vspace{7pt}}}
% Exercice résolu : énoncé en haut, \tcblower puis la solution sur fond crème
\newcounter{popexo}\newcounter{popexr}
\newtcolorbox{exoresolu}[1][]{popbase,colframe=poporange,colbacklower=popcream,
  segmentation style={popink,line width=1.2pt,dashed},
  before upper={\stepcounter{popexr}\popboxhead{Exercice résolu \thepopexr}{poporange}{#1}},
  before lower={\poppill{Solution}{popink}{popcream}\par\vspace{6pt}}}
% Exercice (non résolu en place) — la correction est dans la partie Corrigés
\newtcolorbox{exercice}[1][]{popbase,colframe=popink,
  before upper={\stepcounter{popexo}\popboxhead{Exercice \thepopexo}{popink}{#1}}}
% Corrigé
\newtcolorbox{corrige}[1][]{popbase,colframe=popgreen,colback=popgreenL,
  before upper={\popboxhead{Corrigé}{popgreen}{#1}}}
% Objectifs / prérequis / bilan
\newtcolorbox{objectifs}{popbase,colframe=popink,colback=poporangeL,
  before upper={\popboxhead{Objectifs de la leçon}{popink}{}}}
\newtcolorbox{prerequis}{popbase,colframe=popink,colback=popblueL,
  before upper={\popboxhead{Prérequis}{popblue}{}}}
\newtcolorbox{bilan}{popbase,colframe=popink,colback=white,boxrule=2.8pt,
  before upper={\popboxhead{Fiche bilan}{poporange}{L'essentiel à connaître pour l'examen}}}
% Figure encadrée avec légende
\newtcolorbox{popfigure}[1][]{enhanced,breakable=false,colback=white,colframe=popline,boxrule=1.4pt,arc=8pt,
  left=10pt,right=10pt,top=10pt,bottom=8pt,before skip=10pt,after skip=12pt,halign=center,
  lower separated=false,halign lower=center,
  fontlower=\popsemi\fontsize{9}{11.5}\selectfont\color{popmuted},
  #1}
\newcommand{\legende}[1]{\par\vspace{6pt}{\popsemi\fontsize{9}{11.5}\selectfont\color{popmuted}#1}}

% ============================================================
% Signature OnBuch+
% ============================================================
\newcommand{\poplogoinline}[1]{{\poplogo\fontsize{#1}{#1}\selectfont\color{popink}OnBuch\color{poporange}+}}
\newcommand{\signatureonbuch}{%
  \begin{tcolorbox}[enhanced,colback=popink,colframe=popink,boxrule=2.2pt,arc=10pt,
      drop shadow=poporange,left=14pt,right=14pt,top=10pt,bottom=10pt,before skip=0pt,after skip=0pt]
    \tikz[baseline=-0.7ex]\node[circle,fill=popgreen,draw=popcream,line width=1.3pt,minimum size=22pt,inner sep=0]{\popcheck{white}};%
    \hspace{9pt}\begin{minipage}[c]{0.62\linewidth}
      {\popxbold\fontsize{12}{13}\selectfont\color{popcream}Signé\ \textbullet\ {\poplogo OnBuch\color{poporange}+}}\par\vspace{2pt}
      {\raggedright\popsemi\fontsize{8}{10.5}\selectfont\color{popline}\MakeUppercase{Équipe pédagogique OnBuch+}\\\MakeUppercase{Conforme au programme officiel MINESEC}\par}
    \end{minipage}\hfill
    {\poplogo\fontsize{22}{22}\selectfont\color{popcream}OnBuch\color{poporange}+}
  \end{tcolorbox}%
}

% ============================================================
% Page de garde
% #1 titre · #2 matière · #3 niveau · #4 module · #5 n° leçon · #6 durée ·
% #7 description / objectifs (texte ou itemize)
% ============================================================
\newcommand{\pagegarde}[7]{%
  \thispagestyle{empty}
  \newgeometry{a4paper,margin=1.7cm,top=1.6cm,bottom=1.4cm}%
  \begin{tikzpicture}[remember picture,overlay]
    \fill[poporange!18] ([xshift=-3.2cm,yshift=-3.6cm]current page.north east) circle (4.6cm);
    \fill[poppurple!14] ([xshift=2.4cm,yshift=7.6cm]current page.south west) circle (3.8cm);
  \end{tikzpicture}%
  {\poplogo\fontsize{30}{30}\selectfont\color{popink}OnBuch\color{poporange}+}\hfill
  \raisebox{0.6ex}{\poppill{Cours signé \textbullet\ Premium}{popink}{popcream}}\par
  \vspace{1pt}\popsquiggle{poppurple}\par
  \vspace{8pt}
  \begin{tcolorbox}[enhanced,colback=poporange,colframe=popink,boxrule=2.8pt,arc=13pt,
      drop shadow=popink,left=18pt,right=18pt,top=12pt,bottom=13pt,after skip=10pt]
    \poppill{#2 \textbullet\ #3}{popink}{popcream}\hspace{5pt}\poppill{#4}{white}{popink}\par\vspace{8pt}
    {\popdisplay\fontsize{14}{14}\selectfont\color{popink}LEÇON #5}\par\vspace{4pt}
    {\raggedright\hyphenpenalty=10000\exhyphenpenalty=10000\popdisplay\fontsize{25}{27}\selectfont\color{popcream}\MakeUppercase{#1}\par}
    \vspace{8pt}
    {\popxbold\fontsize{9.5}{9.5}\selectfont\color{popink}\MakeUppercase{Durée conseillée : #6}}
  \end{tcolorbox}
  \begin{tcolorbox}[enhanced,colback=white,colframe=popink,boxrule=2.2pt,arc=10pt,
      drop shadow=popink,left=16pt,right=16pt,top=10pt,bottom=10pt,after skip=8pt]
    \poppill{Dans cette leçon}{poppurple}{white}\par\vspace{5pt}
    \popsemi\fontsize{9.3}{12.6}\selectfont\color{popink}#7
  \end{tcolorbox}
  \noindent\begin{minipage}[t]{0.487\linewidth}
    \begin{tcolorbox}[enhanced,colback=popgreen,colframe=popink,boxrule=2.2pt,arc=10pt,
        drop shadow=popink,left=11pt,right=11pt,top=10pt,bottom=10pt,height=3.1cm,valign=center]
      \tcbox[colback=white,colframe=popink,boxrule=1.3pt,arc=5pt,boxsep=0pt,nobeforeafter,
        left=3pt,right=3pt,top=3pt,bottom=3pt,valign=center]{\qrcode[height=1.9cm]{https://play.google.com/store/apps/details?id=com.luvvix.onbuch&hl=fr}}%
      \hspace{8pt}\begin{minipage}[c]{0.5\linewidth}
        {\popdisplay\fontsize{11}{12.5}\selectfont\color{white}\MakeUppercase{Télécharge OnBuch}}\par\vspace{4pt}
        {\raggedright\popsemi\fontsize{8.4}{11}\selectfont\color{white}Gratuit \textbullet\ Google Play\\Tuteur IA, annales, concours, résultats\par}
      \end{minipage}
    \end{tcolorbox}
  \end{minipage}\hfill
  \begin{minipage}[t]{0.487\linewidth}
    \begin{tcolorbox}[enhanced,colback=poppurple,colframe=popink,boxrule=2.2pt,arc=10pt,
        drop shadow=popink,left=11pt,right=11pt,top=10pt,bottom=10pt,height=3.1cm,valign=center]
      \tikz[baseline=-0.7ex]\node[circle,fill=white,draw=popink,line width=1.3pt,minimum size=1.6cm,inner sep=0]{\popdisplay\fontsize{24}{24}\selectfont\color{poppurple}+};%
      \hspace{8pt}\begin{minipage}[c]{0.6\linewidth}
        {\popdisplay\fontsize{11}{12.5}\selectfont\color{white}\MakeUppercase{Passe à OnBuch+}}\par\vspace{4pt}
        {\raggedright\popsemi\fontsize{8.4}{11}\selectfont\color{white}Tous les cours signés, Tuteur IA illimité, fiches PDF — onglet OnBuch+ de l'app\par}
      \end{minipage}
    \end{tcolorbox}
  \end{minipage}
  \vspace{8pt}
  \popcontenu
  \vfill
  \signatureonbuch
  \restoregeometry
  \newpage
}

% Rangée « Ce cours contient » (page de garde)
\newcommand{\poptile}[3]{% #1 couleur #2 gros texte #3 légende
  \begin{tcolorbox}[enhanced,colback=#1,colframe=popink,boxrule=2pt,arc=9pt,shadow={2.5pt}{-2.5pt}{0pt}{popink},
      left=6pt,right=6pt,top=9pt,bottom=9pt,halign=center,valign=center,height=1.75cm,width=\linewidth,nobeforeafter]
    {\popdisplay\fontsize{12.5}{13}\selectfont\color{white}\MakeUppercase{#2}}\par\vspace{3pt}
    {\centering\popsemi\fontsize{7.8}{9.5}\selectfont\color{white}#3\par}
  \end{tcolorbox}}
\newcommand{\popcontenu}{%
  \par\noindent{\popxboldls\fontsize{7.5}{7.5}\selectfont\color{popmuted}CE COURS SIGNÉ CONTIENT}\par\vspace{7pt}
  \noindent\begin{minipage}[t]{0.235\linewidth}\poptile{popblue}{Cours}{complet, illustré, pas à pas}\end{minipage}\hfill
  \begin{minipage}[t]{0.235\linewidth}\poptile{poporange}{Méthodes}{et exercices résolus}\end{minipage}\hfill
  \begin{minipage}[t]{0.235\linewidth}\poptile{poppink}{Exercices}{type examen + corrigés}\end{minipage}\hfill
  \begin{minipage}[t]{0.235\linewidth}\poptile{popgreen}{Fiche bilan}{l'essentiel à retenir}\end{minipage}\par}

% Sommaire
\newtcolorbox{sommaire}{enhanced,colback=white,colframe=popink,boxrule=2.2pt,arc=10pt,
  drop shadow=popink,left=18pt,right=18pt,top=13pt,bottom=13pt,before skip=6pt,after skip=20pt,
  before upper={\poppill{Sommaire}{poppurple}{white}\par\vspace{9pt}\popsemi\fontsize{10.6}{14.5}\selectfont\color{popink}}}

% Signature de fin de cours — poussée tout en bas de la dernière page
\newcommand{\findecours}{%
  \par\vfill\nopagebreak
  \begin{center}\popsquiggle{poporange}\end{center}\vspace{4pt}
  \signatureonbuch
}
\AtBeginDocument{\def\llbracket{\mathopen{[\mkern-3.5mu[}}\def\rrbracket{\mathclose{]\mkern-3.5mu]}}} % intervalles d'entiers (glyphe ⟦ absent de la police)
% Glyphes absents de Fira Math : redéfinis à partir de glyphes disponibles
\AtBeginDocument{%
  \def\setminus{\mathbin{\backslash}}\let\smallsetminus\setminus
  \def\vdots{\mathord{\vbox{\baselineskip3.2pt\lineskiplimit0pt\kern4pt\hbox{.}\hbox{.}\hbox{.}}}}%
  \def\ddots{\mathinner{\mkern1mu\raise6.5pt\hbox{.}\mkern2mu\raise3.6pt\hbox{.}\mkern2mu\raise0.7pt\hbox{.}\mkern1mu}}%
  \def\triangle{\mathord{\scalebox{0.9}{$\symup{\Delta}$}}}%
  \def\nabla{\mathord{\raisebox{\depth}{\scalebox{1}[-1]{$\symup{\Delta}$}}}}%
  \def\checkmark{\popcheck{popdark}}%
  \def\star{\mathbin{\ast}}\def\bigstar{\mathord{\ast}}%
  \def\diamond{\mathbin{\scalebox{0.6}{$\Diamond$}}}%
  \def\bigcup{\mathop{\mathchoice{\raisebox{-0.3em}{\scalebox{1.7}{$\cup$}}}{\scalebox{1.25}{$\cup$}}{\cup}{\cup}}\displaylimits}%
  \def\bigcap{\mathop{\mathchoice{\raisebox{-0.3em}{\scalebox{1.7}{$\cap$}}}{\scalebox{1.25}{$\cap$}}{\cap}{\cap}}\displaylimits}%
  \def\lesseqqgtr{\mathrel{\lessgtr}}\def\bigoplus{\mathop{\oplus}\displaylimits}%
}
\providecommand{\ssi}{si et seulement si} % abréviation fréquente
\AtBeginDocument{\providecommand{\wideparen}{\overparen}} % arc de cercle

%% FIGFIX-BEGIN v5 — correctifs globaux des figures (docs/onbuch-generation/tools/figfix)
\usepackage{adjustbox}\usepackage{etoolbox}\tcbuselibrary{raster}
% toute figure TikZ/pgfplots plus large que la ligne est réduite à la largeur disponible
\BeforeBeginEnvironment{tikzpicture}{\begin{adjustbox}{max width=\linewidth}}
\AfterEndEnvironment{tikzpicture}{\end{adjustbox}}
% légendes pgfplots lisibles par-dessus les courbes
\pgfplotsset{every axis/.append style={legend style={fill=white,fill opacity=0.92,text opacity=1,draw=black!15,inner sep=3pt}}}
% étiquettes d'arêtes (syntaxe edge["..."]) avec fond blanc
\tikzset{every edge quotes/.append style={fill=white,inner sep=1.2pt}}
%% FIGFIX-END
\begin{document}

\pagegarde{Les catastrophes : classification, gestion et prévention}{SVTEEHB}{Tle TI}{Module III : L'Éducation à l'Environnement et au Développement Durable}{N09}{3 h}{Cette leçon définit la notion de catastrophe, propose une classification fondée sur l'origine (naturelle ou humaine) et présente les grands types de catastrophes illustrés par des exemples camerounais et mondiaux. Elle aborde enfin la gestion des catastrophes à travers la prévision, la prédiction, la prévention et les gestes qui sauvent.
\vspace{4pt}
\begin{multicols}{2}\small
\begin{itemize}
\item À la fin de cette leçon, je sais définir une catastrophe et la distinguer d'un simple accident ou d'un aléa
\item Je sais classer une catastrophe donnée selon son origine (naturelle ou humaine) en justifiant mon choix
\item Je sais citer et décrire les principales catastrophes d'origine naturelle (séismes, éruptions, inondations, glissements de terrain, explosions limniques, cyclones, sécheresses)
\item Je sais citer et décrire les principales catastrophes d'origine humaine (accidents industriels et technologiques, pollutions majeures, feux, conflits, accidents de transport)
\item Je sais distinguer prévision, prédiction et prévention dans la gestion des catastrophes
\item Je sais interpréter des documents (cartes de risques, tableaux de données, courbes de fréquence) pour évaluer la vulnérabilité d'une région
\end{itemize}\end{multicols}}

\begin{sommaire}
\begin{multicols}{2}
\begin{enumerate}
\item Définition et caractéristiques d'une catastrophe
\item Classification des catastrophes
\item Les catastrophes d'origine naturelle
\item Les catastrophes d'origine humaine
\item La gestion des catastrophes : prévision, prédiction, prévention
\item Les gestes qui sauvent et la conduite à tenir
\item Activité d'intégration
\item Méthodes et exercices résolus
\item Exercices
\item Corrigés des exercices
\item Fiche bilan
\end{enumerate}
\end{multicols}
\end{sommaire}

\begin{objectifs}
\begin{itemize}
\item À la fin de cette leçon, je sais définir une catastrophe et la distinguer d'un simple accident ou d'un aléa
\item Je sais classer une catastrophe donnée selon son origine (naturelle ou humaine) en justifiant mon choix
\item Je sais citer et décrire les principales catastrophes d'origine naturelle (séismes, éruptions, inondations, glissements de terrain, explosions limniques, cyclones, sécheresses)
\item Je sais citer et décrire les principales catastrophes d'origine humaine (accidents industriels et technologiques, pollutions majeures, feux, conflits, accidents de transport)
\item Je sais distinguer prévision, prédiction et prévention dans la gestion des catastrophes
\item Je sais interpréter des documents (cartes de risques, tableaux de données, courbes de fréquence) pour évaluer la vulnérabilité d'une région
\item Je sais élaborer un outil de sensibilisation (affiche, dépliant, sketch) sur la gestion d'une catastrophe
\item Je sais décrire et pratiquer les gestes qui sauvent en cas de catastrophe (position latérale de sécurité, garrot, réaction face à un incendie ou une inondation)
\end{itemize}
\end{objectifs}

\begin{prerequis}
\begin{itemize}
\item Notions sur la structure interne de la Terre et la tectonique des plaques (classe de Première)
\item Notion d'écosystème et d'équilibre naturel (classes antérieures)
\item Notions sur la pollution de l'air, de l'eau et des sols (Module III, leçons précédentes)
\item Notion de développement durable et d'éducation à l'environnement (début du Module III)
\end{itemize}
\end{prerequis}

\newpage

\coursec{Définition et caractéristiques d'une catastrophe}

\textbf{Situation d'accroche.} Le 28 octobre 2019, à Bafoussam (région de l'Ouest du Cameroun), un pan de colline détrempé par des pluies torrentielles s'effondre en pleine nuit sur un quartier : plusieurs habitations sont ensevelies et plus de 40 personnes perdent la vie. Trente-trois ans plus tôt, le 21 août 1986, le lac Nyos, dans la région du Nord-Ouest, libérait brutalement un nuage de dioxyde de carbone qui asphyxiait environ \num{1746} personnes et des milliers de têtes de bétail. Chaque année, les inondations de Yaoundé et de l'Extrême-Nord, les feux de brousse de l'Adamaoua et les accidents industriels rappellent que les catastrophes, naturelles ou humaines, font partie de notre quotidien.

\textbf{Problématique.} Qu'est-ce qui distingue une catastrophe d'un simple accident ? Quels critères permettent de la caractériser ? Peut-on prévoir, prévenir ces événements, et comment réagir pour sauver des vies ? Pour répondre à ces questions, commençons par définir rigoureusement la notion de catastrophe et les concepts qui l'entourent : l'aléa, le risque et la vulnérabilité.

\courssub{Qu'est-ce qu'une catastrophe ?}

\textbf{Intuition d'abord.} Observe autour de toi : un incendie qui détruit une cuisine est un drame familial, mais le quartier continue de vivre ; en revanche, lorsque le lac Nyos se dégaze, ce sont des villages entiers qui disparaissent, les hôpitaux de la région sont débordés, les secours doivent venir de loin. La différence n'est pas seulement une question de nombre de victimes : c'est une question d'\emph{ampleur} et de \emph{dépassement des capacités de réaction}. C'est cette intuition que la définition scientifique va formaliser.

\begin{definition}[Catastrophe]
Une \cle{catastrophe} est un événement \cle{soudain}, d'origine \cle{naturelle} (séisme, inondation, éruption volcanique, glissement de terrain\ldots) ou \cle{humaine} (accident industriel, effondrement de bâtiment, pollution majeure, conflit\ldots), qui provoque des \cle{dommages importants} : pertes en vies humaines, destructions matérielles considérables et perturbations durables de l'environnement, au point de \cle{dépasser les capacités locales de réaction} (secours, hôpitaux, moyens matériels) et de nécessiter une aide extérieure.
\end{definition}

\begin{aretenir}[Les trois critères cumulatifs d'une catastrophe]
Un événement est qualifié de catastrophe lorsqu'il réunit trois critères :
\begin{enumerate}
\item la \cle{soudaineté} : l'événement survient brutalement, en quelques secondes à quelques heures, laissant peu de temps pour réagir ;
\item l'\cle{ampleur des dégâts} : nombreuses victimes, destructions étendues, atteinte durable des milieux ;
\item le \cle{dépassement des capacités locales de réaction} : les moyens locaux (centres de santé, protection civile, populations) sont insuffisants, ce qui impose l'intervention de l'État, voire de la communauté internationale.
\end{enumerate}
\end{aretenir}

\courssub{Aléa, risque, catastrophe : trois notions à ne pas confondre}

\textbf{Observation.} Le Mont Cameroun (\SI{4095}{m}), volcan actif, est entré en éruption en 1999 et en 2000. La lave a détruit des plantations, mais les coulées se sont arrêtées avant d'atteindre les quartiers densément peuplés de Buéa. Le même volcan, avec la même activité, aurait causé une catastrophe bien plus grave si la ville avait été touchée. Cela montre que le danger ne dépend pas seulement du phénomène naturel, mais aussi de ce qui se trouve exposé.

\begin{definition}[Aléa, vulnérabilité, enjeux et risque]
\begin{itemize}
\item L'\cle{aléa} est un événement dangereux \emph{possible}, caractérisé par sa nature, son intensité et sa probabilité d'occurrence dans un lieu et un temps donnés (ex. : une éruption du Mont Cameroun, une crue du Logone).
\item La \cle{vulnérabilité} est le degré de fragilité des personnes et des biens exposés (habitations construites sur des pentes instables, absence de système d'alerte, pauvreté).
\item Les \cle{enjeux} sont les personnes, les biens et les activités exposés à l'aléa (populations, cultures, hôpitaux, routes).
\item Le \cle{risque} est la combinaison de l'aléa, de la vulnérabilité et des enjeux : c'est la possibilité de dommages résultant de la rencontre entre un aléa et des enjeux vulnérables. On écrit souvent, de façon schématique :
\[ \text{Risque} = \text{Aléa} \times \text{Vulnérabilité} \times \text{Enjeux} \]
\item La \cle{catastrophe} est l'aléa \emph{réalisé} : le risque s'est concrétisé et a causé des dommages majeurs.
\end{itemize}
\end{definition}

\begin{attention}[Piège classique au Bac]
Ne confonds jamais \emph{aléa} et \emph{risque}. L'aléa est le phénomène dangereux lui-même (le séisme, la crue) : il existe même en l'absence de toute présence humaine. Le risque n'existe que s'il y a des enjeux vulnérables exposés : un séisme de magnitude 7 au milieu d'un désert inhabité est un aléa réalisé, mais pas une catastrophe. Inversement, une crue modérée dans un quartier précaire densément peuplé (comme certains quartiers de Yaoundé en bordure de marécages) peut devenir une catastrophe.
\end{attention}

\begin{popfigure}
\begin{center}
\begin{tikzpicture}[font=\small]
% Aléa
\draw[fill=popblue!25, draw=popblue, thick] (0,0) circle (1.6);
\node[popdark, font=\bfseries] at (0,0.9) {ALÉA};
\node[align=center, popdark] at (0,0) {phénomène\\ dangereux\\ possible};
% Enjeux vulnérables
\draw[fill=poporange!25, draw=poporange, thick] (2.4,0) circle (1.6);
\node[popdark, font=\bfseries] at (2.4,0.9) {ENJEUX};
\node[align=center, popdark] at (2.4,0) {personnes, biens\\ vulnérables\\ exposés};
% Zone de risque
\node[align=center, poppurple, font=\bfseries] at (1.2,-0.75) {RISQUE};
% Flèche
\draw[-{Stealth[length=3mm]}, very thick, poppurple] (1.2,-2.1) -- (1.2,-3.1);
\node[right, popmuted] at (1.4,-2.6) {l'aléa se réalise};
% Catastrophe
\draw[fill=poppink!30, draw=poppink, very thick, rounded corners] (-1.4,-5.2) rectangle (3.8,-3.7);
\node[align=center, popdark, font=\bfseries] at (1.2,-4.15) {CATASTROPHE};
\node[align=center, popdark] at (1.2,-4.75) {dommages majeurs : victimes, destructions,\\ dépassement des capacités locales};
\end{tikzpicture}
\end{center}
\legende{Relation entre aléa, risque et catastrophe. Le risque naît de la rencontre (zone commune) entre un aléa et des enjeux vulnérables ; la catastrophe est la concrétisation de ce risque, lorsque l'aléa se réalise et que les dégâts dépassent les capacités locales de réaction.}
\end{popfigure}

\courssub{Catastrophe ou accident : une question d'échelle}

\textbf{Démarche d'investigation.} Comparons deux événements réels à l'aide d'un tableau de données.

\begin{center}
\begin{tabularx}{\linewidth}{|l|X|X|}
\hline
\rowcolor{popblueL} \textbf{Critère} & \textbf{Accident de la route (ex. : collision sur l'axe Douala--Yaoundé)} & \textbf{Explosion limnique du lac Nyos (1986)} \\
\hline
Étendue spatiale & Localisée (quelques mètres) & Plusieurs villages sur environ \SI{25}{km} autour du lac \\
\hline
Victimes & Quelques morts ou blessés & Environ \num{1746} morts, \num{3500} têtes de bétail \\
\hline
Capacités locales & Gendarmerie et hôpital local suffisent & Hôpitaux régionaux débordés, aide nationale et internationale indispensable \\
\hline
Durée des perturbations & Quelques heures & Années (déplacement des populations, dégazage artificiel du lac) \\
\hline
\end{tabularx}
\end{center}

\textbf{Interprétation.} L'accident et la catastrophe sont tous deux des événements soudains causant des dommages, mais ils se distinguent par l'\emph{échelle} : nombre de victimes, étendue spatiale des dégâts, et surtout capacité des structures locales à faire face. Un accident reste gérable avec les moyens locaux ; une catastrophe les submerge.

\begin{attention}[Erreur fréquente]
Confondre \emph{accident} et \emph{catastrophe}. Un accident devient une catastrophe lorsque l'ampleur des dégâts \cle{dépasse les capacités locales de réaction}. Ainsi, l'effondrement d'un immeuble à Douala qui tue plusieurs dizaines de personnes et submerge les services d'urgence de la ville bascule du statut d'accident à celui de catastrophe. Le critère n'est pas un nombre absolu de victimes, mais le rapport entre les dégâts et les moyens disponibles localement.
\end{attention}

\courssub{Exemples de catastrophes au Cameroun}

Le territoire camerounais, par sa diversité géologique et climatique, est exposé à de nombreux aléas. Voici les catastrophes majeures qui ont marqué l'histoire récente du pays.

\begin{exemplebox}[L'explosion limnique du lac Nyos (21 août 1986)]
Le lac Nyos, lac de cratère volcanique de la région du Nord-Ouest, accumulait en profondeur du dioxyde de carbone \ce{CO2} d'origine magmatique, dissous sous pression. Dans la nuit du 21 août 1986, le lac s'est brutalement \emph{dégazé} : un nuage de \ce{CO2}, gaz incolore, inodore et plus dense que l'air, s'est répandu dans les vallées voisines. Résultat : environ \num{1746} personnes et \num{3500} têtes de bétail sont mortes asphyxiées, les survivants présentant des lésions cutanées et des troubles respiratoires. C'est l'une des catastrophes naturelles les plus meurtrières de l'histoire du Cameroun. Depuis, des colonnes de dégazage artificiel ont été installées pour évacuer progressivement le \ce{CO2} du fond du lac : c'est un exemple de \emph{prévention}.
\end{exemplebox}

\begin{savaistu}[Deux lacs tueurs sur trois sont au Cameroun]
Trois lacs au monde sont connus pour présenter un risque d'\cle{explosion limnique} : le lac Nyos et le lac Monoun, tous deux au Cameroun, et le lac Kivu (entre le Rwanda et la République démocratique du Congo). Le 15 août 1984, le lac Monoun avait déjà libéré un nuage de \ce{CO2} qui tua 37 personnes — un signal d'alarme qui, hélas, n'avait pas été suffisamment pris en compte avant le drame de Nyos deux ans plus tard. Cela illustre l'importance de la \emph{mémoire des catastrophes} dans la prévention.
\end{savaistu}

D'autres catastrophes jalonnent l'actualité camerounaise :
\begin{itemize}
\item le \textbf{glissement de terrain de Bafoussam} (octobre 2019) : plus de 40 morts, des habitations ensevelies après des pluies intenses sur des pentes urbanisées ;
\item les \textbf{inondations récurrentes} de Yaoundé (quartiers construits dans les bas-fonds marécageux) et de l'Extrême-Nord (crues du Logone, déplacement de dizaines de milliers de personnes, destruction des champs) ;
\item les \textbf{éruptions du Mont Cameroun} (1999, 2000) et les \textbf{feux de brousse} de la savane adamaouaïenne ;
\item les \textbf{catastrophes d'origine humaine} : effondrements de bâtiments, explosions de dépôts de carburant, pollutions industrielles.
\end{itemize}

\begin{popfigure}
\begin{center}
\begin{tikzpicture}[font=\small, scale=1.05]
% Contour très simplifié du Cameroun
\draw[thick, popgreen!70!black, fill=popgreen!15]
  (0.6,0) -- (2.2,0.2) -- (3.2,1.0) -- (3.0,2.2) -- (3.4,3.2) -- (2.8,4.4)
  -- (2.2,5.6) -- (2.6,6.6) -- (2.2,7.4) -- (1.4,7.6) -- (0.9,6.8)
  -- (0.4,5.6) -- (0.8,4.4) -- (0.2,3.4) -- (0.5,2.2) -- (0.1,1.1) -- cycle;
\node[popmuted, font=\footnotesize] at (1.7,-0.5) {Golfe de Guinée};
% Nyos (Nord-Ouest)
\fill[poppink] (0.9,3.6) circle (0.09);
\node[right, popdark, align=left, font=\footnotesize] at (1.05,3.75) {\textbf{Lac Nyos} (1986) : explosion limnique, $\approx$ \num{1746} morts};
\fill[poppink] (1.15,3.25) circle (0.07);
\node[right, popdark, font=\footnotesize] at (1.3,3.1) {Lac Monoun (1984) : 37 morts};
% Bafoussam (Ouest)
\fill[poporange] (1.35,2.55) circle (0.09);
\node[right, popdark, align=left, font=\footnotesize] at (1.5,2.55) {\textbf{Bafoussam} (2019) : glissement de terrain, $>$ 40 morts};
% Mont Cameroun
\fill[popgold] (0.55,1.35) circle (0.09);
\node[right, popdark, font=\footnotesize] at (0.7,1.2) {Mont Cameroun : éruptions (1999, 2000)};
% Yaoundé
\fill[popblue] (1.7,1.55) circle (0.09);
\node[right, popdark, font=\footnotesize] at (1.85,1.7) {Yaoundé : inondations urbaines};
% Extrême-Nord
\fill[poppurple] (2.1,6.6) circle (0.09);
\node[right, popdark, align=left, font=\footnotesize] at (2.25,6.6) {Extrême-Nord : crues du Logone,\\ inondations récurrentes};
\end{tikzpicture}
\end{center}
\legende{Carte simplifiée du Cameroun localisant quelques catastrophes majeures : explosions limniques des lacs Nyos et Monoun (Nord-Ouest), glissement de terrain de Bafoussam (Ouest), éruptions du Mont Cameroun (Sud-Ouest), inondations de Yaoundé (Centre) et de l'Extrême-Nord.}
\end{popfigure}

\courssub{Les conséquences des catastrophes}

Une catastrophe ne se mesure pas seulement à son bilan immédiat. Ses effets se déploient dans le temps et dans plusieurs domaines, souvent interconnectés.

\begin{center}
\begin{tabularx}{\linewidth}{|l|X|}
\hline
\rowcolor{popblueL} \textbf{Type de conséquence} & \textbf{Exemples} \\
\hline
\textbf{Humaines} & Morts, blessés, disparus, traumatismes psychologiques, déplacements forcés de populations (ex. : survivants de Nyos relogés) \\
\hline
\textbf{Matérielles} & Destruction d'habitations, de routes, de ponts, d'écoles, de cultures et de cheptel (ex. : maisons ensevelies à Bafoussam) \\
\hline
\textbf{Économiques} & Perte des récoltes et du bétail, interruption des activités commerciales, coût des secours et de la reconstruction, appauvrissement des ménages \\
\hline
\textbf{Sanitaires} & Risques d'épidémies (choléra après les inondations par contamination de l'eau), pénurie d'eau potable, rupture des soins \\
\hline
\textbf{Environnementales} & Pollution des sols et des eaux, destruction des habitats naturels, érosion aggravée, perturbation durable des écosystèmes \\
\hline
\end{tabularx}
\end{center}

\begin{exoresolu}[Classer les conséquences d'une catastrophe]
\textbf{Énoncé.} Après les inondations de l'Extrême-Nord, on observe : (a) 30 villages détruits ; (b) une épidémie de choléra ; (c) la perte de \num{12000} hectares de cultures de mil ; (d) la contamination des puits par les eaux de crue ; (e) des centaines de familles déplacées. Classe chaque conséquence selon sa nature (humaine, matérielle, économique, sanitaire ou environnementale).
\tcblower
\textbf{Solution détaillée.}
\begin{itemize}
\item (a) Villages détruits : conséquence \emph{matérielle} (destruction d'habitations et d'infrastructures).
\item (b) Épidémie de choléra : conséquence \emph{sanitaire} (maladie liée à la contamination de l'eau).
\item (c) Perte de \num{12000} hectares de mil : conséquence \emph{économique} (perte de la production agricole et des revenus), avec des répercussions humaines (insécurité alimentaire).
\item (d) Contamination des puits : conséquence \emph{environnementale} (dégradation d'une ressource en eau), qui engendre la conséquence sanitaire (b).
\item (e) Familles déplacées : conséquence \emph{humaine} (déplacement forcé de populations).
\end{itemize}
Remarque : les catégories se recoupent — une même catastrophe produit des effets en chaîne (environnement $\rightarrow$ santé $\rightarrow$ économie). C'est pourquoi la gestion des catastrophes exige une approche globale.
\end{exoresolu}

\begin{aretenir}[L'essentiel de la section]
\begin{itemize}
\item Une \cle{catastrophe} est un événement soudain, naturel ou humain, causant des dommages majeurs qui dépassent les capacités locales de réaction.
\item Trois critères cumulatifs : \cle{soudaineté}, \cle{ampleur des dégâts}, \cle{dépassement des capacités locales}.
\item Risque $=$ Aléa $\times$ Vulnérabilité $\times$ Enjeux ; la catastrophe est l'aléa réalisé sur des enjeux vulnérables.
\item Un accident devient catastrophe par son \emph{échelle}, non par sa nature.
\item Les conséquences sont humaines, matérielles, économiques, sanitaires et environnementales, et s'enchaînent les unes aux autres.
\end{itemize}
\end{aretenir}

\begin{exercice}[Pour t'entraîner]
Pour chacun des événements suivants, indique s'il s'agit d'un aléa, d'un risque ou d'une catastrophe, en justifiant : (1) le Mont Cameroun est un volcan actif ; (2) des quartiers de Buéa sont construits sur d'anciennes coulées de lave ; (3) l'éruption de 1999 a détruit des plantations ; (4) le lac Nyos contient du \ce{CO2} dissous en profondeur ; (5) le dégazage de 1986 a tué environ \num{1746} personnes. Le corrigé sera fourni dans la fiche d'exercices de la leçon.
\end{exercice}

\coursec{Classification des catastrophes}

Dans la section précédente, nous avons défini la catastrophe comme un événement soudain ou progressif qui provoque des dommages importants (pertes humaines, destructions matérielles, perturbations économiques et sociales) dépassant les capacités locales de réponse. Mais toutes les catastrophes ne se ressemblent pas : un séisme, une marée noire, une épidémie de choléra et un effondrement de pont n'ont ni les mêmes causes, ni les mêmes dynamiques, ni les mêmes moyens de prévention. Pour agir efficacement, il faut donc \emph{classer} les catastrophes. C'est précisément l'objet de cette section : dégager les critères de classification, construire les grands groupes, et apprendre à classer correctement n'importe quel événement catastrophique, y compris les cas délicats.

\courssub{Le critère fondamental : l'origine de la catastrophe}

Partons d'une observation simple. Le 21 août 1986, le lac Nyos (région du Nord-Ouest du Cameroun) a brutalement dégazé environ \num{1,6} million de tonnes de \ce{CO2} d'origine volcanique, asphyxiant environ \num{1700} personnes et des milliers de têtes de bétail : personne n'avait rien fait pour provoquer cet événement, il résulte d'un phénomène géologique naturel. À l'inverse, les incendies répétés de dépôts de carburant ou les accidents industriels dans les zones portuaires de Douala résultent directement d'activités humaines. L'élément qui distingue ces deux événements, c'est \cle{l'origine} du phénomène déclencheur.

\begin{definition}[Critère fondamental de classification]
Le critère fondamental de classification des catastrophes est \cle{l'origine} de l'élément déclencheur, c'est-à-dire la cause première de l'événement. On distingue ainsi deux grands groupes :
\begin{itemize}
\item les catastrophes d'\cle{origine naturelle}, qui résultent de phénomènes physiques, géologiques, climatiques ou biologiques se produisant \textbf{sans intervention directe de l'Homme} ;
\item les catastrophes d'\cle{origine humaine} (ou anthropiques), qui résultent directement des activités, des erreurs ou des décisions de l'Homme.
\end{itemize}
\end{definition}

\begin{attention}[Ce que le critère porte réellement]
Le critère de classification porte sur la cause du \textbf{déclencheur}, pas sur la nature des dégâts. Un séisme qui détruit des immeubles reste une catastrophe \emph{naturelle} : ce sont des biens humains qui sont détruits, mais c'est un phénomène géologique qui déclenche. Inversement, une explosion d'usine reste une catastrophe \emph{humaine} même si ses polluants se dispersent dans la nature.
\end{attention}

\courssub{Premier grand groupe : les catastrophes d'origine naturelle}

\begin{definition}[Catastrophe d'origine naturelle]
Une catastrophe d'origine naturelle est un événement catastrophique provoqué par un phénomène de la nature (appelé \cle{aléa naturel}) survenant indépendamment de toute action humaine directe. L'Homme peut en être victime, mais il n'en est pas la cause.
\end{definition}

On distingue classiquement quatre sous-types :

\begin{itemize}
\item \textbf{Catastrophes géologiques (ou telluriques)} : séismes (tremblements de terre), éruptions volcaniques, tsunamis, glissements de terrain d'origine purement naturelle. \emph{Exemples} : l'activité du Mont Cameroun (\SI{4095}{m}), volcan actif dont les éruptions récentes (1999, 2000) ont menacé les plantations et villages de la plaine de Buea ; le séisme de magnitude \num{9,1} au Japon en 2011.
\item \textbf{Catastrophes climatiques et météorologiques} : cyclones, tempêtes, tornades, inondations d'origine pluviale, sécheresses, vagues de chaleur ou de froid. \emph{Exemples} : les inondations récurrentes de Yaoundé et de Douala pendant la saison des pluies ; la sécheresse et l'avancée du désert dans l'Extrême-Nord du Cameroun.
\item \textbf{Catastrophes biologiques naturelles} : épidémies et pandémies dues à des agents pathogènes (choléra, méningite, COVID-19), invasions de criquets pèlerins détruisant les cultures du Sahel.
\item \textbf{Catastrophes cosmiques (ou extraterrestres)} : chute de météorites, tempêtes solaires perturbant les réseaux électriques et de communication. Elles sont rares mais potentiellement majeures.
\end{itemize}

\courssub{Deuxième grand groupe : les catastrophes d'origine humaine}

\begin{definition}[Catastrophe d'origine humaine]
Une catastrophe d'origine humaine (ou anthropique) est un événement catastrophique provoqué directement par une activité, une erreur, une négligence ou une décision humaine. On parle aussi de catastrophe \cle{technologique} lorsqu'elle met en jeu des installations industrielles ou des technologies.
\end{definition}

On en distingue plusieurs sous-types :

\begin{itemize}
\item \textbf{Catastrophes industrielles et technologiques} : explosions d'usines ou de dépôts d'hydrocarbures, rejets toxiques, accidents nucléaires (Tchernobyl, 1986 ; Fukushima, 2011), marées noires (naufrage de pétroliers), ruptures de barrages mal entretenus.
\item \textbf{Catastrophes liées aux transports} : grandes catastrophes ferroviaires (déraillements), aériennes, maritimes ou routières de masse. \emph{Exemple camerounais} : l'accident ferroviaire d'Eséka (octobre 2016), qui a fait environ 79 morts et plus de 500 blessés.
\item \textbf{Catastrophes sanitaires d'origine humaine} : intoxications alimentaires de masse, pollutions chimiques ou bactériologiques de l'eau potable liées à l'insalubrité, épidémies favorisées par la mauvaise gestion des déchets urbains (épidémies de choléra récurrentes dans les quartiers précaires de Douala et de Yaoundé).
\item \textbf{Catastrophes liées aux conflits et aux actes de violence} : guerres, attentats, génocides, qui provoquent des pertes humaines massives, des déplacements de populations (réfugiés, déplacés internes) et la destruction des infrastructures.
\item \textbf{Catastrophes structurelles} : effondrements d'immeubles ou de ponts dus à des défauts de conception, de construction ou de contrôle.
\end{itemize}

\courssub{Les cas intermédiaires : catastrophes naturelles aggravées par l'Homme}

La réalité est souvent plus subtile que la dichotomie précédente. De nombreux événements ont un \textbf{déclencheur naturel} mais une \textbf{ampleur amplifiée} par l'action humaine. C'est un point essentiel, très souvent exigé dans les évaluations.

\begin{exemplebox}[Inondations urbaines à Yaoundé et Douala]
Les pluies diluviennes sont un aléa \emph{naturel}. Mais à Yaoundé comme à Douala, les inondations meurtrières sont aggravées par :
\begin{itemize}
\item la \textbf{déforestation} des collines et des bassins versants, qui supprime l'infiltration de l'eau et accélère le ruissellement ;
\item l'\textbf{urbanisation anarchique} : constructions dans les bas-fonds et les lits majeurs des cours d'eau (zones naturellement inondables) ;
\item le \textbf{colmatage des caniveaux} par les déchets plastiques, qui bloque l'évacuation des eaux.
\end{itemize}
Le déclencheur est naturel (la pluie), mais les facteurs aggravants sont humains : on classe l'événement comme une catastrophe \emph{naturelle aggravée par l'action humaine}.
\end{exemplebox}

\begin{exemplebox}[Glissements de terrain sur les pentes de l'Ouest camerounais]
Dans les régions de l'Ouest et du Nord-Ouest (reliefs escarpés, sols volcaniques), les glissements de terrain surviennent lors de fortes pluies (déclencheur naturel), mais ils sont \emph{favorisés} par les constructions sur les pentes raides, la coupe de la végétation qui stabilisait les sols et l'absence de réseaux de drainage. Là encore, l'origine première est naturelle, l'aggravation est humaine.
\end{exemplebox}

\begin{attention}[Erreur fréquente au Bac]
Ne classe \textbf{jamais} automatiquement toutes les inondations comme des catastrophes purement naturelles. Beaucoup d'inondations urbaines (Yaoundé, Douala, Bafoussam) sont \emph{aggravées} par l'occupation humaine des zones inondables, la déforestation et la mauvaise gestion des déchets. Dans une copie, précise toujours le déclencheur \emph{puis} les facteurs aggravants : c'est ce raisonnement en deux temps qui est attendu et valorisé.
\end{attention}

\courssub{Les critères secondaires de classification}

L'origine est le critère fondamental, mais d'autres critères secondaires permettent d'affiner la description d'une catastrophe.

\textbf{Critère 1 : l'échelle spatiale (étendue géographique).}
\begin{itemize}
\item \textbf{Locale} : elle touche un village, un quartier, une ville (ex. : incendie d'un marché, glissement de terrain sur une colline).
\item \textbf{Nationale ou régionale} : elle affecte une grande partie d'un pays ou plusieurs pays voisins (ex. : sécheresse sahélienne touchant l'Extrême-Nord du Cameroun, le Tchad et le Niger).
\item \textbf{Mondiale} : elle concerne plusieurs continents (ex. : pandémie de COVID-19 à partir de 2020, pandémie de VIH/Sida).
\end{itemize}

\textbf{Critère 2 : la durée et la cinétique (mode d'apparition).}
\begin{itemize}
\item \textbf{Catastrophes brutales (ou soudaines)} : elles surviennent en quelques secondes à quelques heures — séismes, explosions, tsunamis, éruptions limniques (lac Nyos, 1986 : le dégazage a duré quelques heures).
\item \textbf{Catastrophes rampantes (ou progressives)} : elles s'installent lentement, sur des mois ou des années — sécheresses, famines, désertification, épidémies lentes, pollutions chroniques des sols et des nappes phréatiques. Leur danger est d'être souvent détectées tardivement.
\end{itemize}

\textbf{Critère 3 : la prévisibilité.}
\begin{itemize}
\item \textbf{Prévisibles à échéance courte} : cyclones et crues (suivis météorologiques et hydrologiques, prévisions à quelques jours).
\item \textbf{Partiellement surveillables} : éruptions volcaniques (signes précurseurs : sismicité, dégazage, déformations du sol — le Mont Cameroun fait l'objet d'une surveillance).
\item \textbf{Actuellement imprévisibles en date et lieu précis} : les séismes — on connaît les zones à risque, mais on ne sait pas prédire le jour ni l'heure.
\end{itemize}

\courssub{Synthèse : l'arbre de classification}

\begin{popfigure}
\begin{center}
\begin{tikzpicture}[
grow=right,
level distance=4.5cm,
level 1/.style={sibling distance=4.0cm},
level 2/.style={sibling distance=0.9cm, level distance=5.0cm},
every node/.style={font=\small},
edge from parent/.style={draw=popmuted, thick, -{Stealth[length=2mm]}}
]
\node[draw=popdark, fill=popdark, text=white, rounded corners=2pt, font=\small\bfseries, align=center, text width=2.5cm] {Catastrophes}
child { node[draw=poporange, fill=poporangeL, rounded corners=2pt, align=center, font=\small\bfseries, text width=2.8cm] {Origine humaine\\(anthropique)}
  child { node[draw=popmuted, fill=popcream, rounded corners=2pt, align=center, text width=3.0cm, font=\small] {Conflits, guerres,\\attentats} }
  child { node[draw=popmuted, fill=popcream, rounded corners=2pt, align=center, text width=3.0cm, font=\small] {Sanitaires : pollutions,\\épidémies liées à\\l'insalubrité} }
  child { node[draw=popmuted, fill=popcream, rounded corners=2pt, align=center, text width=3.0cm, font=\small] {Transports : ferroviaire\\(Eséka, 2016), aérien,\\maritime, routier} }
  child { node[draw=popmuted, fill=popcream, rounded corners=2pt, align=center, text width=3.0cm, font=\small] {Industrielles et\\technologiques :\\explosions, marées noires,\\accidents nucléaires} }
}
child { node[draw=popblue, fill=popblueL, rounded corners=2pt, align=center, font=\small\bfseries, text width=2.8cm] {Origine naturelle\\(sans intervention\\directe de l'Homme)}
  child { node[draw=popmuted, fill=popcream, rounded corners=2pt, align=center, text width=3.0cm, font=\small] {Cosmiques : météorites,\\tempêtes solaires} }
  child { node[draw=popmuted, fill=popcream, rounded corners=2pt, align=center, text width=3.0cm, font=\small] {Biologiques : épidémies,\\invasions de criquets} }
  child { node[draw=popmuted, fill=popcream, rounded corners=2pt, align=center, text width=3.0cm, font=\small] {Climatiques : cyclones,\\inondations, sécheresses,\\tempêtes} }
  child { node[draw=popmuted, fill=popcream, rounded corners=2pt, align=center, text width=3.0cm, font=\small] {Géologiques : séismes,\\éruptions volcaniques\\(Mont Cameroun),\\tsunamis, glissements\\de terrain} }
};
\end{tikzpicture}
\end{center}
\legende{Arbre de classification des catastrophes selon le critère fondamental de l'origine. Le premier embranchement sépare les catastrophes d'origine naturelle (en bleu) des catastrophes d'origine humaine (en orange) ; chaque groupe se subdivise en sous-types illustrés par des exemples.}
\end{popfigure}

\begin{center}
\begin{tabularx}{\linewidth}{|l|X|X|}
\hline
\rowcolor{popblueL} \textbf{Origine} & \textbf{Exemples camerounais} & \textbf{Exemples mondiaux} \\
\hline
Naturelle géologique & Éruptions du Mont Cameroun (1999, 2000) ; dégazage du lac Nyos (1986, $\approx$ \num{1700} morts) & Séisme de magnitude \num{9,1} au Japon (2011) ; éruption du Vésuve (79 apr. J.-C.) \\
\hline
Naturelle climatique & Inondations de Yaoundé et Douala ; sécheresse de l'Extrême-Nord & Cyclone Katrina (États-Unis, 2005) ; sécheresses sahéliennes \\
\hline
Naturelle biologique & Épidémies de choléra et de méningite ; invasions de criquets & Pandémie de COVID-19 (2020) ; grippe espagnole (1918) \\
\hline
Humaine technologique & Accidents de dépôts de carburant ; pollution industrielle de zones portuaires (Douala) & Tchernobyl (1986) ; marée noire de l'\emph{Exxon Valdez} (1989) \\
\hline
Humaine (transports, structures) & Déraillement d'Eséka (2016, $\approx$ 79 morts) ; effondrements d'immeubles & Naufrage du \emph{Titanic} (1912) ; effondrement du pont de Gênes (2018) \\
\hline
Humaine (conflits) & Déplacements de populations liés aux conflits (Nord-Ouest, Sud-Ouest, Extrême-Nord) & Guerres mondiales ; génocide rwandais (1994) \\
\hline
Mixte (naturelle aggravée par l'Homme) & Inondations urbaines amplifiées par l'urbanisation anarchique et la déforestation ; glissements de terrain sur pentes construites (Ouest) & Inondations amplifiées par l'imperméabilisation des sols ; Fukushima (catastrophe en cascade) \\
\hline
\end{tabularx}
\end{center}

\courssub{Méthode de classement et catastrophes en cascade}

\begin{methode}[Classer une catastrophe en trois étapes]
\begin{enumerate}
\item \textbf{Identifier l'élément déclencheur} (la cause première) : quel phénomène a mis en route l'événement ? Un séisme ? Une pluie ? Une explosion ? Une décision humaine ?
\item \textbf{Déterminer la nature du déclencheur} : relève-t-il d'un phénomène naturel (géologique, climatique, biologique, cosmique) ou d'une activité humaine (industrielle, technique, sanitaire, conflictuelle) ? La catastrophe appartient alors au groupe correspondant.
\item \textbf{Signaler les facteurs aggravants} : rechercher si l'action humaine a amplifié les dégâts (constructions en zone inondable ou sismique, déforestation, absence de normes, insalubrité) ou, au contraire, les a limités (digues, plans d'évacuation, surveillance). Mentionner aussi, si possible, les critères secondaires : échelle spatiale, durée, prévisibilité.
\end{enumerate}
\end{methode}

\begin{exemplebox}[Une catastrophe en cascade : le Japon, 11 mars 2011]
Un séisme de magnitude \num{9,1} (catastrophe \emph{naturelle} géologique) frappe la côte nord-est du Japon. Il génère un tsunami dont les vagues dépassent \SI{10}{m} (toujours naturel). Le tsunami submerge les systèmes de refroidissement de la centrale nucléaire de Fukushima-Daiichi, provoquant la fusion de trois réacteurs et des rejets radioactifs : c'est une catastrophe \emph{technologique}. Un seul enchaînement mobilise donc les deux grands groupes : l'aléa naturel est le déclencheur, l'accident nucléaire en est la conséquence humaine et technologique. On parle de \cle{catastrophe en cascade}.
\end{exemplebox}

\begin{savaistu}[Les catastrophes « NaTech »]
Les spécialistes de la gestion des risques parlent de catastrophe \cle{NaTech} (contraction de \emph{Natural} et \emph{Technological}) lorsqu'un aléa naturel déclenche un accident technologique : un séisme qui rompt un oléoduc, une inondation qui submerge une usine chimique, la foudre qui enflamme un dépôt de carburant. Fukushima (2011) est l'exemple NaTech le plus célèbre. Au Cameroun, la vigilance sur les installations industrielles et portuaires de Douala, ville exposée aux fortes pluies et située en zone côtière, relève exactement de cette logique de prévention des effets en cascade.
\end{savaistu}

\begin{exoresolu}[Classer trois événements]
Énoncé. Classer chacun des événements suivants en précisant le groupe, le sous-type et les éventuels facteurs aggravants : (a) le dégazage du lac Nyos (1986) ; (b) les inondations meurtrières de quartiers de Yaoundé construits dans les bas-fonds ; (c) le déraillement d'Eséka (2016).
\tcblower
Solution détaillée.
\begin{itemize}
\item (a) \textbf{Lac Nyos} : le déclencheur est un dégazage massif de \ce{CO2} d'origine volcanique, phénomène géologique sans intervention humaine. C'est une catastrophe d'\emph{origine naturelle}, sous-type \emph{géologique}, d'échelle locale et d'apparition brutale.
\item (b) \textbf{Inondations de Yaoundé} : le déclencheur est naturel (pluies intenses de la saison humide), donc le groupe est celui des catastrophes d'\emph{origine naturelle} (sous-type climatique). Mais les facteurs aggravants sont humains : urbanisation anarchique des bas-fonds (zones inondables), déforestation des collines, colmatage des caniveaux par les déchets. On conclut : catastrophe naturelle \emph{aggravée par l'action humaine} (cas intermédiaire).
\item (c) \textbf{Eséka} : le déclencheur est lié à l'exploitation ferroviaire (vitesse, état de la voie ou du matériel, organisation) : c'est une catastrophe d'\emph{origine humaine}, sous-type \emph{accident de transport}, d'échelle locale et d'apparition brutale.
\end{itemize}
\end{exoresolu}

\begin{aretenir}[L'essentiel de la classification]
\begin{itemize}
\item Le critère fondamental est \textbf{l'origine du déclencheur} : naturelle (géologique, climatique, biologique, cosmique) ou humaine (industrielle, transports, sanitaire, conflits, structures).
\item De nombreux événements sont \textbf{mixtes} : déclencheur naturel + facteurs aggravants humains (inondations urbaines, glissements de terrain sur pentes construites).
\item Une catastrophe peut en déclencher une autre : c'est la \textbf{catastrophe en cascade} (séisme $\rightarrow$ tsunami $\rightarrow$ accident nucléaire de Fukushima ; on parle aussi de NaTech).
\item Les critères secondaires (échelle spatiale, durée brutale ou rampante, prévisibilité) affinent la description sans remplacer le critère d'origine.
\item Méthode en trois temps : déclencheur $\rightarrow$ nature du déclencheur $\rightarrow$ facteurs aggravants.
\end{itemize}
\end{aretenir}

Maintenant que nous savons classer les catastrophes, la section suivante examinera en détail le premier grand groupe : les catastrophes d'origine naturelle, leurs mécanismes et leurs manifestations au Cameroun et dans le monde.

\coursec{Les catastrophes d'origine naturelle}

\noindent Après avoir établi une classification générale des catastrophes (aléa, vulnérabilité, risque) et distingué les origines naturelle, humaine et mixte, nous nous penchons maintenant sur celles qui trouvent leur origine dans la dynamique propre de la planète Terre, sans action directe de l'être humain comme cause déclenchante initiale. Ces phénomènes, bien que naturels, voient leur impact dramatique souvent amplifié par la vulnérabilité des sociétés (densité de population, urbanisation non maîtrisée, pauvreté, défaut d'aménagement). Nous allons les passer en revue par grandes catégories, en illustrant chaque mécanisme par des exemples vécus au Cameroun ou dans le monde, et en adoptant la rigueur de la démarche d'investigation scientifique.

\courssub{Catastrophes géologiques : séismes, volcanisme et tsunamis}

\textbf{Rappel sur la tectonique des plaques (savoir admis de Première).} La lithosphère est fragmentée en plaques rigides qui se déplacent les unes par rapport aux autres (quelques cm/an). Aux frontières de plaques (divergentes, convergentes, transformantes), s'accumulent des contraintes élastiques dans les roches. Leur rupture brutale libère une énergie considérable : c'est l'origine des \cle{séismes} et de l'essentiel du \cle{volcanisme}.

\begin{experience}[Mécanisme du séisme : rebond élastique]
\textbf{Observation :} Après un séisme, on observe des décalages le long de failles (ex. faille de San Andreas). Les roches de part et d'autre de la faille ont accumulé une déformation élastique.
\textbf{Hypothèse :} La rupture se produit lorsque la contrainte dépasse la résistance de la roche. L'énergie élastique stockée est libérée sous forme d'ondes sismiques.
\textbf{Modélisation (expérience de cours) :} On place deux blocs de bois (représentant les plaques) sur une table, reliés par un élastique à un dynamomètre. On tire lentement : le bloc ne bouge pas d'abord (accumulation contrainte), puis glisse brutalement (rupture) quand la force de frottement est vaincue. Le dynamomètre chute : libération d'énergie.
\textbf{Interprétation :} Le glissement brutal correspond au séisme. Les ondes sismiques (P, S, de surface) se propagent depuis le \cle{foyer} (hypocentre) vers l'\cle{épicentre} (projection à la surface).
\textbf{Conclusion :} La magnitude (échelle de Richter ou de moment $M_w$) quantifie l'énergie libérée ; l'intensité (échelle MSK ou EMS) décrit les effets en surface (dégâts).
\end{experience}

\noindent Au Cameroun, le risque sismique est faible à modéré, lié aux failles de la chaîne volcanique (ligne du Cameroun) et à la faille du Bénoué. Cependant, les séismes lointains (ex. Haïti 2010, Turquie 2023) rappellent l'importance de la construction parasismique.

\begin{exemplebox}[Le volcanisme camerounais : le Mont Cameroun]
Le \cle{Mont Cameroun} (4\,100\,m), point culminant d'Afrique de l'Ouest, est un volcan de point chaud (\emph{hotspot}) aligné sur la \emph{ligne du Cameroun}. C'est le volcan le plus actif de la région.
\begin{itemize}
    \item \textbf{Éruption de 1999 :} Deux fissures s'ouvrent sur le flanc Sud à 2\,650\,m et 1\,400\,m d'altitude. Coulées de lave fluide (basalte) détruisant des plantations, approche de Buea. Évacuation de populations. Durée : ~3 semaines.
    \item \textbf{Éruption de 2000 :} Fissure à 1\,900\,m. Coulées plus rapides, menace directe sur Buea. Arrêt miraculeux à quelques centaines de mètres de la ville.
\end{itemize}
Ces éruptions de type \emph{hawaïen} (effusives, laves fluides, peu explosives) permettent l'évacuation, mais détruisent tout sur leur passage (infrastructures, cultures, écosystèmes).
\end{exemplebox}

\begin{savaistu}[Le Mont Cameroun : sentinelle de l'Afrique de l'Ouest]
Avec ses éruptions récentes (1982, 1999, 2000, 2012), le Mont Cameroun est le volcan le plus actif d'Afrique de l'Ouest. Il est surveillé en continu par l'Observatoire Volcanologique du Mont Cameroun (OVMC) de l'Université de Buea : sismomètres, GPS (déformation du sol), analyse des gaz (\ce{SO2}, \ce{CO2}), caméras thermiques. Cette surveillance permet la \cle{prévision} à court terme (jours/semaines) et la gestion de crise.
\end{savaistu}

\noindent Les \cle{tsunamis} (vagues de raz-de-marée) sont souvent la conséquence de séismes sous-marins (subduction) ou d'effondrements volcaniques (ex. Krakatoa 1883, Anak Krakatoa 2018). Le déplacement brutal d'un grand volume d'eau génère une onde qui se propage à grande vitesse (~700 km/h en haute mer) et s'amplifie en arrivant sur la côte (hauteur de plusieurs mètres à dizaines de mètres). Le littoral camerounais (Kribi, Limbe, Douala) est exposé aux tsunamis générés dans le Golfe de Guinée ou l'Atlantique Sud, bien que le risque soit considéré comme modéré.

\courssub{Les explosions limniques : le cas emblématique du lac Nyos}

\begin{experience}[Enquête sur la catastrophe du lac Nyos (21 août 1986) — Démarche d'investigation]
\textbf{Observations (témoignages, constat médico-légal, mesures physico-chimiques) :}
\begin{itemize}
    \item Nuit du 21 au 22 août 1986 : un bruit sourd entendu, odeur de poudre à canon / œuf pourri (\ce{H2S} traces).
    \item Un nuage blanc, bas, suit les vallées (direction Nord) sur ~25 km.
    \item 1\,746 morts humains, 3\,500 têtes de bétail. Cadavres sans blessures externes, couleur rosée (lividité), signes d'asphyxie aiguë. Survivants : troubles neurologiques, brûlures cutanées légères (acidité), coma.
    \item Lac Nyos : eau devenue trouble, rougeâtre (oxyde de fer \ce{Fe^{3+}}), niveau baissé d'~1 m, vague de 25 m sur les rives. Température de surface inchangée.
    \item Profil chimique (mesures 1986-1989) : concentration en \ce{CO2} dissous croissante avec la profondeur, saturation atteinte au fond (~200 m). Concentration en \ce{CO2} dans l'air au-dessus du lac : normale quelques jours après.
\end{itemize}
\textbf{Hypothèses initiales :} Éruption volcanique sous-lacustre ? Éboulement ? Attaque chimique (gaz de combat) ?
\textbf{Preuves décisives (équipe internationale, 1986-1989) :}
\begin{enumerate}
    \item Absence de cendres, de lave, de séismes précurseurs $\rightarrow$ exclut éruption magmatique classique.
    \item Isotopes du carbone ($\delta^{13}$C) du \ce{CO2} lacustre : signature \emph{magmatique} (manteau), non organique. Le \ce{CO2} remonte par des fractures depuis une chambre magmatique peu profonde.
    \item Stratification stable du lac (méromictique) : pas de mélange saisonnier (pas de \emph{turnover}). Le \ce{CO2} s'accumule indéfiniment au fond sous pression hydrostatique (~20 bar à 200 m).
    \item Dégazage brutal : un déclencheur (éboulement, pluie froide, séisme mineur) provoque un soulèvement d'eau profonde $\rightarrow$ baisse de pression $\rightarrow$ nucléation bulles $\rightarrow$ emballement (effet champagne) $\rightarrow$ colonne d'eau/gaz jaillissant à la surface.
\end{enumerate}
\textbf{Interprétation :} Le lac Nyos est un \cle{lac méromictique} volcanique. Le \ce{CO2} d'origine mantellique se dissout dans les eaux profondes froides. La stratification empêche l'évacuation naturelle. Un événement mineur déclenche un \cle{dégazage limnique} catastrophique.
\textbf{Conclusion :} Le nuage mortel était du \ce{CO2} (densité 1,5 $\times$ air) + gouttelettes d'eau + \ce{H2S} traces. Il a épousé le relief (vallées), déplaçant l'air (\ce{O2}), causant la mort par \cle{asphyxie hypoxique} (pas empoisonnement).
\end{experience}

\begin{popfigure}
\begin{tikzpicture}[scale=0.9, every node/.style={font=\small}]
    % Echelle profondeur
    \draw[->] (-1,0) -- (-1,8) node[left, pos=0.5, rotate=90] {Profondeur (m)};
    \foreach \y/\label in {0/0, 2/50, 4/100, 6/150, 8/200} {
        \draw (-1.1,\y) -- (-0.9,\y) node[left] {\label};
    }
    
    % Coupe du lac
    % Rives
    \draw[thick, fill=brown!30] (-4,0) -- (-0.5,0) -- (-0.5,8) -- (-4,8) -- cycle;
    \draw[thick, fill=brown!30] (4.5,0) -- (8,0) -- (8,8) -- (4.5,8) -- cycle;
    
    % Eau - Epilimnion (0-50m)
    \draw[fill=popblueL!50, opacity=0.6] (-0.5,0) -- (4.5,0) -- (4.5,2) -- (-0.5,2) -- cycle;
    \node[align=center] at (2,1) {\textbf{Epilimnion}\\(0--50 m)\\\emph{Eau oxygénée, mélangée}\\$T \approx 24^\circ$C};
    
    % Thermocline / Metalimnion
    \draw[fill=popblue!30, opacity=0.7] (-0.5,2) -- (4.5,2) -- (4.5,3.5) -- (-0.5,3.5) -- cycle;
    \node[align=center] at (2,2.75) {\textbf{Métalimnion / Thermocline}\\(50--90 m)\\\emph{Barrière à la convection}};
    
    % Hypolimnion (90-200m) - Eau profonde saturée en CO2
    \draw[fill=popdark!20, opacity=0.8] (-0.5,3.5) -- (4.5,3.5) -- (4.5,8) -- (-0.5,8) -- cycle;
    \node[align=center, text width=4.5cm] at (2,5.75) {\textbf{Hypolimnion}\\(>90 m)\\\emph{Eau stagnante, anoxique, riche en \ce{Fe^{2+}}, \ce{NH4+}}\\\textbf{\ce{CO2} dissous : saturation $\approx$ 20 bar}\\(~500 mmol/L)};
    
    % Bulles CO2 au fond (source magmatique)
    \foreach \i in {1,...,15} {
        \pgfmathsetmacro{\x}{-0.2 + rand*4.4}
        \pgfmathsetmacro{\y}{7.2 + rand*0.6}
        \pgfmathsetmacro{\r}{0.05 + rand*0.05}
        \draw[fill=poporange!80!popdark, opacity=0.8] (\x,\y) circle (\r);
    }
    \node[poporange, align=center] at (2,7.5) {\textbf{Source magmatique}\\(\ce{CO2}, chaleur)};
    
    % Tuyau de dégazage (après 2001)
    \draw[very thick, popgreen] (2,0) -- (2,8);
    \draw[popgreen, thick] (2,7.8) -- (2,8.2) -- (2.5,8.2) -- (2.5,7.8) -- cycle; % Tête
    \node[popgreen, right] at (2.6,7.5) {Tuyau de dégazage (2001/2011)};
    
    % Dégazage brutal (flèches)
    \draw[->, very thick, poporange] (2,3.5) -- (2,1.5) node[midway, right] {\textbf{Dégazage brutal}};
    \draw[->, very thick, poporange] (2,1.5) .. controls (3,1) and (5,0.5) .. (6,0) node[right, align=left] {Nuage de \ce{CO2}\\(densité > air)\\\emph{Écoulement vallonné}};
    
    % Légende couleurs (positionnée dans la zone rives gauche)
    \begin{scope}[xshift=-4.5cm, yshift=-0.5cm]
        \draw[fill=popblueL!50] (0,0) rectangle (0.5,0.3) node[right, xshift=0.2] {Eau de surface};
        \draw[fill=popdark!20] (0,-0.5) rectangle (0.5,-0.2) node[right, xshift=0.2] {Eau profonde saturée \ce{CO2}};
        \draw[fill=poporange!80!popdark] (0,-1) rectangle (0.5,-0.7) node[right, xshift=0.2] {Bulles \ce{CO2} magmatique};
    \end{scope}
\end{tikzpicture}
\legende{Coupe schématique du lac Nyos montrant la stratification méromictique stable, l'accumulation de \ce{CO2} d'origine magmatique dans l'hypolimnion (eau profonde), le tuyau de dégazage artificiel installé en 2001 (renforcé 2011) et le mécanisme du dégazage brutal (explosion limnique) générant un nuage de \ce{CO2} dense s'écoulant dans les vallées.}
\end{popfigure}

\begin{definition}[Explosion limnique]
Une \textbf{explosion limnique} (ou dégazage limnique) est le relargage brutal et massif de gaz dissous (principalement du dioxyde de carbone \ce{CO2}, parfois du méthane \ce{CH4} ou du sulfure d'hydrogène \ce{H2S}) accumulé dans les eaux profondes d'un lac volcanique méromictique (stratifié de façon permanente), suite à un dépassement de la saturation ou à un déclencheur mécanique (éboulement, séisme, pluie froide).
\end{definition}

\begin{propriete}[Comportement du \ce{CO2} et mécanisme létal (admis)]
Le dioxyde de carbone (\ce{CO2}) a une densité molaire de 44 g/mol, contre ~29 g/mol pour l'air sec. Sa densité relative par rapport à l'air est donc d'environ \textbf{1,5}. Lors d'une explosion limnique, le \ce{CO2} libéré forme une \textbf{nappe gazeuse rasante} qui s'écoule par gravité dans les vallées et les bas-fonds (comportement d'un fluide plus dense). Il \textbf{ne tue pas par toxicité chimique} (le \ce{CO2} n'est pas un poison au sens strict), mais par \textbf{asphyxie par déplacement de l'oxygène} (hypoxie sévère) : la concentration en \ce{O2} chute en dessous du seuil vital (~10-12 \%), entraînant perte de conscience rapide et mort par arrêt cardio-respiratoire. La présence de traces de \ce{H2S} (odeur œuf pourri) peut aggraver le tableau (inhibition cytochrome c oxydase), mais le \ce{CO2} est la cause majeure.
\end{propriete}

\begin{exemplebox}[Données chiffrées : Lac Nyos et Lac Monoun]
\begin{itemize}
    \item \textbf{Lac Monoun (1984) :} 37 morts. Premier signal d'alerte méconnu.
    \item \textbf{Lac Nyos (21 août 1986) :} Libération estimée à \textbf{1,6 million de tonnes de \ce{CO2}} (volume ~0,8 km$^3$ à CNTP) en quelques heures. Vague de 25 m sur les rives. Nuage descendu à 50-100 km/h dans la vallée, tuant jusqu'à 25 km.
    \item \textbf{Solution technique :} Installation d'un \textbf{tuyau de dégazage} (auto-amorçage par la poussée d'Archimède sur la colonne eau-gaz) en \textbf{2001}, puis deux tuyaux supplémentaires en \textbf{2011}. Objectif : maintenir la pression partielle de \ce{CO2} au fond en dessous de la saturation, empêcher toute nouvelle accumulation critique. Surveillance continue (capteurs pression, température, conductivité).
\end{itemize}
\end{exemplebox}

\begin{attention}[Erreur fréquente : « Gaz toxique » vs « Asphyxiant »]
Ne jamais écrire : « Le \ce{CO2} est un gaz toxique qui a empoisonné les villageois. »
\textbf{Correction rigoureuse :} Le \ce{CO2} est un gaz \textbf{asphyxiant simple} (inerte, non toxique chimiquement aux concentrations ambiantes). Il tue par \textbf{privation d'oxygène} (hypoxie) due à son accumulation au ras du sol (densité > air). La distinction est cruciale : un gaz toxique (ex. \ce{CN-}, Sarin) agit à très faible dose par mécanisme biochimique ; le \ce{CO2} agit par effet physique de dilution de l'\ce{O2} à haute concentration (> 30-50 \%). Les survivants du lac Nyos n'ont pas de séquelles toxicologiques spécifiques, seulement des séquelles neurologiques post-anoxiques.
\end{attention}

\courssub{Catastrophes climatiques et météorologiques}

Ces catastrophes résultent de phénomènes atmosphériques extrêmes. Leur fréquence et/ou intensité tendent à augmenter dans le contexte du \cle{changement climatique} (réchauffement global).

\begin{itemize}
    \item \textbf{Cyclones tropicaux (ouragans, typhons) :} Systèmes dépressionnaires organisés sur océans chauds ($T > 26,5^\circ$C), vents > 118 km/h. Dégâts : vents violents, pluies diluviennes, \cle{surcote} (élévation du niveau marin). Rarement touchent directement le Cameroun (position équatoriale, zone de convergence intertropicale), mais résidus dépressionnaires apportent de fortes pluies.
    \item \textbf{Tornades :} Tourbillons violents, échelle locale (quelques 100 m), vents > 300 km/h. Associés à orages supercellulaires. Rares au Cameroun.
    \item \textbf{Orages violents / Tempêtes :} Grêle, rafales descendantes (\emph{downbursts}), foudre. Fréquents en saison des pluies (mars-octobre). Risque : électrocution, incendies, destructions cultures/habitats précaires.
    \item \textbf{Vagues de chaleur (canicules) :} Températures maximales > seuils saisonniers pendant plusieurs jours/nuits. Impact sanitaire : coups de chaleur, déshydratation, surmortalité personnes âgées/malades chroniques. Aggravées par îlots de chaleur urbains (Yaoundé, Douala).
    \item \textbf{Sécheresses et désertification (Nord Cameroun - Régions Nord, Extrême-Nord, Adamaoua) :} Déficit pluviométrique prolongé (pluviométrie < 700 mm/an, irrégularité interannuelle forte). Conséquences : baisse nappes phréatiques, tarissement cours d'eau (Logone, Chari), perte récoltes (mil, sorgho, maïs), mortalité bétail, insécurité alimentaire, migrations, conflits agriculteurs/éleveurs. La \cle{désertification} (dégradation des terres en zones arides/semi-arides par facteurs climatiques + activités humaines : surpâturage, déboisement) avance vers le sud.
\end{itemize}

\begin{popfigure}
\begin{tikzpicture}
\begin{axis}[
    width=\linewidth,
    height=8cm,
    xlabel={Année},
    ylabel={Nombre de catastrophes naturelles déclarées},
    title={Évolution mondiale des catastrophes naturelles (type EM-DAT / CRED)},
    xmin=1970, xmax=2020,
    ymin=0, ymax=450,
    xtick={1970,1980,1990,2000,2010,2020},
    ytick={0,100,200,300,400},
    grid=major,
    grid style={popline},
    legend pos=north west,
    legend cell align=left,
    /pgf/number format/use comma,
    /pgf/number format/1000 sep={\,}
]
% Données approximatives type EM-DAT (total reported disasters)
\addplot[popcurve, thick, mark=*, mark size=2, popblue] coordinates {
    (1970, 78) (1975, 110) (1980, 155) (1985, 200) (1990, 250) (1995, 310) (2000, 350) (2005, 400) (2010, 380) (2015, 350) (2020, 410)
};
\addlegendentry{Total catastrophes (toutes causes)}

\addplot[popcurve, thick, mark=square*, mark size=2, poporange] coordinates {
    (1970, 20) (1975, 30) (1980, 45) (1985, 60) (1990, 85) (1995, 120) (2000, 160) (2005, 200) (2010, 190) (2015, 180) (2020, 220)
};
\addlegendentry{Climatiques / Météorologiques (inondations, tempêtes, sécheresses)}

\addplot[popcurve, thick, mark=triangle*, mark size=2, popgreen] coordinates {
    (1970, 15) (1975, 20) (1980, 25) (1985, 30) (1990, 35) (1995, 40) (2000, 45) (2005, 50) (2010, 55) (2015, 50) (2020, 55)
};
\addlegendentry{Géologiques (séismes, volcans, tsunamis)}
\end{axis}
\end{tikzpicture}
\legende{Évolution du nombre de catastrophes naturelles recensées dans la base EM-DAT (CRED, UCLouvain) depuis 1970. On observe une forte augmentation du nombre total d'événements déclarés, due pour partie à une meilleure déclaration (effet de surveillance) et pour partie à l'augmentation réelle des événements hydro-météorologiques (inondations, tempêtes, vagues de chaleur) liée au changement climatique. Les catastrophes géologiques restent stables (aléas tectoniques constants).}
\end{popfigure}

\courssub{Catastrophes hydrologiques : inondations et crues subites}

Les inondations sont la catastrophe la plus fréquente et la plus coûteuse au Cameroun.

\begin{itemize}
    \item \textbf{Inondations urbaines (Yaoundé, Douala, Bafoussam, Garoua) :} Causées par pluies intenses + \textbf{imperméabilisation des sols} (béton, bitume) + \textbf{déficit/défaut de réseaux d'assainissement} (caniveaux bouchés par déchets plastiques, sédiments) + \textbf{occupation illicite des zones inondables} (lits majeurs, marécages asséchés). Exemple : Douala, juin 2023 : 200 mm en 24h, quartiers entiers submergés, pertes humaines, choléra.
    \item \textbf{Inondations fluviales (Plaine du Logone, Vallée du Bénoué, Fleuve Sanaga) :} Crue lente, prévisible (semaines), due à pluies amont (bassin versant). Population adaptée (architecture sur pilotis, cultures de décrue), mais crues exceptionnelles (ex. 2012, 2022 Logone) dépassent seuils historiques, déplacent des dizaines de milliers de personnes (réfugiés climatiques internes).
    \item \textbf{Crues subites (\emph{flash floods}) :} Bassins versants pentus, petits, imperméables (roche nue, sols brûlés). Temps de concentration < 6h. Dangereuses car imprévisibles à très court terme. Exemple : Monts Mandara, Monts de l'Ouest.
\end{itemize}

\courssub{Glissements de terrain et effondrements}

\begin{exemplebox}[Bafoussam, nuit du 28 au 29 octobre 2019]
Après plusieurs jours de pluies diluviennes (cumuls > 300 mm en 72h), un pan de colline s'effondre sur le quartier \emph{Gouache} (pente > 35\%, sols latéritiques profonds, altération intense). Bilan : \textbf{43 morts}, dizaines de maisons détruites, route coupée.
\textbf{Facteurs aggravants identifiés (rapport MINATD/ONU) :}
\begin{enumerate}
    \item \textbf{Géologie :} Roche mère fracturée, saprolite épaisse (argiles gonflantes), nappe phréatique haute.
    \item \textbf{Climat :} Pluies exceptionnelles (saturation en eau $\rightarrow$ perte cohésion, augmentation pression interstitielle).
    \item \textbf{Anthropique :} \textbf{Déforestation} (coupe bois de chauffe, extension cultures), \textbf{urbanisation anarchique} (construction en bas de pente, talutage sans soutènement), absence drainage.
\end{enumerate}
\end{exemplebox}

\begin{propriete}[Critères de stabilité des versants (simplifié)]
La stabilité d'un versant dépend de l'équilibre entre \textbf{forces motrices} (poids du sol $P$, poussée eau $U$) et \textbf{forces résistantes} (cohésion $c$, frottement interne $\tan \phi$, racinaires végétation). Le \textbf{Facteur de Sécurité} $F_s = \frac{\text{Résistance}}{\text{Mobilisation}}$. Si $F_s < 1$, rupture. La déforestation supprime le renforcement racinaire ($c_{racines}$) et l'interception pluviométrique $\rightarrow$ infiltration maximale $\rightarrow$ saturation rapide $\rightarrow$ chute $F_s$.
\end{propriete}

\courssub{Feux de brousse d'origine naturelle}

\begin{definition}[Feu de brousse naturel]
Incendie de végétation (savane, forêt, brousse) déclenché par la \textbf{foudre} (éclairs « secs » sans pluie ou avec pluie insuffisante pour éteindre le départ de feu), en période sèche (novembre-mars au Nord, décembre-février au Sud). À distinguer strictement des feux \textbf{anthropiques} (brûlis agricoles, chasse, incendies criminels, mégots, lignes électriques) qui représentent > 90 \% des départs de feu au Cameroun.
\end{definition}

\noindent Rôle écologique paradoxal : en savane, le feu naturel (ou précoce contrôlé) maintient l'écosystème (régénération herbes, contrôle ligneux, cycle nutriments). En forêt dense humide, le feu est toujours destructeur (non adapté).

\courssub{Épidémies et épizooties : catastrophes biologiques naturelles}

Le programme les classe parmi les catastrophes d'origine naturelle lorsque le pathogène émerge ou réémerge sans manipulation humaine (zoonoses, vecteurs climatiques).

\begin{itemize}
    \item \textbf{Choléra (\emph{Vibrio cholerae}) :} Maladie diarrhéique aiguë, transmission fécale-orale (eau/aliments souillés). \textbf{Épidémies récurrentes} au Cameroun (Nord, Extrême-Nord, Littoral, Sud-Ouest). Déclenchées/favorisées par : inondations (contamination nappes/puits), sécheresses (concentration populations points d'eau), insalubrité urbaine, déplacements populations (conflits Boko Haram, crise NOSO). Exemple : Épidémie 2021-2023 : > 15\,000 cas, > 300 morts.
    \item \textbf{Méningite à méningocoque (\emph{Neisseria meningitidis} sérogroupe A, C, W, X) :} \textbf{Ceinture de la méningite} (Sahel) traverse l'Extrême-Nord et le Nord. Pic saisonnier : \textbf{saison sèche chaude et poussiéreuse (janvier-mai)}. Le vent d'Harmattan (poussières) irrite la muqueuse pharyngée $\rightarrow$ facilite invasion bactérienne. Transmission aérienne (gouttelettes). Létalité 10-15 \% sans traitement, séquelles neurologiques (surdi, épilepsie). Vaccination conjuguée (MenAfriVac sérogroupe A) a fait chuter les épidémies de groupe A, mais émergence sérogroupes C, W, X.
    \item \textbf{Autres :} Fièvre jaune (sylvatique $\rightarrow$ urbaine via \emph{Aedes}), Ebola (réservoir chauves-souris, zoonose), Mpox (variole du singe), grippe aviaire (épizootie $\rightarrow$ risque zoonotique).
\end{itemize}

\begin{methode}[Conduite à tenir face à une alerte épidémique (choléra/méningite)]
\begin{enumerate}
    \item \textbf{Surveillance / Alerte :} Déclaration obligatoire (maladies à déclaration obligatoire - MDO) au District de Santé. Seuil d'alerte : 5 cas/semaine (choléra) ou 10 cas/100\,000 hab/semaine (méningite).
    \item \textbf{Confirmation biologique :} Prélèvement (selles/rectal pour choléra, LCR pour méningite) $\rightarrow$ Labo national (Centre Pasteur Yaoundé/Garoua).
    \item \textbf{Riposte :} Traitement gratuit (CTA - Centre de Traitement du Choléra : réhydratation ORS/IV, antibio Doxycycline/Azithromycine ; Méningite : Ceftriaxone IV d'urgence). Vaccination de masse réactive (si vaccin dispo). WASH (Eau, Hygiène, Assainissement) : chloration points d'eau, latrines, sensibilisation.
    \item \textbf{Communication :} Messages clairs (signes d'alerte, où aller, gratuité), lutte contre rumeurs/stigmatisation.
\end{enumerate}
\end{methode}

\courssub{Cartographie des zones à risques naturels au Cameroun}

\begin{popfigure}
\begin{tikzpicture}[scale=0.85, every node/.style={font=\footnotesize}]
    % Contour simplifié Cameroun (approximatif)
    \draw[thick, popdark] 
        (0,0) -- (1,0.2) -- (2,0.5) -- (3,1) -- (3.5,1.5) -- (3.8,2) -- (3.5,2.5) -- (3,3) -- (2.5,3.5) -- (2,4) -- (1.5,4.5) -- (1,5) -- (0.5,5.5) -- (0,6) -- (-0.5,5.5) -- (-1,5) -- (-1.5,4.5) -- (-2,4) -- (-2.5,3.5) -- (-3,3) -- (-3.5,2.5) -- (-3.8,2) -- (-3.5,1.5) -- (-3,1) -- (-2,0.5) -- (-1,0.2) -- cycle;
    
    % Villes repères
    \node[circle, fill=popblue, inner sep=1.5pt, label=above:\textcolor{popdark}{Douala}] (Douala) at (-1.2, 1.8) {};
    \node[circle, fill=popblue, inner sep=1.5pt, label=above:\textcolor{popdark}{Yaoundé}] (Yaounde) at (-0.5, 2.5) {};
    \node[circle, fill=popblue, inner sep=1.5pt, label=right:\textcolor{popdark}{Bafoussam}] (Bafoussam) at (0.5, 2.8) {};
    \node[circle, fill=popblue, inner sep=1.5pt, label=above:\textcolor{popdark}{Garoua}] (Garoua) at (-1.5, 4.5) {};
    \node[circle, fill=popblue, inner sep=1.5pt, label=right:\textcolor{popdark}{Maroua}] (Maroua) at (0.5, 5.0) {};
    \node[circle, fill=popblue, inner sep=1.5pt, label=left:\textcolor{popdark}{Buea}] (Buea) at (-2.8, 1.5) {};
    \node[circle, fill=popblue, inner sep=1.5pt, label=left:\textcolor{popdark}{Kribi}] (Kribi) at (-3.2, 0.5) {};
    
    % Zones de risque - Volcanisme / Limnique (Sud-Ouest)
    \draw[fill=poporange!30, opacity=0.6] (-3.5, 1.2) circle (0.6);
    \node[poporange, align=center] at (-3.5, 1.2) {\textbf{Volcanisme}\\(Mont Cameroun)\\\textbf{Risque limnique}\\(Lacs Nyos, Monoun)};
    
    % Inondations urbaines (Douala, Yaoundé)
    \draw[fill=popblue!30, opacity=0.6] (-1.2, 1.8) ellipse (0.5 and 0.3);
    \node[popblue, align=center] at (-1.2, 1.4) {\textbf{Inondations urbaines}\\(Douala, Yaoundé,\\Bafoussam)};
    
    % Inondations fluviales (Logone, Bénoué)
    \draw[fill=popblue!20!popgreen, opacity=0.6] (-1.5, 4.5) ellipse (1.2 and 0.5);
    \node[popblue!70!popdark, align=center] at (-1.5, 4.0) {\textbf{Inondations fluviales}\\(Logone, Chari, Bénoué)};
    
    % Sécheresse / Désertification / Méningite (Grand Nord)
    \draw[fill=popgold!40, opacity=0.6] (-0.5, 5.5) rectangle (2.5, 6.5);
    \node[popdark, align=center] at (1, 6.0) {\textbf{Sécheresse / Désertification}\\\textbf{Méningite (Ceinture)}\\\textbf{Insécurité alimentaire}};
    
    % Glissements de terrain (Ouest, Nord-Ouest, Sud-Ouest)
    \draw[fill=poppurple!30, opacity=0.6] (0.5, 2.8) ellipse (0.8 and 0.4);
    \node[poppurple, align=center] at (0.5, 2.4) {\textbf{Glissements terrain}\\(Bafoussam, Bamenda,\\Kumba, Monts Mandara)};
    
    % Risque sismique modéré (Ligne du Cameroun)
    \draw[dashed, thick, poppink] (-3.5, 1.2) -- (-2.5, 2.5) -- (-1, 3.2) -- (0.5, 3.8) -- (2, 4.5);
    \node[poppink, rotate=30] at (-1.5, 3.0) {Faille / Ligne du Cameroun (Risque sismique modéré)};
    
    % Légende (placée en bas à droite de la carte)
    \begin{scope}[xshift=4.5cm, yshift=-1.5cm]
        \draw[fill=poporange!30] (0,0) rectangle (0.4,0.2) node[right, xshift=0.2] {Volcanisme / Limnique};
        \draw[fill=popblue!30] (0,-0.4) rectangle (0.4,-0.2) node[right, xshift=0.2] {Inondations urbaines};
        \draw[fill=popblue!20!popgreen] (0,-0.8) rectangle (0.4,-0.6) node[right, xshift=0.2] {Inondations fluviales};
        \draw[fill=popgold!40] (0,-1.2) rectangle (0.4,-1.0) node[right, xshift=0.2] {Sécheresse / Désertification / Méningite};
        \draw[fill=poppurple!30] (0,-1.6) rectangle (0.4,-1.4) node[right, xshift=0.2] {Glissements de terrain};
        \draw[dashed, thick, poppink] (0,-1.9) -- (0.4,-1.9) node[right, xshift=0.2] {Faille / Risque sismique};
    \end{scope}
\end{tikzpicture}
\legende{Carte simplifiée des principaux aléas naturels au Cameroun. Le pays présente une diversité de risques : volcanisme et risque limnique unique au monde sur la Ligne du Cameroun (Sud-Ouest) ; inondations urbaines chroniques (Douala, Yaoundé) et fluviales saisonnières (Logone, Bénoué) ; sécheresse, désertification et méningite dans le Grand Nord ; glissements de terrain dans les zones de relief (Ouest, Nord-Ouest) ; risque sismique modéré le long de la Ligne du Cameroun.}
\end{popfigure}

\courssub{Les gestes qui sauvent face aux aléas naturels majeurs}

\begin{methode}[Conduite à tenir et gestes qui sauvent (savoir-faire Bac)]
\begin{enumerate}
    \item \textbf{Séisme (pendant) :} \textbf{Se baisser, s'abriter, s'agripper} (Drop, Cover, Hold On). Sous table solide, loin des fenêtres/miroirs. Pas d'ascenseur. Dehors : s'éloigner bâtiments, poteaux, arbres.
    \item \textbf{Séisme (après) :} Vérifier blessures, couper gaz/électricité si fuite suspectée. Écouter radio (consignes). Prévoir répliques. Ne pas téléphoner (saturer réseaux) sauf urgence vitale.
    \item \textbf{Inondation / Crue subite :} \textbf{Ne jamais marcher ni conduire dans l'eau} (30 cm emportent voiture, 15 cm font tomber adulte). Rejoindre point haut. Couper électricité si eau menace installation. Éviter contact eau souillée (choléra, leptospirose).
    \item \textbf{Explosion limnique (lac Nyos type) :} \textbf{Monter en altitude immédiatement} (colline, étage élevé). \ce{CO2} plus dense que l'air = s'accumule en bas. Ne pas s'allonger. Respirer dans tissu humide si passage nuage inévitable. Alerter voisins en criant « Gaz ! Montez ! ».
    \item \textbf{Glissement de terrain :} Signes avant-coureurs : fissures sol, arbres penchés, eau trouble sources, bruits souterrains. \textbf{Évacuer immédiatement} perpendiculairement à la direction du glissement. Ne pas retourner chercher affaires.
    \item \textbf{Éruption volcanique (cendres) :} Masque (ou tissu humide) sur nez/bouche. Lunettes (pas lentilles). Toit : enlever cendres (poids + corrosion). Fermer portes/fenêtres, boucher aérations.
    \item \textbf{Canicule :} Boire 1,5-2 L eau/jour (sans soif). Mouiller corps. Fermer volets jour, aérer nuit. Surveiller personnes vulnérables (âgés, nourrissons, malades chroniques). Signes alerte : crampes, épuisement (sueurs froides, nausées), coup de chaleur (fièvre > 40°C, confusion, peau sèche/chaude) = URGENCE VITALE (appeler 112/119, refroidir activement).
\end{enumerate}
\end{methode}

\courssub{Synthèse et exercice résolu}

\begin{aretenir}[Les catastrophes d'origine naturelle : points clés]
\begin{itemize}
    \item \textbf{Géologiques :} Séismes (failles, plaques), Volcanisme (Mont Cameroun : effusif, surveillé), Tsunamis (côte).
    \item \textbf{Limniques :} Explosion limnique = dégazage brutal \ce{CO2} lac méromictique (Nyos, Monoun). Mécanisme : saturation fond $\rightarrow$ déclencheur $\rightarrow$ emballement. \ce{CO2} dense $\rightarrow$ nappe rasante $\rightarrow$ asphyxie (pas poison).
    \item \textbf{Climatiques/Météorologiques :} Cyclones, tornades, orages, vagues de chaleur, \textbf{sécheresses/désertification (Grand Nord)}.
    \item \textbf{Hydrologiques :} Inondations urbaines (imperméabilisation, déchets), fluviales (Logone), crues subites.
    \item \textbf{Gravitaires :} Glissements terrain (Bafoussam 2019) : pluie + pente + déforestation + urbanisation.
    \item \textbf{Biologiques :} Choléra (eau, inondations), Méningite (Harmattan, ceinture Sahel), Fièvre jaune, Ebola, Mpox.
    \item \textbf{Gestion :} Prévision (surveillance volcanique, météo, épidémiologique), Prédiction (modélisation court terme), Prévention (aménagement, vaccination, WASH, dégazage Nyos, reboisement, construction parasismique).
\end{itemize}
\end{aretenir}

\begin{exoresolu}[Classification de catastrophes camerounaises]
\textbf{Énoncé :} Classer les événements suivants selon leur origine principale (Naturelle / Humaine / Mixte) et leur catégorie (Géologique, Limnique, Climatique, Hydrologique, Gravitaire, Biologique). Justifier brièvement.
\begin{enumerate}
    \item Éruption du Mont Cameroun en 2000.
    \item Inondation de Douala en juin 2023 (pluies + caniveaux bouchés).
    \item Catastrophe du lac Nyos en 1986.
    \item Épidémie de choléra 2021-2023 dans le Nord.
    \item Glissement de terrain de Bafoussam en 2019.
    \item Sécheresse 2011-2012 dans l'Extrême-Nord (déficit pluie + surpâturage).
\end{enumerate}

\textbf{Solution détaillée :}
\begin{enumerate}
    \item \textbf{Origine : Naturelle.} \textbf{Catégorie : Géologique (Volcanisme).} Cause : dynamique magmatique point chaud. L'Homme n'a pas déclenché l'éruption (mais l'occupation du flanc aggrave la vulnérabilité).
    \item \textbf{Origine : Mixte (Naturelle + Humaine).} \textbf{Catégorie : Hydrologique (Inondation urbaine).} Déclencheur naturel : pluie intense. Facteurs aggravants majeurs anthropiques : imperméabilisation, absence drainage, déchets, habitat en zone inondable.
    \item \textbf{Origine : Naturelle.} \textbf{Catégorie : Limnique (Explosion limnique).} Accumulation \ce{CO2} magmatique sur décennies/siècles. Déclencheur probable naturel (éboulement, renversement). \textbf{Attention :} le tuyau de dégazage est une action humaine de \emph{prévention}, pas la cause.
    \item \textbf{Origine : Naturelle (pathogène) + Humaine (propagation).} \textbf{Catégorie : Biologique (Épidémie).} \emph{Vibrio cholerae} existe naturellement (réservoir aquatique). L'explosion épidémique est favorisée par facteurs humains : inondations, insalubrité, déplacements forcés, accès eau potable insuffisant.
    \item \textbf{Origine : Mixte.} \textbf{Catégorie : Gravitaire (Glissement de terrain).} Déclencheur naturel : pluies extrêmes (saturation). Facteurs de conditionnement anthropiques : déforestation (perte racinaires), urbanisation pente raide sans génie civil, talutage sauvage.
    \item \textbf{Origine : Mixte.} \textbf{Catégorie : Climatique (Sécheresse) $\rightarrow$ Désertification.} Déficit pluviométrique = aléa climatique naturel (variabilité + changement climatique). \textbf{Désertification} = dégradation terres = processus mixte (climat + surpâturage, déboisage, mauvaises pratiques agricoles).
\end{enumerate}
\end{exoresolu}

\begin{exercice}[Application : Analyse de risque local]
Choisis une catastrophe naturelle survenue dans ta région ou une ville que tu connais (ex. : inondation à Yaoundé, glissement à Bafoussam, sécheresse à Maroua, éruption vue au Mont Cameroun).
\begin{enumerate}
    \item Identifie l'\textbf{aléa} (phénomène physique), la \textbf{vulnérabilité} (facteurs humains, sociaux, économiques) et le \textbf{risque} (conjonction aléa $\times$ vulnérabilité).
    \item Propose \textbf{trois mesures de prévention structurelles} (ouvrages, aménagement) et \textbf{trois mesures non structurelles} (réglementation, éducation, assurance, système d'alerte) adaptées à ce contexte.
    \item Dessine un \textbf{schéma fonctionnel} (boucle de rétroaction) montrant comment le changement climatique peut amplifier cet aléa à l'horizon 2050.
\end{enumerate}
\end{exercice}

\begin{corrige}[Exercice : Analyse de risque local]
\textit{Exemple de réponse attendue pour une inondation urbaine à Yaoundé (quartier Mvog-Ada / Mfoundi) :}
\begin{enumerate}
    \item \textbf{Aléa :} Crue rapide rivière Mfoundi (pluie intense, bassin pentu, temps de concentration court). \textbf{Vulnérabilité :} Habitat précaire en lit majeur, caniveaux bouchés (déchets plastiques), absence plan d'urbanisme, population dense pauvre. \textbf{Risque :} Inondations récurrentes annuelles, pertes vies humaines, choléra, destruction biens.
    \item \textbf{Prévention structurelle :} (1) Curage/élargissement lit rivière + bassins de rétention amont. (2) Réseau d'assainissement pluvial séparatif dimensionné (retour période 10-20 ans). (3) Relogement populations zones à risque très fort (expropriation indemnisée).
    \item \textbf{Prévention non structurelle :} (1) Plan Local d'Urbanisme (PLU) contraignant : zone inconstructible = lit majeur + bande 50m. (2) Système d'alerte précoce (radar météo + sirènes + SMS) + exercices d'évacuation. (3) Gestion déchets : collecte porte-à-porte, tri, centres de valorisation $\rightarrow$ zéro déchet dans caniveaux. (4) Assurance catastrophe paramétrique pour ménages vulnérables.
    \item \textbf{Schéma fonctionnel (texte) :} Réchauffement global $\rightarrow$ + Évaporation $\rightarrow$ + Humidité atmosphérique $\rightarrow$ + Intensité pluies extrêmes (Clausius-Clapeyron ~7\%/°C) $\rightarrow$ + Fréquence crues centennales $\rightarrow$ + Dégâts (si vulnérabilité inchangée) $\rightarrow$ (Boucle positive) Pression sur budgets $\rightarrow$ - Investissement prévention $\rightarrow$ + Vulnérabilité. \textbf{Rupture possible :} Adaptation (aménagement, alerte, assurance) $\rightarrow$ - Vulnérabilité $\rightarrow$ Risque stabilisé.
\end{enumerate}
\end{corrige}

\coursec{Les catastrophes d'origine humaine}

Après avoir examiné les catastrophes dont le déclencheur est un phénomène naturel (séismes, éruptions volcaniques, cyclones, inondations), nous abordons maintenant celles qui trouvent leur cause principale dans l'activité humaine. Si la nature fournit parfois le cadre (un relief instable, un climat sec), c'est l'action, la négligence ou la volonté de l'Homme qui transforme une situation en drame. Comprendre ces catastrophes, c'est comprendre nos vulnérabilités technologiques, organisationnelles et sociales pour mieux les réduire.

\courssub{Catastrophes industrielles et technologiques}

\begin{definition}[Catastrophe technologique]
Une \cle{catastrophe technologique} (ou catastrophe industrielle) est un événement brutal, imprévu et d'ampleur limitée dans l'espace et le temps, provoqué par la défaillance d'une installation, d'un procédé ou d'une activité humaine (usine, transport de matières dangereuses, barrage, centrale nucléaire), et qui entraîne des dommages graves pour les personnes, les biens et l'environnement.
\end{definition}

Ces catastrophes résultent souvent de l'interaction entre une \cle{défaillance technique} (rupture de confinement, erreur de manœuvre, défaut de conception) et des \cle{facteurs organisationnels} (maintenance insuffisante, formation inadéquate, pression de production, non-respect des procédures de sécurité). La démarche d'investigation (type \emph{arbre des causes} ou méthode \emph{5M}) montre systématiquement qu'il n'y a pas une cause unique, mais une chaîne de défaillances.

\begin{exemplebox}[Exemples marquants de catastrophes technologiques]
\begin{itemize}
    \item \textbf{Tchernobyl (Ukraine, 26 avril 1986) :} Lors d'un test de sécurité mal préparé sur le réacteur n°4 de type RBMK, une excursion de puissance incontrôlée entraîne une explosion de vapeur puis un incendie du graphite. Le cœur du réacteur est découvert, rejetant dans l'atmosphère une quantité massive de radioéléments (iode-\num{131}, césium-\num{137}, strontium-\num{90}). Le nuage contamine l'Europe entière. Bilan : \num{31} morts directs (opérateurs, pompiers), des milliers de cancers thyroïdiens et leucémies attribués à l'exposition, zone d'exclusion de \SI{2600}{km^2} encore inhabitée.
    \item \textbf{Fukushima Daiichi (Japon, 11 mars 2011) :} Un séisme de magnitude \num{9,0} déclenche un tsunami de \SI{14}{\meter} à \SI{15}{\meter} qui submerge les groupes électrogènes de secours. Perte totale de l'alimentation électrique $\rightarrow$ arrêt du refroidissement des cœurs et des piscines de désactivation $\rightarrow$ fusion des cœurs (réacteurs 1, 2, 3) et explosions d'hydrogène $\rightarrow$ rejets radioactifs atmosphériques et marins. Évacuation de \num{160000} personnes.
    \item \textbf{Bhopal (Inde, 3 décembre 1984) :} Fuite de \SI{40}{\tonne} de \ce{CH3NCO} (isocyanate de méthyle, MIC) dans l'usine Union Carbide. Entrée d'eau dans un réservoir de stockage $\rightarrow$ réaction exothermique incontrôlée $\rightarrow$ surpression $\rightarrow$ ouverture du disque de rupture. Nuage toxique au ras du sol sur les bidonvilles voisins. Bilan officiel : \num{3800} morts immédiats, \num{15000} à \num{20000} morts à long terme, \num{500000} exposés.
    \item \textbf{Toulouse - AZF (France, 21 septembre 2001) :} Explosion d'un stock de \SI{300}{\tonne} de nitrate d'ammonium (engrais) dans l'usine AZF. Onde de choc équivalente à \SI{20}{\tonne} à \SI{40}{\tonne} de TNT. \num{31} morts, \num{2500} blessés, \num{10000} bâtiments endommagés, vitres brisées jusqu'à \SI{3}{\km}.
    \item \textbf{Rupture de barrages :} Brumadinho (Brésil, 2019, barrage de résidus miniers, \num{270} morts), Mariana (Brésil, 2015), ou le barrage de Malpasset (France, 1959, \num{423} morts). Au Cameroun, la surveillance des barrages hydroélectriques (Song Loulou, Lagdo, Memve'ele) est une priorité nationale.
\end{itemize}
\end{exemplebox}

\begin{exoresolu}[Analyse de la chaîne des causes : l'accident d'AZF Toulouse]
\textbf{Énoncé :} À partir du rapport d'enquête officiel, identifier les maillons de la chaîne causale ayant conduit à l'explosion du 21 septembre 2001.

\textbf{Solution détaillée :}
\begin{enumerate}
    \item \textbf{Fait déclencheur (hypothèse retenue) :} Contact entre le nitrate d'ammonium (oxydant) et une substance combustible (probablement du dichloroisocyanurate de sodium, DCCNa, stocké à proximité dans le hangar 221) $\rightarrow$ formation d'un mélange détonant sensible au choc/à la chaleur.
    \item \textbf{Défaillances techniques :} Stockage non conforme (produits incompatibles côte à côte), absence de cloisonnement pare-feu, ventilation insuffisante, détecteurs d'incendie absents ou hors service.
    \item \textbf{Défaillances organisationnelles :} Procédures de réception et de stockage non respectées, turnover du personnel, formation aux risques chimiques lacunaire, sous-traitance mal encadrée, culture de sécurité faible (audits internes n'ont pas détecté l'anomalie).
    \item \textbf{Facteurs contextuels :} Pression de production (stocks élevés), site urbain dense autour de l'usine (plan d'exposition aux risques insuffisant).
    \item \textbf{Conséquence :} Explosion $\rightarrow$ onde de choc $\rightarrow$ projection de débris $\rightarrow$ effondrement de bâtiments $\rightarrow$ victimes et dégâts.
\end{enumerate}
\emph{Enseignement :} La sécurité industrielle repose sur la \cle{défense en profondeur} (barrières successives : conception, contrôle, procédures, formation, organisation, plan d'urgence). La rupture d'une seule barrière ne doit pas suffire à provoquer la catastrophe.
\end{exoresolu}

\begin{attention}[Piège : confusion accident / intention]
Ne confonds pas \emph{catastrophe technologique} (accidentelle, non voulue) et \emph{acte de malveillance / terrorisme / sabotage} (volontaire). Le programme les range parfois ensemble dans les « catastrophes d'origine humaine », mais leurs mécanismes et leur prévention diffèrent radicalement. Une catastrophe technologique relève de la \cle{maîtrise du risque industriel} (normes, inspections, culture de sécurité) ; un acte volontaire relève de la \cle{sécurité civile} et de la \cle{défense}.
\end{attention}

\courssub{Pollutions majeures et catastrophes écologiques}

Ces catastrophes se distinguent par leur caractère souvent \cle{progressif} ou \cle{diffus} à l'origine, puis \cle{brutal} lors de la manifestation aiguë (marée noire, déversement toxique). Elles atteignent les \cle{milieux récepteurs} (air, eau, sol, sédiments) et les \cle{chaînes trophiques}.

\begin{experience}[Marée noire : dynamique et impacts (modèle expérimental simplifié)]
\textbf{Matériel :} Bac transparent (eau de mer), huile colorée (simulant pétrole brut), paille/feuilles (rivage), plume (oiseau), détergent (dispersant), tamis (barrage flottant).
\textbf{Protocole :}
\begin{enumerate}
    \item Verser l'huile à la surface : observation de l'étalement (nappe fine, irisée).
    \item Simuler le vent/agitation (soufflerie douce) : la nappe dérive, s'émulsionne (formation « chocolat mousse » eau-dans-huile), viscosité augmente, récupération difficile.
    \item Déposer sur le rivage (paille) : imprégnation profonde, persistance mois/années.
    \item Tremper la plume : perte d'étanchéité, d'isolation $\rightarrow$ hypothermie, ingestion par lissage $\rightarrow$ intoxication.
    \item Tester dispersant : fragmentation en gouttelettes $\rightarrow$ colonne d'eau contaminée (toxique pour plancton, poissons), mais surface nettoyée.
    \item Tester barrage + récupérateur (tamis + pompe) : efficace par mer calme, inefficace par vague > \SI{1,5}{\meter}.
\end{enumerate}
\textbf{Interprétation :} La lutte est une course contre la montre : \textbf{confinement + récupération mécanique} (priorité) > dispersion chimique (dernier recours) > nettoyage rivage (long, coûteux, destructeur pour la microfaune). La \cle{prévention} (double coque navires, routage, exercices) reste la seule stratégie vraiment efficace.
\end{experience}

\begin{exemplebox}[Le scandale des déchets toxiques de Koko (Nigeria, 1988)]
En 1988, la société italienne Ecomar, affrétée par des négociants suisses, déverse \textbf{\SI{3884}{\tonne}} de déchets toxiques (résidus de production de PCB, solvants chlorés, boues acides) dans le port de \textbf{Koko}, au Nigeria, à quelques centaines de kilomètres des côtes camerounaises. Les fûts, présentés comme « engrais », fuient sous le soleil tropical. Contamination des sols, nappes phréatiques, mangroves. Populations riveraines : éruptions cutanées, troubles respiratoires, avortements spontanés. L'Italie rapatrie les déchets en 1989 sous pression internationale. Cet événement a accéléré l'adoption de la \textbf{Convention de Bâle (1989)} sur le contrôle des mouvements transfrontaliers de déchets dangereux. Au Cameroun, la loi n°96/12 du 5 août 1996 relative à la gestion de l'environnement transpose ces principes.
\end{exemplebox}

\courssub{Accidents de transport collectifs}

Le transport de masse (avion, train, bateau, car) concentre beaucoup de vies dans un espace confiné. La \cle{kinétique de l'accident} (énergie cinétique $E_c = \frac{1}{2}mv^2$) dicte la gravité.

\begin{center}
\begin{tabularx}{\linewidth}{|l|X|X|X|}\hline
\textbf{Mode} & \textbf{Scénarios types} & \textbf{Facteurs humains majeurs} & \textbf{Exemple camerounais} \\ \hline
\textbf{Aérien} & Perte de contrôle, collision, CFIT (impact au sol sans perte de contrôle), incendie & Erreur pilote, fatigue, maintenance défaillante, météo, communication ATC & Crash Kenya Airways Douala (2007), Camair-Co (incidents) \\ \hline
\textbf{Ferroviaire} & Déraillement, collision, passage à niveau & Vitesse excessive, défaut voie/matériel, signalisation, fatigue agent & Déraillement Eseka (2016, \num{79} morts) : vitesse + état voie \\ \hline
\textbf{Maritime/Fluvial} & Naufrage, collision, incendie, chavirage & Surcharge massive, absence gilets, météo, état coque, navigation de nuit & Naufrages fréquents sur le Lac Tchad, Wouri, Sanaga (pirogues surchargées) \\ \hline
\textbf{Routier (cars/poids lourds)} & Sortie de route, collision frontale, incendie & Vitesse, alcool/stupéfiants, fatigue, état véhicule (freins, pneus), infrastructure, indiscipline & Accidents quotidiens sur axes Yaoundé-Douala, Bafoussam-Bamenda \\ \hline
\end{tabularx}
\end{center}

\begin{aretenir}[Gravité des accidents de transport]
La gravité dépend de : (1) \textbf{Énergie libérée} (masse $\times$ vitesse²), (2) \textbf{Kinématique des occupants} (ceinture, structure déformable), (3) \textbf{Environnement} (obstacles latéraux, ravins, feu), (4) \textbf{Délai de secours} (heure d'or). Au Cameroun, la \cle{Commission Nationale de la Sécurité Routière} et le \cle{Ministère des Transports} luttent contre la surcharge, l'importation de véhicules vétustes et l'indiscipline.
\end{aretenir}

\courssub{Incendies et explosions urbains}

L'urbanisation rapide et souvent anarchique des métropoles africaines (Douala, Yaoundé) crée des \cle{pièges à feu} : marchés densément bâtis en matériaux combustibles (bois, tôles, plastiques), ruelles inaccessibles aux engins de secours, réseaux électriques vétustes ou piratés, stockage illicite de carburant/butane.

\begin{itemize}
    \item \textbf{Feux de marchés :} Marché Mokolo (Yaoundé), Marché Central (Douala), Marché Bafoussam. Origines : court-circuit, mégot, feu de cuisine, acte de malveillance. Propagation ultra-rapide (charge de feu énorme, effet cheminée dans les allées). Conséquences : pertes économiques dramatiques pour des milliers de commerçants (souvent informels), pollution atmosphérique aiguë (dioxines, furanes, HAP), traumatismes psychologiques.
    \item \textbf{Dépôts de carburant / stations-service :} Risque d'explosion BLEVE (Boiling Liquid Expanding Vapor Explosion) si un réservoir sous pression est engobé par un feu. Onde de choc + boule de feu. Exemple : dépôt de Bépanda (Douala), incidents récurrents stations-service.
    \item \textbf{Feux de brousse d'origine anthropique :} Au Cameroun, \textbf{> 80\% des feux de végétation} sont d'origine humaine.
\end{itemize}

\begin{savaistu}[Le brûlis agricole : une pratique ancestrale devenue catastrophe écologique]
Au Cameroun, comme en zone soudano-sahélienne (Nord, Adamaoua) et en zone forestière (Sud, Est), le \cle{brûlis} est utilisé pour : (1) nettoyer les champs, (2) minéraliser la matière organique (apport rapide de cendres riches en K, Ca, Mg), (3) chasser le gibier, (4) ouvrir les pistes. \textbf{Mais :} la répétition annuelle, la sécheresse croissante et l'extension des cultures transforment cet outil en fléau.
\begin{itemize}
    \item \textbf{Pertes :} 3 à 4 millions d'hectares brûlés/an au Cameroun (estimations MINEPDED).
    \item \textbf{Impacts sols :} Destruction de la matière organique, destruction de la microfaune/microflore, lessivage/érosion des cendres (perte nutriments), durcissement (croûte de battance), \textbf{appauvrissement durable} $\rightarrow$ baisse des rendements $\rightarrow$ extension des cultures $\rightarrow$ cercle vicieux.
    \item \textbf{Biodiversité :} Mortalité directe (faune lente, nids, graines), fragmentation habitats, favorisation espèces pyrophiles invasives (\emph{Imperata cylindrica}, \emph{Chromolaena odorata}).
    \item \textbf{Climat :} Émissions massives de CO₂, CO, CH₄, N₂O, particules (aérosols carbone noir) $\rightarrow$ réchauffement régional, perturbation régime des pluies.
    \item \textbf{Santé :} Pics de pollution de l'air (PM\num{2,5}) $\rightarrow$ infections respiratoires aiguës, asthme, maladies cardiovasculaires (hôpitaux saturés en saison sèche).
\end{itemize}
\textbf{Alternatives :} Agriculture de conservation (couverture permanente, semis direct sous couverture végétale), agroforesterie, compostage, pâturage tournant contrôlé. La \cle{Stratégie Nationale de Gestion Durable des Feux de Végétation} (2018) vise à encadrer le feu « utile » et réprimer le feu « nuisible ».
\end{savaistu}

\courssub{Catastrophes sanitaires liées à l'Homme}

L'insalubrité urbaine et les défaillances des services de base (eau, assainissement, collecte déchets) transforment des pathologies endémiques en \cle{épidémies explosives}.

\begin{itemize}
    \item \textbf{Choléra (\emph{Vibrio cholerae} O1/O139) :} Maladie des « mains sales » et de l'eau souillée. Facteurs : absence d'eau potable, latrines inexistantes ou à ciel ouvert, inondations (ruissellement fécal), densité de population (quartiers précaires : Nkolbisson, Mimboman, New-Bell, Logbaba). Transmission fécale-orale directe ou via eau/aliments contaminés. Incubation courte (2 h - 5 j) $\rightarrow$ épidémie fulgurante. Traitement : \textbf{réhydratation orale (SRO) ou IV urgente} + antibiotique (doxycycline/azithromycine) pour réduire portage. Prévention : \textbf{EAH (Eau, Assainissement, Hygiène)}, vaccination orale (OCV) en zone à risque.
    \item \textbf{Intoxications alimentaires collectives (IAC) :} \emph{Salmonella}, \emph{Staphylococcus aureus}, \emph{Clostridium perfringens}, \emph{Bacillus cereus}, toxines algales. Contexte : restauration collective (cantines, mariages, funérailles), rupture chaîne du froid, hygiène manipulateurs défaillante. Surveillance : déclaration obligatoire, enquête épidémiologique (courbe épidémique, questionnaire alimentaire, analyses microbiologiques aliments/selles), retrait produits, traitement symptomatologique.
    \item \textbf{Maladies émergentes à potentiel épidémique (Fièvres hémorragiques virales - Ebola, Marburg) :} Liées à l'interface homme/faune sauvage (braconnage, consommation viande de brousse, déforestation). Surveillance \emph{One Health} (santé humaine, animale, environnementale) cruciale au Cameroun (zones forestières Sud/Est).
\end{itemize}

\begin{methode}[Conduite à tenir face à une suspicion d'IAC (cantine scolaire, fête)]
\begin{enumerate}
    \item \textbf{Alerter} immédiatement le chef d'établissement / organisateur $\rightarrow$ Centre de Santé le plus proche $\rightarrow$ District de Santé.
    \item \textbf{Sécuriser} : arrêter la distribution du repas suspect, isoler les restes (échantillons stériles, froid +\SI{4}{\celsius}) pour analyses.
    \item \textbf{Prise en charge médicale} : triage (gravité déshydratation), réhydratation (SRO/IV), antispasmodiques, surveillance signes de gravité (choc, hémolyse \emph{E. coli} O157:H7).
    \item \textbf{Enquête} : définition de cas, recherche facteur commun (plat incriminé), prélèvements selles (cas) + aliments/restes + écouvillons mains/manipulateurs/surfaces.
    \item \textbf{Mesures correctives} : exclusion manipulateurs porteurs, désinfection locaux, révision procédures HACCP (Analyse des dangers et points critiques pour leur maîtrise), formation.
    \item \textbf{Communication} : transparence vis-à-vis familles, médias, éviter rumeurs.
\end{enumerate}
\end{methode}

\courssub{Catastrophes liées aux conflits}

Les conflits armés (guerres, insurrections, terrorisme) génèrent des catastrophes humanitaires complexes : \cle{violences directes} (morts, blessures, mutilations, violences sexuelles), \cle{déplacements forcés} de populations (PDI - Personnes Déplacées Internes, réfugiés), \cle{destruction d'infrastructures} (santé, éducation, eau, électricité, routes, marchés), \cle{effondrement des services} (vaccination, prénatal, traitement VIH/TB/paludisme), \cle{insécurité alimentaire} aiguë (champs abandonnés, pillages, blocus), \cle{traumatismes psychologiques} massifs.

\begin{exemplebox}[Crises humanitaires au Cameroun (contexte 2017-2024)]
\begin{itemize}
    \item \textbf{Régions du Nord-Ouest et du Sud-Ouest (Crise anglophone) :} Depuis 2017, conflit entre forces de défense et groupes armés séparatistes. \textbf{> 700 000 PDI}, \textbf{> 60 000 réfugiés} au Nigeria. Écoles/centres de santé fermés/détruits (stratégie « \emph{lockdown} », attaques ciblées). Recrutement d'enfants. Violence sexuelle. Économie locale effondrée (plantations CDC, PAMOL, petits commerces).
    \item \textbf{Région de l'Extrême-Nord (Insurrection Boko Haram / ISWAP) :} Depuis 2014. \textbf{> 300 000 PDI}, \textbf{> 100 000 réfugiés} nigérians. Attentats kamikazes (filles/femmes enlevées), raids villages, vol bétail, destruction cultures. Insécurité alimentaire sévère (phase 3-4 CIP). Pression sur ressources naturelles (bois, eau) dans zone d'accueil (Mokolo, Kousseri, Maroua). Risque épidémique (choléra, rougeole, paludisme) accru par promiscuité camps.
\end{itemize}
Ces situations relèvent de la \cle{protection civile} et du \cle{Droit International Humanitaire} (Conventions de Genève) : distinction combattants/civils, proportionnalité, précaution, accès humanitaire impartial.
\end{exemplebox}

\courssub{Facteurs humains aggravants des catastrophes}

Une catastrophe « naturelle » (inondation, glissement de terrain, sécheresse) devient souvent une catastrophe \emph{majeure} à cause de l'action humaine. C'est le concept de \cle{construction sociale du risque} : \textbf{Risque = Aléa $\times$ Vulnérabilité $\times$ Enjeu}. L'Homme agit sur la Vulnérabilité et l'Enjeu.

\begin{popfigure}
\begin{tikzpicture}[
    node distance=1.2cm and 1.5cm,
    box/.style={rectangle, draw=popdark, thick, fill=popcream, rounded corners, minimum width=3.2cm, minimum height=1.1cm, align=center, font=\small},
    arrow/.style={-Stealth, thick, popdark},
    plus/.style={circle, draw=poporange, fill=poporange!20, minimum size=0.6cm, font=\tiny\bfseries, inner sep=1pt}
]
% Nœuds principaux
\node[box] (defor) {Déforestation /\newline Déboisement anarchique};
\node[box, right=of defor] (eros) {Érosion hydrique /\newline éolienne accélérée};
\node[box, right=of eros] (sol) {Perte fertilité /\newline Appauvrissement sols};
\node[box, below=of eros] (glis) {Glissements de terrain /\newline Coulées de boue};
\node[box, below=of sol] (inond) {Inondations rapides /\newline Ruissellement accru};
\node[box, left=of glis] (urb) {Urbanisation anarchique /\newline Occupation zones à risque};
\node[box, above=of urb] (pauv) {Pauvreté /\newline Absence d'alternatives};

% Flèches principales
\draw[arrow] (defor) -- (eros);
\draw[arrow] (eros) -- (sol);
\draw[arrow] (eros) -- (glis);
\draw[arrow] (eros) -- (inond);
\draw[arrow] (urb) -- (glis);
\draw[arrow] (urb) -- (inond);
\draw[arrow] (pauv) -- (defor);
\draw[arrow] (pauv) -- (urb);

% Boucles de rétroaction (pointillés)
\draw[arrow, dashed, poppurple, bend left=30] (sol) to node[plus, pos=0.5] {+} (pauv);
\draw[arrow, dashed, poppurple, bend right=30] (glis) to node[plus, pos=0.5] {+} (pauv);
\draw[arrow, dashed, poppurple, bend right=30] (inond) to node[plus, pos=0.5] {+} (pauv);

% Légendes des boucles
\node[align=center, font=\tiny, text=poppurple] at (2.5, -3.2) {Boucles de rétroaction positive :\newline catastrophe $\rightarrow$ pauvreté $\rightarrow$ pression sur ressources};
\end{tikzpicture}
\legende{Schéma de la chaîne des facteurs humains aggravants : la déforestation et l'urbanisation anarchique, souvent motivées par la pauvreté, déclenchent l'érosion qui favorise glissements de terrain et inondations, qui à leur tour aggravent la pauvreté (boucles de rétroaction positive).}
\end{popfigure}

\begin{aretenir}[Principaux facteurs aggravants d'origine humaine]
\begin{enumerate}
    \item \textbf{Urbanisation anarchique :} Occupation des zones inondables (lits majeurs, bas-fonds), pentes instables, couloirs d'écoulement. Absence de réseaux d'assainissement (caniveaux bouchés, inexistants). Imperméabilisation des sols (béton) $\rightarrow$ coefficient de ruissellement $\uparrow$ $\rightarrow$ crue éclair.
    \item \textbf{Déforestation / Dégradation des sols :} Perte du pouvoir tampon des forêts (interception pluie, infiltration, cohésion sols par racines). Brûlis, coupe bois de feu/charbon, extension agricole non durable.
    \item \textbf{Non-respect des normes :} Construction sans étude géotechnique, sans permis, matériaux de récupération. Industries sans étude d'impact, sans plan d'urgence interne (PUI), rejets non traités.
    \item \textbf{Pauvreté et inégalités :} Populations forcées d'habiter zones à risque (coût foncier faible). Absence d'assurance, d'épargne, de filets de protection sociale $\rightarrow$ incapacité à se préparer, à évacuer, à se reconstruire. Dépendance aux ressources naturelles non durables.
    \item \textbf{Gouvernance faible :} Application laxiste des lois (urbanisme, environnement, travail, mines), corruption, manque de moyens des services de contrôle (MINEPDED, MINHDU, MINT, MINSANTE), planification urbaine dépassée, coordination interministérielle défaillante en crise.
\end{enumerate}
\end{aretenir}

\courssub{Comparaison systématique : Catastrophes naturelles vs Catastrophes humaines}

\begin{popfigure}
\begin{tikzpicture}[
    cell/.style={rectangle, draw=popline, thin, minimum height=0.9cm, align=center, font=\small, inner xsep=2pt},
    header/.style={cell, fill=popblue, font=\small\bfseries, text=white},
    subheader/.style={cell, fill=popblueL, font=\small\bfseries, text=popdark},
    row1/.style={fill=popblueL!30},
    row2/.style={fill=white}
]
% Largeurs colonnes
\def\wA{3.2cm}
\def\wB{4.5cm}
\def\wC{4.5cm}

% En-tête
\node[header, minimum width=\wA] (h1) at (0,0) {Critère};
\node[header, minimum width=\wB, anchor=west] (h2) at (h1.east) {Catastrophes naturelles};
\node[header, minimum width=\wC, anchor=west] (h3) at (h2.east) {Catastrophes humaines};

% Ligne 1 : Déclencheur
\node[cell, row1, minimum width=\wA, anchor=north] (c11) at (h1.south) {Déclencheur principal};
\node[cell, row1, minimum width=\wB, anchor=north west] (c12) at (c11.east) {Phénomène géologique, météorologique, biologique \textbf{indépendant de la volonté humaine} (séisme, volcan, cyclone, tsunami, épidémie zoonotique émergente)};
\node[cell, row1, minimum width=\wC, anchor=north west] (c13) at (c12.east) {Action, négligence, erreur, volonté humaine (technique, organisation, conflit, pollution, urbanisation)};

% Ligne 2 : Prévisibilité
\node[cell, row2, minimum width=\wA, anchor=north] (c21) at (c11.south) {Prévisibilité / Prédiction};
\node[cell, row2, minimum width=\wB, anchor=north west] (c22) at (c21.east) {Variable : \textbf{cyclones} (jours), \textbf{inondations} (heures/jours), \textbf{séismes} (imprévisibles à court terme, aléa probabiliste), \textbf{volcans} (semaines/jours si monitorés)};
\node[cell, row2, minimum width=\wC, anchor=north west] (c23) at (c22.east) {Souvent \textbf{prévisibles à long terme} (usines vieillissantes, zones à risque connues, conflits latents). \textbf{Prédiction court terme difficile} (erreur humaine soudaine, sabotage), mais \textbf{scénarios identifiables} (études de dangers)};

% Ligne 3 : Prévention
\node[cell, row1, minimum width=\wA, anchor=north] (c31) at (c21.south) {Prévention / Réduction vulnérabilité};
\node[cell, row1, minimum width=\wB, anchor=north west] (c32) at (c31.east) {\textbf{Impossible d'empêcher l'aléa}. Réduction enjeu/vulnérabilité : aménagement territoire (PPRI), construction parasismique/anticyclonique, systèmes alerte précoce (SAP), éducation, assurances, repli stratégique};
\node[cell, row1, minimum width=\wC, anchor=north west] (c33) at (c32.east) {\textbf{Possible d'empêcher ou réduire l'aléa} : maintenance, normes conception, formation, culture sécurité, régulation, traités internationaux, résolution conflits, gestion durable ressources. Réduction enjeu : même leviers que naturels};

% Ligne 4 : Responsabilité
\node[cell, row2, minimum width=\wA, anchor=north] (c41) at (c31.south) {Responsabilité / Cadre juridique};
\node[cell, row2, minimum width=\wB, anchor=north west] (c42) at (c41.east) {Responsabilité de l'État : protection populations (obligation de moyens), aménagement, secours. Pas de « coupable » de l'aléa. Droit humanitaire (catastrophes) ; assurance solidaire (CatNat)};
\node[cell, row2, minimum width=\wC, anchor=north west] (c43) at (c42.east) {Responsabilité \textbf{pénale et civile} identifiable (exploitant, employeur, transporteur, commanditaire conflit). Pollueur-payeur. Droit pénal (mise en danger, homicide involontaire), droit de l'environnement, DIH (crimes de guerre). Réparation intégrale};

% Ligne 5 : Exemples camerounais
\node[cell, row1, minimum width=\wA, anchor=north] (c51) at (c41.south) {Exemples au Cameroun};
\node[cell, row1, minimum width=\wB, anchor=north west] (c52) at (c51.east) {Éruption Mont Cameroun (1999, 2000), Éruption limnique Lac Nyos (1986), Inondations Logone/Chari (annuel), Séismes région Kribi/Campo, Glissements terrain (Mbankolo, Bafoussam)};
\node[cell, row1, minimum width=\wC, anchor=north west] (c53) at (c52.east) {Incendies marchés (Mokolo, Central), Naufrages (Lac Tchad, Wouri), Accidents cars (Eseka train 2016, axes routiers), Pollution industrielle (usines Douala/Bonabéri), Crises N-O/S-O/Extrême-Nord, Feux de brousse anthropiques, Déchets Koko (proximité)};

\end{tikzpicture}
\legende{Tableau comparatif : catastrophes naturelles vs catastrophes d'origine humaine. La distinction n'est pas toujours étanche (ex: inondation aggravée par urbanisation), mais elle guide la stratégie de gestion (subir/adapter vs éliminer/contrôler la source).}
\end{popfigure}

\begin{exoresolu}[Classification d'événements récents au Cameroun]
\textbf{Énoncé :} Classer les événements suivants dans le tableau : (1) Éruption du Mont Cameroun 2000, (2) Incendie marché Mokolo 2023, (3) Naufrage pirogue sur le Wouri 2022 (surcharge), (4) Inondation Yaoundé quartier Mimboman juin 2024 (ruissellement + caniveaux bouchés), (5) Attaque Boko Haram village Extrême-Nord 2023, (6) Glissement terrain Mbankolo 2019 (après fortes pluies + coupe arbres), (7) Fuite chlore usine SOCAVER Douala (hypothétique), (8) Épidémie choléra Nord 2022 (eau insalubre).

\textbf{Solution détaillée :}
\begin{center}
\begin{tabular}{|l|l|l|}\hline
\textbf{Origine Naturelle (Aléa dominant)} & \textbf{Origine Humaine (Cause dominante)} & \textbf{Mixte / Naturelle aggravée par l'Homme} \\ \hline
(1) Éruption Mont Cameroun 2000 & (2) Incendie marché Mokolo (électrique/malveillance) & (4) Inondation Mimboman (pluie + urbanisation/caniveaux) \\
(5) Attaque Boko Haram (conflit = volonté humaine) & (3) Naufrage Wouri (surcharge = négligence) & (6) Glissement Mbankolo (pluie + déforestation) \\
& (7) Fuite chlore usine (défaillance technique) & (8) Choléra Nord (bactérie + insalubrité humaine) \\
\hline
\end{tabular}
\end{center}
\emph{Analyse :} La plupart des « catastrophes naturelles » au Cameroun sont en réalité \textbf{mixtes} : l'aléa météorologique (pluie extrême) rencontre une vulnérabilité construite par l'Homme (urbanisation, déforestation, absence d'assainissement). La gestion du risque impose d'agir sur les facteurs humains.
\end{exoresolu}

\begin{attention}[Erreur fréquente au Bac : « Les catastrophes naturelles sont inévitables »]
\textbf{Faux.} L'aléa (la pluie, le séisme) est inévitable. La \textbf{catastrophe} (le bilan en vies et en biens) ne l'est pas. Elle résulte de la rencontre entre l'aléa et une société vulnérable. La prévention (aménagement, construction, éducation, alerte) réduit drastiquement le bilan. Exemple : séisme magnitude 7,0 Haïti 2010 $\rightarrow$ 200 000+ morts (construction informelle) ; séisme magnitude 9,0 Japon 2011 $\rightarrow$ ~20 000 morts (majoritairement tsunami, construction parasismique a tenu). Au Cameroun, l'application du \textbf{Code de l'Urbanisme} et du \textbf{Plan de Prévention des Risques (PPR)} est la clé.
\end{attention}

\courssub{Gestes qui sauvent : Conduite à tenir générale face à une catastrophe}

\begin{methode}[Réflexes vitaux et conduite à tenir (Savoir-faire Bac)]
\begin{enumerate}
    \item \textbf{Avant (Préparation) :} Connaître les risques locaux (PPR, plan communal de sauvegarde), constituer un \textbf{kit d'urgence} (eau \SI{2}{\liter}/pers/jour $\times$ 3 j, aliments non périssables, lampe torche, radio piles, trousse secours, documents d'identité photocopiés, sifflet), repérer les issues de secours et points de rassemblement (famille, école, travail).
    \item \textbf{Pendant (Protection immédiate) :}
    \begin{itemize}
        \item \textbf{Séisme :} « Se baisser, s'abriter, s'agripper » (sous table solide, loin des vitres) ; si dehors : s'éloigner bâtiments/poteaux.
        \item \textbf{Inondation :} Monter en hauteur (étage, toit), ne JAMAIS traverser une route inondée (\SI{30}{\centi\meter} d'eau emportent une voiture), couper électricité/gaz.
        \item \textbf{Incendie :} Rampant sous la fumée (air moins toxique), porte fermée (freine propagation), alerter \num{118} (pompiers) ou \num{112} (urgence européenne), ne pas prendre ascenseur.
        \item \textbf{Accident chimique/industriel :} S'abriter (confinement) : rentrer, fermer portes/fenêtres/ventilation, sceller interstices (ruban adhésif, linge mouillé), écouter radio consignes.
        \item \textbf{Attentat/Tir actif :} \textbf{Fuir} (sortie la plus proche), \textbf{Se cacher} (pièce verrouillée, éteindre lumière/téléphone silencieux), \textbf{Alerter} (numéro d'urgence) quand en sécurité.
    \end{itemize}
    \item \textbf{Après (Secours et résilience) :} Ne pas bouger les blessés graves sauf danger immédiat (incendie, effondrement imminent). Appliquer \textbf{PAS} (Protéger, Alerter, Secourir) : \textbf{P}roteger la zone (balisage, triangle), \textbf{A}lerter les secours (\num{112}, \num{118}, \num{119}) en donnant : \textbf{Où} (précis), \textbf{Quoi} (nature), \textbf{Combien} (victimes), \textbf{État} (conscient, saigne, respire), \textbf{S}ecourir (PLS si inconscient respirant, massage cardiaque + DAE si arrêt cardiaque, compression point de compression si hémorragie externe massive). Ne pas téléphoner inutilement (saturer réseaux). Suivre consignes autorités (évacuation, confinement, distribution eau/aliments). Soutien psychologique (cellules d'urgence médico-psychologique - CUMP).
\end{enumerate}
\end{methode}

\courssub{Synthèse et transition vers la gestion}

Les catastrophes d'origine humaine sont \textbf{évitables} par essence, car leur cause réside dans nos choix techniques, organisationnels, politiques et éthiques. Leur étude révèle une constante : la \cle{rupture de barrières de sécurité} (physiques, procédurales, humaines) que la \cle{culture de la prévention} doit maintenir debout. La section suivante détaillera la \textbf{gestion intégrée des catastrophes} : comment passer de la réaction à l'anticipation, de la prévision à la résilience, à travers les piliers \textbf{Prévision -- Prédiction -- Prévention -- Préparation -- Intervention -- Reconstruction}.

\coursec{La gestion des catastrophes : prévision, prédiction, prévention}

Après avoir identifié et classé les catastrophes selon leur origine (naturelle ou humaine) dans la section précédente, nous devons à présent comprendre comment les sociétés s'organisent pour y faire face. Une catastrophe n'est pas une fatalité absolue : si l'aléa (le phénomène dangereux) est souvent inévitable, le risque (la rencontre entre l'aléa et des enjeux vulnérables) peut être maîtrisé. C'est tout l'objet de la \cle{gestion des catastrophes}, qui repose sur trois piliers fondamentaux : la prévision, la prédiction et la prévention, articulés autour d'un cycle continu d'actions avant, pendant et après l'événement.

\courssub{Les trois piliers de la réduction du risque : prévision, prédiction, prévention}

Ces trois termes sont souvent confondus dans le langage courant, mais ils désignent des démarches scientifiques et opérationnelles distinctes, complémentaires et opérant à des échelles de temps différentes. Une confusion entre eux conduit à des erreurs d'appréciation du danger et à des décisions inadaptées.

\begin{definition}[Prévision]
La \textbf{prévision} est l'estimation, à \textbf{long terme} (années, décennies), de la \textbf{probabilité} de survenue d'un phénomène dangereux et de sa \textbf{localisation} probable, fondée sur l'analyse statistique des données historiques, géologiques et météorologiques (cartes d'aléas, courbes de récurrence). Elle ne donne ni la date, ni l'heure, ni l'intensité exacte de l'événement futur.
\end{definition}

\begin{definition}[Prédiction]
La \textbf{prédiction} est l'annonce, à \textbf{court terme} (heures, jours, semaines), de la survenue \textbf{imminente} d'un événement, avec une estimation de sa \textbf{date}, de son \textbf{lieu} précis et de son \textbf{intensité}. Elle repose sur la surveillance en temps réel de précurseurs physiques (sismicité, déformation du sol, pression atmosphérique, niveau des cours d'eau).
\end{definition}

\begin{definition}[Prévention]
La \textbf{prévention} est l'ensemble des \textbf{mesures structurelles et non structurelles} mises en œuvre \textbf{en amont} (avant la catastrophe) pour \textbf{réduire la vulnérabilité} des enjeux (populations, bâtiments, infrastructures, environnement) et \textbf{limiter les dégâts} potentiels. Elle transforme le territoire pour le rendre plus résilient.
\end{definition}

\begin{attention}[Erreur fréquente : Confondre prévision et prédiction]
Ne confondez jamais ces deux notions.
\begin{itemize}
    \item \textbf{Prévision} = \og Où et avec quelle fréquence cela arrive-t-il généralement ? \fg (Statistique, long terme). Exemple : \og La zone de Foumban a une probabilité élevée de secousses sismiques sur 50 ans. \fg
    \item \textbf{Prédiction} = \og Quand, où exactement et avec quelle force cela va-t-il arriver maintenant ? \fg (Physique temps réel, court terme). Exemple : \og Une éruption du Mont Cameroun est attendue dans les 48 prochaines heures sur le flanc Sud-Est. \fg
\end{itemize}
\textbf{Piège du Bac :} On ne peut pas \og prévoir \fg un séisme pour demain (prédiction impossible à ce jour), mais on peut \og prévoir \fg que la région de Kribi est soumise à un aléa sismique modéré (prévision statistique). La prévision guide l'urbanisme ; la prédiction déclenche l'alerte et l'évacuation.
\end{attention}

\begin{experience}[Construction d'une carte d'aléa par prévision statistique (démarche d'investigation)]
\textbf{Observation :} Les archives du service météorologique de Douala montrent que le quartier de Bepanda est inondé en moyenne 3 fois par an depuis 20 ans, avec des hauteurs d'eau variables.
\textbf{Problématique :} Comment estimer la probabilité qu'une crue exceptionnelle (hauteur $> 1,5$ m) se produise l'an prochain ?
\textbf{Hypothèse :} La fréquence historique des crues suit une loi statistique (loi de Gumbel ou de Weibull) permettant d'extrapoler la période de retour $T$ d'une crue de hauteur $H$.
\textbf{Exploitation de document (Tableau de relevés annuels max de hauteur d'eau à Bepanda, 2000-2020) :}
\begin{center}
\begin{tabularx}{\linewidth}{|c|X|}\hline
\textbf{Année} & \textbf{Hauteur max (m)} \\ \hline
2000 & 0,8 \\ 2001 & 1,2 \\ 2002 & 0,6 \\ 2003 & 1,5 \\ 2004 & 0,9 \\ 2005 & 1,1 \\ 2006 & 0,7 \\ 2007 & 1,8 \\ 2008 & 1,0 \\ 2009 & 1,3 \\ \hline
2010 & 0,8 \\ 2011 & 1,4 \\ 2012 & 0,9 \\ 2013 & 1,6 \\ 2014 & 1,1 \\ 2015 & 0,7 \\ 2016 & 1,2 \\ 2017 & 1,9 \\ 2018 & 1,0 \\ 2019 & 1,3 \\ 2020 & 1,5 \\ \hline
\end{tabularx}
\end{center}
\textbf{Traitement :} On classe les hauteurs par ordre décroissant. La crue de 1,9 m (2017) est la plus forte (rang $r=1$) sur $N=21$ ans. Sa période de retour estimée : $T = \frac{N+1}{r} = \frac{22}{1} = 22$ ans. La crue de 1,8 m (2007) a $r=2$, $T=11$ ans. Une crue de 1,5 m a été observée 3 fois (rangs 3, 4, 5 environ), $T \approx 5$ à $7$ ans.
\textbf{Interprétation :} La \textbf{prévision} indique qu'une crue $\ge 1,5$ m a une période de retour d'environ 5-7 ans (probabilité annuelle $\approx 15\%$ à $20\%$). Cela ne dit pas \emph{quand} elle aura lieu l'an prochain, mais que le risque est réel et récurrent.
\textbf{Conclusion :} La prévision statistique permet de classer les zones en aléa fort, moyen, faible pour l'urbanisme (carte d'aléa), base de la prévention.
\end{experience}

\begin{popfigure}
\begin{tikzpicture}[scale=0.9, every node/.style={font=\small}]
% Axes
\draw[->, thick] (0,0) -- (11,0) node[right] {Temps};
\draw[->, thick] (0,0) -- (0,6) node[above] {Niveau d'action / Information};

% Prévision (Long terme - Fond)
\fill[popblue!15] (0.5,0.5) rectangle (10.5,1.5);
\draw[popblue, thick] (0.5,0.5) rectangle (10.5,1.5);
\node[popblue, align=center] at (5.5,1.0) {\textbf{PRÉVISION} (Long terme : années -- décennies)\\\textit{Probabilité, localisation, cartes d'aléas, urbanisme}};

% Prévention (Long terme - Action)
\fill[popgreen!15] (0.5,2.0) rectangle (10.5,3.0);
\draw[popgreen, thick] (0.5,2.0) rectangle (10.5,3.0);
\node[popgreen, align=center] at (5.5,2.5) {\textbf{PRÉVENTION} (Continu, en amont)\\\textit{Normes parasismiques, reboisement, dégazage Nyos, zonage}};

% Prédiction (Court terme - Alerte)
\fill[poporange!15] (7.0,3.5) rectangle (10.5,4.5);
\draw[poporange, thick] (7.0,3.5) rectangle (10.5,4.5);
\node[poporange, align=center] at (8.75,4.0) {\textbf{PRÉDICTION} (Court terme : heures -- jours)\\\textit{Alerte météo, surveillance volcanique, crue imminente}};

% Flèches de liaison
\draw[->, thick, popmuted] (5.5,1.5) -- (5.5,2.0) node[midway, right] {guide};
\draw[->, thick, popmuted] (5.5,3.0) -- (5.5,3.5) node[midway, right] {prépare la réponse à};
\draw[->, thick, popmuted] (8.75,3.5) .. controls (9.5,3.0) and (9.5,1.5) .. (10,1.0) node[right, align=left] {affine les modèles\\de prévision};

% Échelle de temps
\draw[decorate, decoration={brace, amplitude=5pt}, popdark] (0, -0.3) -- (11, -0.3) node[midway, below=5pt] {Échelle de temps : Long terme \hfill Court terme};
\end{tikzpicture}
\legende{Articulation des trois piliers : la prévision (statistique, long terme) guide la prévention (aménagement, continu) ; la prédiction (temps réel, court terme) déclenche la gestion de crise. La prédiction nourrit en retour les modèles de prévision.}
\end{popfigure}

\begin{exemplebox}[Exemples concrets au Cameroun]
\begin{itemize}
    \item \textbf{Prévision (Séisme) :} La carte d'aléa sismique du Cameroun (établie par l'IRGM et le MINRESI) classe la zone du Mont Cameroun, de la plaine de la Mémé et de la région de Kribi en aléa \og moyen \fg à \og fort \fg. Cela ne prédit pas le séisme de demain, mais impose des normes parasismiques (Eurocode 8 adapté) pour les nouveaux bâtiments.
    \item \textbf{Prédiction (Éruption volcanique) :} Le réseau de surveillance du Mont Cameroun (sismomètres, GPS, inclinomètres, gaz) géré par l'Observatoire Volcanologique du Cameroun (OVC) a permis de \textbf{prédire} l'éruption de 2000 (début le 28 mai) quelques jours à l'avance grâce à une crise sismique précurseure et une déformation du flanc, permettant l'évacuation de Bakingili.
    \item \textbf{Prédiction (Inondation) :} La vigilance météorologique de la DMN (Direction de la Météorologie Nationale) émet des bulletins d'alerte (vert, jaune, orange, rouge) 24h à 48h avant de fortes pluies annoncées par les modèles numériques, permettant aux populations des bas-fonds de Yaoundé (Mvog-Ada, Ekounou) de se préparer.
    \item \textbf{Prévention (Lac Nyos) :} Le dégazage artificiel par tuyaux (installés en 2001, 2011, 2019) est une mesure de \textbf{prévention structurelle} majeure : il réduit \emph{en continu} la quantité de $\ce{CO2}$ dissous dans les eaux profondes, abaissant la probabilité d'une nouvelle éruption limnique catastrophique (comme celle de 1986 : 1 746 morts).
    \item \textbf{Prévention (Urbain) :} Le Plan d'Urbanisme Directeur (PUD) de Yaoundé interdit théoriquement la construction en zone inondable (lit majeur de la Mfoundi, marécages) et impose des largeurs de rues pour l'accès des secours.
\end{itemize}
\end{exemplebox}

\courssub{Le cycle de gestion des catastrophes : trois phases temporelles}

La gestion ne se limite pas à l'urgence. Elle suit un \textbf{cycle continu} où la fin d'une catastrophe prépare la résilience face à la suivante.

\begin{popfigure}
\begin{tikzpicture}[scale=0.95, every node/.style={font=\small}]
% Styles
\tikzset{
    phase/.style={draw, thick, rounded corners, minimum width=3.2cm, minimum height=2.8cm, align=center, font=\small},
    arrow/.style={->, >=Stealth, thick, popdark},
    title/.style={font=\bfseries\large, text=popdark},
    action/.style={font=\footnotesize, text=popmuted},
    highlight/.style={font=\bfseries\footnotesize, text=popblue}
}

% Nœuds des 3 phases
\node[phase, fill=popblue!10, draw=popblue, align=center] (avant) at (0,0){
    \textbf{AVANT}\\(Prévention -- Planification -- Éducation)\\[4pt]
    \textit{Mesures structurelles :}\\
    \bullet Normes construction (parasismique, cyclonique)\\
    \bullet Reboisement pentes / bandes riveraines\\
    \bullet Dégazage Lac Nyos, barrages de retenue\\
    \bullet Zonage (interdiction zone inondable)\\[4pt]
    \textit{Mesures non structurelles :}\\
    \bullet Cartes de risques, Plans ORSEC / PCS\\
    \bullet Exercices simulation (écoles, entreprises)\\
    \bullet Sensibilisation (radio, affiches, chefs traditionnels)\\
    \bullet Systèmes alerte précoce (sirènes, SMS)
};

\node[phase, fill=poporange!10, draw=poporange, align=center] (pendant) at (7,0){
    \textbf{PENDANT}\\(Alerte -- Évacuation -- Secours)\\[4pt]
    \textit{Déclenchement :}\\
    \bullet Alerte précoce (météo, volcan, crue)\\
    \bullet Sirènes, SMS, radio communautaire, porte-à-porte\\[4pt]
    \textit{Évacuation :}\\
    \bullet Mise à l'abri (sites de repli identifiés)\\
    \bullet Aide aux personnes vulnérables (PMR, âgés)\\[4pt]
    \textit{Secours immédiats :}\\
    \bullet Premiers secours (PSC1), triage victimes\\
    \bullet Protection civile, Pompiers, Forces de défense\\
    \bullet Distribution eau, vivres, abris d'urgence
};

\node[phase, fill=popgreen!10, draw=popgreen, align=center] (apres) at (3.5,-4){
    \textbf{APRÈS}\\(Secours d'urgence -- Reconstruction -- Bilan)\\[4pt]
    \textit{Urgence (J+1 à J+30) :}\\
    \bullet Recherche/sauvetage, soins médicaux\\
    \bullet Hébergement temporaire, eau/assainissement\\
    \bullet Sécurité civile, lutte contre pillages\\[4pt]
    \textit{Reconstruction (Mois/Années) :}\\
    \bullet Réhabilitation infrastructures (routes, écoles)\\
    \bullet \textbf{Reconstruction résiliente} (normes améliorées)\\
    \bullet Soutien psychologique, relance économique\\[4pt]
    \textit{Retour d'expérience (REX) :}\\
    \bullet Bilan gestion de crise\\
    \bullet Mise à jour cartes risques, plans ORSEC\\
    \bullet Alimentation base de données nationale
};

% Flèches cycliques
\draw[arrow, popblue, bend left=15] (avant.east) to node[above, highlight] {Préparation continue} (pendant.west);
\draw[arrow, poporange, bend left=15] (pendant.south) to node[right, highlight] {Fin alerte / Bilan} (apres.east);
\draw[arrow, popgreen, bend left=15] (apres.west) to node[left, highlight] {Leçons apprises / Résilience} (avant.south west);

% Titre central
\node[title] at (3.5, 1.5) {\textbf{CYCLE DE GESTION DES CATASTROPHES}};
\node[action] at (3.5, 1.0) {Approche intégrée : Réduction du risque (Avant) $\rightarrow$ Gestion de crise (Pendant) $\rightarrow$ Relèvement (Après)};
\end{tikzpicture}
\legende{Le cycle de gestion des catastrophes : une boucle continue où la phase \og Après \fg (Retour d'Expérience - REX) alimente la phase \og Avant \fg (mise à jour des préventions, plans et formations).}
\end{popfigure}

\begin{aretenir}[Les trois phases de la gestion]
\begin{enumerate}
    \item \textbf{Avant (Réduction du risque / Prévention) :} C'est la phase la plus efficace et la moins coûteuse en vies humaines. Elle repose sur la \textbf{connaissance du risque} (cartes), la \textbf{réglementation} (urbanisme, normes), l'\textbf{aménagement} (ouvrages de protection, reboisement) et la \textbf{préparation} (plans d'urgence, exercices, éducation).
    \item \textbf{Pendant (Gestion de crise / Réponse) :} Déclenchée par l'alerte (prédiction ou survenue brutale). Objectifs : \textbf{sauver des vies} (évacuation, secours médicaux), \textbf{limiter les dégâts secondaires} (incendies, épidémies), \textbf{assurer les besoins vitaux} (eau, abri, nourriture). Coordination sous l'autorité du Préfet / Maire (PCS - Plan Communal de Sauvegarde) ou du Ministre de l'Intérieur (ORSEC).
    \item \textbf{Après (Relèvement / Reconstruction) :} Phase longue. \textbf{Urgence humanitaire} puis \textbf{reconstruction durable} (\og \emph{Build Back Better} \fg : reconstruire mieux qu'avant, en intégrant les leçons du sinistre). Le \textbf{Retour d'Expérience (REX)} est obligatoire pour mettre à jour les cartes de risques et les plans d'urgence.
\end{enumerate}
\end{aretenir}

\courssub{Acteurs et outils de la gestion au Cameroun}

La gestion des catastrophes au Cameroun est coordonnée par l'État mais implique une multiplicité d'acteurs selon le principe de \textbf{subsidiarité} (l'échelon local gère, l'échelon supérieur appuie si dépassé).

\begin{center}
\begin{tabularx}{\linewidth}{|l|X|X|}\hline
\rowcolor{popblueL}
\textbf{Acteur} & \textbf{Rôle principal} & \textbf{Exemple d'intervention} \\ \hline
\textbf{Direction de la Protection Civile (DPC) / MINAT} & Coordination nationale, plans ORSEC, formation secours, gestion abris & Activation Plan ORSEC \og Éruption Mont Cameroun \fg ; formation moniteurs PSC1 \\ \hline
\textbf{Communes (Maires) / PCS} & Premier niveau de réponse, Plan Communal de Sauvegarde (PCS), alerte population & Évacuation quartiers inondés Yaoundé IV ; ouverture salles de classe comme abris \\ \hline
\textbf{Croix-Rouge Camerounaise (CRC)} & Secours d'urgence, premiers soins, restauration liens familiaux, appui logistique & Distribution kits hygiène après inondation Logone-et-Chari ; postes de secours \\ \hline
\textbf{Forces de Défense et de Sécurité (Police, Gendarmerie, BIR, Génie Militaire)} & Maintien ordre, recherche/sauvetage (GRIMP, plongeurs), logistique lourde (ponts, routes) & Déblaiement routes après glissement terrain Dschang ; évacuation héliportée \\ \hline
\textbf{Services Techniques (MINTP, MINEPDED, MINSANTE, DMN, IRGM/OVC)} & Expertise technique : surveillance (météo, volcans), santé publique, environnement, infrastructures & Bulletin vigilance DMN ; surveillance sismique IRGM ; chloration eau MINSANTE \\ \hline
\textbf{Populations / ONG / Chefs traditionnels / Médias} & Acteurs de terrain : respect consignes, entraide, relais information, connaissance locale & Vigilance comités de quartier ; radio communautaire Buea alerte éruption ; ONG reboisement \\ \hline
\end{tabularx}
\end{center}

\begin{propriete}[Outils opérationnels de la gestion]
\begin{itemize}
    \item \textbf{Cartes de risques (Aléa $\times$ Enjeux $\times$ Vulnérabilité) :} Documents cartographiques (souvent SIG) identifiant les zones à risque (inondation, glissement, éruption, industriel). Base du zonage réglementaire (PUD, POS). \emph{Exemple : Carte de risque inondation Yaoundé (voir figure ci-dessous).}
    \item \textbf{Plan ORSEC (Organisation de la Réponse de SÉcurité Civile) :} Dispositif national de mobilisation des moyens de l'État (niveau départemental, régional, national). Déclenché par le Préfet ou le Ministre de l'Intérieur.
    \item \textbf{Plan Communal de Sauvegarde (PCS) :} Document opérationnel communal : recensement des risques, moyens locaux, sites d'hébergement, annuaire acteurs, fiches réflexes. Obligatoire pour toute commune.
    \item \textbf{Systèmes d'Alerte Précoce (SAP) :} Chaîne \og Capteur $\rightarrow$ Centre d'analyse $\rightarrow$ Décision $\rightarrow$ Diffusion (Sirènes, SMS \emph{Cell Broadcast}, Radio, Réseaux sociaux, Porte-voix) $\rightarrow$ Population \fg. Objectif : \textbf{gagner du temps} pour l'évacuation.
    \item \textbf{Exercices de simulation (Exercices cadres, exercices terrain) :} Testent les plans, forment les acteurs, sensibilisent le public. \emph{Exercice annuel \og Mont Cameroun \fg à Buea/Limbe ; exercices tsunami zone côtière.}
\end{itemize}
\end{propriete}

\begin{popfigure}
\begin{tikzpicture}[scale=0.85, every node/.style={font=\small}]
% Légende couleurs
\node[align=left, font=\bfseries\small] at (0.2, 6.5) {Légende :};
\fill[popblue!80] (0.2, 6.1) rectangle ++(0.6,0.3) node[right=2pt, midway] {Aléa Fort (Lit mineur / Marécage)};
\fill[poporange!80] (0.2, 5.7) rectangle ++(0.6,0.3) node[right=2pt, midway] {Aléa Moyen (Lit majeur 50-100 ans)};
\fill[popgreen!80] (0.2, 5.3) rectangle ++(0.6,0.3) node[right=2pt, midway] {Aléa Faible (Hors crue centennale)};
\fill[poppink!80] (0.2, 4.9) rectangle ++(0.6,0.3) node[right=2pt, midway] {Enjeux forts (Habitat dense, École, Hôpital)};
\draw[popred, ultra thick, ->] (0.2, 4.5) -- ++(0.6,0) node[right=2pt, midway] {Chemin évacuation vers site repli};
\draw[popblue, thick, dashed] (0.2, 4.1) -- ++(0.6,0) node[right=2pt, midway] {Digue / Ouvrage protection};

% Carte simplifiée Quartier Mvog-Ada (Yaoundé)
% Rivière Mfoundi
\draw[popblue, line width=3pt, decorate, decoration={snake, amplitude=1pt, segment length=10pt}] (1,5) .. controls (3,4) and (5,3) .. (8,2) node[right, popblue, pos=0.1] {Mfoundi};

% Zone Aléa Fort (Lit mineur / Marécage inondable fréquent)
\fill[popblue!50] (1.5,4.5) .. controls (2.5,3.8) and (4,3.2) .. (5,2.5) .. controls (4,3.5) and (2.5,4.2) .. (1.5,4.5);
\node[popblue, font=\bfseries\tiny] at (3,3.5) {ALÉA FORT};

% Zone Aléa Moyen (Lit majeur)
\fill[poporange!40] (0.5,5.2) .. controls (3,4.2) and (6,3) .. (9,2) .. controls (7,4) and (4,5) .. (1,5.5) .. controls (0.5,5.3) .. cycle;
\node[poporange, font=\bfseries\tiny] at (5,4.2) {ALÉA MOYEN};

% Zone Aléa Faible (Hauteurs) - Fond vert clair déjà couvert par les zones ci-dessus, on ne redessine pas le fond pour éviter la superposition confuse.
% On s'assure que les zones aléa fort/moyen sont par-dessus.

% Enjeux (Bâtiments)
% École en aléa moyen (vulnérable)
\fill[poppink!80] (6.5, 4.0) rectangle ++(0.8, 0.6) node[midway, font=\tiny, rotate=90] {ÉCOLE};
% Hôpital en aléa faible (bon)
\fill[poppink!80] (8.5, 1.0) rectangle ++(1.0, 0.7) node[midway, font=\tiny] {HÔPITAL};
% Habitat dense en aléa fort (TRÈS VULNÉRABLE)
\foreach \x/\y in {2.0/4.2, 2.5/4.0, 3.0/3.8, 3.5/3.5, 2.2/3.8, 2.8/3.6} {
    \fill[poppink!90] (\x,\y) rectangle ++(0.35,0.35);
}
\node[popred, font=\bfseries\tiny] at (2.8, 4.6) {HABITAT PRÉCAIRE (Vulnérabilité FORTE)};

% Habitat dense en aléa moyen
\foreach \x/\y in {5.5/4.5, 6.0/4.5, 5.5/4.0, 7.0/3.5, 7.5/3.5} {
    \fill[poppink!70] (\x,\y) rectangle ++(0.35,0.35);
}

% Digue de protection (Prévention structurelle)
\draw[popblue, ultra thick, decorate, decoration={zigzag, amplitude=1pt}] (1.2,4.6) -- (4.8,2.8) node[midway, above, sloped, popblue, font=\tiny] {Digue de protection (Prévention)};

% Site de repli (Hauteur, Aléa Faible)
\fill[popgold!80] (8.0, 5.0) rectangle ++(1.2, 0.6) node[midway, font=\bfseries\tiny] {SITE REPLI (Stade/École haute)};
\draw[popred, ultra thick, ->] (3.5, 3.8) .. controls (5, 4.5) and (7, 4.8) .. (8.5, 5.2);

% Échelle
\node[font=\tiny] at (9.5, 0.3) {Échelle approx. 1 cm = 200 m};
\end{tikzpicture}
\legende{Exemple simplifié de carte de risque d'inondation (quartier Mvog-Ada, Yaoundé). Superposition de l'aléa (bleu/orange/vert) et des enjeux (rose). Le risque élevé (zone rouge hachurée) combine aléa fort + habitat précaire. La digue (prévention) réduit l'aléa en amont ; le chemin d'évacuation vers le site de repli (or) est la mesure de sauvegarde en cas de crue majeure (prédiction/alerte).}
\end{popfigure}

\courssub{Sensibilisation, éducation et gestes qui sauvent}

La meilleure prévention technique échoue si les populations ne perçoivent pas le risque ou ne savent pas réagir. L'éducation formelle (programmes scolaires SVT, Géographie, Éducation à la Citoyenneté) et non formelle (médias, chefs traditionnels, associations) est un pilier de la \cle{résilience}.

\begin{methode}[Élaborer un outil de sensibilisation local (Affiche / Dépliant / Spot radio) -- Savoir-faire Bac]
\textbf{Objectif :} Produire un support clair, visuel, en langues locales (français, anglais, langues nationales : bassa, ewondo, douala, fulfulde, etc.) pour une catastrophe identifiée (ex: inondation, éruption, glissement).
\textbf{Étapes :}
\begin{enumerate}
    \item \textbf{Identifier le risque cible et le public :} Ex: Inondation quartier Mvog-Ada ; public = ménages chefs de famille, enfants, personnes âgées.
    \item \textbf{Structurer le message selon la chronologie :}
        \begin{itemize}
            \item \textbf{AVANT} : \og Je repère ma zone risque (carte), je prépare mon kit d'urgence (papiers, lampe, eau, médicaments, radio), je connais le point de rassemblement. \fg
            \item \textbf{PENDANT} : \og J'écoute la radio / SMS alerte. Si ordre évacuation : je coupe électricité/gaz, je prends mon kit, je rejoins le point de rassemblement SANS traverser l'eau (30 cm emportent une voiture, 15 cm font tomber un adulte). J'aide mon voisin vulnérable. \fg
            \item \textbf{APRÈS} : \og Je ne rentre que sur autorisation officielle. Je bois de l'eau traitée (chlorée/bouillie). Je signale les dangers (fils électriques, murs fissurés). Je me fais soigner si blessé. \fg
        \end{itemize}
    \item \textbf{Choisir le format et les codes visuels :} Pictogrammes universels (ISO 7010), couleurs codes (Rouge = Danger/Interdiction, Vert = Secours/Issue, Bleu = Obligation/Info), peu de texte, gros caractères.
    \item \textbf{Tester et diffuser :} Faire relire par des non-experts (compréhension ?). Diffuser : affichage mairies/écoles/marchés/églises/mosquées ; spots radio communautaires (CRTV, radios privées) ; réseaux sociaux (WhatsApp, Facebook) ; causeries éducatives chez les chefs de quartier.
\end{enumerate}
\end{methode}

\begin{experience}[Gestes qui sauvent : Conduite à tenir face à une inondation soudaine (Crue éclair) -- TP / Démonstration]
\textbf{Contexte :} Vous êtes en classe (ou à la maison) près d'un cours d'eau (Mfoundi, Wouri, Logone). La pluie tombe violemment depuis 2h. L'eau monte vite, courant fort, couleur boueuse.
\textbf{Matériel :} Vidéos de crue éclair (ex: Douala 2023, Maroua 2012), mannequin, gilet de sauvetage, corde, bidon 20L vide (flotteur).
\textbf{Protocole \& Observation :}
\begin{enumerate}
    \item \textbf{ALERTE :} Entendre la sirène / voir l'eau monter anormalement / recevoir SMS alerte. \textbf{Ne pas attendre l'ordre d'évacuation si le danger est imminent.}
    \item \textbf{MISE EN SÉCURITÉ IMMÉDIATE :}
        \begin{itemize}
            \item Couper électricité (disjoncteur) et gaz (bouteille) \textbf{avant} que l'eau n'atteigne les prises/appareils (risque électrocution, explosion).
            \item Monter à l'étage / sur le toit / sur un meuble haut \textbf{si l'évacuation est impossible} (eau trop rapide, nuit). \textbf{Ne JAMAIS rester au rez-de-chaussée / sous-sol.}
            \item Prendre le \textbf{Kit d'urgence} (sac à dos prêt : papiers d'identité dans sac plastique zip, lampe torche + piles, téléphone chargé + powerbank, eau 1,5L, barre énergétique, médicaments chroniques, sifflet, couverture de survie).
        \end{itemize}
    \item \textbf{ÉVACUATION (si possible et ordonné) :}
        \begin{itemize}
            \item Suivre les \textbf{itinéraires balisés} (panneaux \og Issue de secours \fg / \og Point de rassemblement \fg), pas les raccourcis.
            \item \textbf{Ne JAMAIS traverser une zone inondée} : à pied (15 cm courant fort = chute), en voiture (30 cm = flottaison/embarquement). \textbf{Tourner le dos au courant} si traversé forcé (dos au courant, pas de côté, bâton/perche pour sonder).
            \item Aider les \textbf{personnes vulnérables} (PMR, femmes enceintes, personnes âgées, enfants) : portage, chaise roulante, brancard improvisé (couverture + deux bâtons).
        \end{itemize}
    \item \textbf{POINT DE RASSEMBLEMENT :} Se signaler aux secours (Protection Civile, Croix-Rouge, Chef de quartier). Ne pas disperser. Attendre les consignes officielles pour le retour.
\end{enumerate}
\textbf{Interprétation :} La survie dépend de la \textbf{préparation (kit, connaissance itinéraires)} et de la \textbf{réaction rapide et disciplinée} (ne pas paniquer, ne pas chercher à sauver des biens matériels). L'entraînement régulier (exercices d'évacuation école/entreprise) automatise ces gestes.
\end{experience}

\begin{savaistu}[Les systèmes d'alerte précoce aux tsunamis : leçon du 26 décembre 2004]
Le tsunami de l'océan Indien (séisme magnitude 9,1 au large de Sumatra) a fait \textbf{environ 230 000 morts} dans 14 pays. L'absence de système d'alerte régional a été fatale : les vagues ont mis 1h30 à 7h pour atteindre les côtes (Thaïlande, Sri Lanka, Inde, Afrique de l'Est), un temps suffisant pour évacuer si l'alerte avait existé.
\medskip
\noindent \textbf{Réponse :} Création du \textbf{Système d'Alerte aux Tsunamis dans l'Océan Indien (IOTWMS)} sous l'égide de l'UNESCO/COI.
\begin{itemize}
    \item \textbf{Réseau sismique temps réel :} Détection séisme $\rightarrow$ estimation magnitude/localisation/foyer en $< 10$ min.
    \item \textbf{Bouées DART (Deep-ocean Assessment and Reporting of Tsunamis) :} Mesurent la hauteur de la vague en haute mer (pression au fond) $\rightarrow$ confirment/infirment le tsunami $\rightarrow$ estiment l'amplitude côtière.
    \item \textbf{Marégraphes côtiers :} Mesurent l'arrivée réelle.
    \item \textbf{Chaîne de diffusion :} Centre régional (Jakarta, Mumbai, Perth) $\rightarrow$ Centres nationaux (au Cameroun : DMN / Protection Civile) $\rightarrow$ Population (Sirènes, SMS, Radio, TV).
\end{itemize}
\textbf{Résultat :} Lors du séisme de 2012 (magnitude 8,6, même zone), l'alerte a fonctionné : évacuation préventive des côtes, pas de victimes du tsunami. \textbf{La prévention par la technologie et la coopération internationale sauve des vies.}
\end{savaistu}

\begin{exoresolu}[Analyse d'une mesure de prévention : Le dégazage du Lac Nyos]
\textbf{Énoncé :} Le lac Nyos (Région du Nord-Ouest, Cameroun) a connu une éruption limnique catastrophique le 21 août 1986, libérant $\approx 1,6 \times 10^6$ tonnes de $\ce{CO2}$, causant 1 746 décès humains et 3 500 têtes de bétail. Depuis 2001, un tuyau de dégazage artificiel (puis deux autres en 2011 et 2019) est installé. Expliquez en quoi cette installation est une mesure de \textbf{prévention structurelle} et comment elle modifie le risque résiduel.
\textbf{Solution détaillée :}
\begin{enumerate}
    \item \textbf{Nature du danger (Aléa) :} Accumulation naturelle de $\ce{CO2}$ d'origine magmatique dans les eaux profondes (stratification stable : eau dense gazeuse au fond, eau légère en surface). Risque de renversement brutal (déclencheur : glissement, pluie forte, séisme) $\rightarrow$ éruption limnique (nuage $\ce{CO2}$ dense, asphyxiant, au ras du sol).
    \item \textbf{Principe du dégazage (Prévention structurelle) :} Un tuyau (diamètre $\approx 0,14$ m puis $0,20$ m) descendu jusqu'au fond ($\approx 200$ m). La pression hydrostatique chasse l'eau gazeuse dans le tuyau $\rightarrow$ détente $\rightarrow$ bulles $\rightarrow$ colonne d'eau gazeuse auto-entretenue (geyser continu). Le $\ce{CO2}$ est évacué \textbf{en continu et de façon contrôlée} vers l'atmosphère (dilution rapide, inoffensif).
    \item \textbf{Effet quantitatif sur l'aléa :}
        \begin{itemize}
            \item 2001 (1 tuyau) : Extraction $\approx 6\,000$ t/an. Stock initial $\approx 15 \times 10^6$ t. Taux extraction $<$ taux recharge naturelle ($\approx 10\,000$ t/an). \textbf{Stock augmentait encore.}
            \item 2011 (2 tuyaux ajoutés, total 3) : Extraction $\approx 35\,000$ t/an $>$ Recharge. \textbf{Stock diminue.}
            \item 2019 (Remplacement/optimisation) : Capacité accrue. Objectif : réduire le stock sous le seuil critique de saturation ($< 5 \times 10^6$ t) pour rendre l'éruption \textbf{physiquement impossible} même en cas de déclencheur.
        \end{itemize}
    \item \textbf{Classification :} C'est de la \textbf{prévention structurelle (active)} : elle agit \emph{sur l'aléa lui-même} (réduction de la source de danger), de façon permanente, par une ouvrage d'ingénierie. Elle complète la prévention non structurelle (interdiction d'habiter les vallées aval, sirènes d'alerte, éducation).
    \item \textbf{Risque résiduel :} Tant que le stock $>$ seuil critique, le risque d'éruption spontanée (sans déclencheur externe) ou déclenchée (séisme local) persiste, mais sa \textbf{probabilité diminue chaque année}. La surveillance continue (capteurs $\ce{CO2}$, température, niveau) reste indispensable pour la \textbf{prédiction} d'une éventuelle reprise d'activité.
\end{enumerate}
\textbf{Conclusion :} Le dégazage du Lac Nyos est l'exemple type de prévention structurelle majeure au Cameroun, transformant un aléa majeur incontrôlé en un risque géré et décroissant, fruit de la coopération scientifique internationale (IRD, Université de Savoie, MINRESI, Protection Civile).
\end{exoresolu}

\begin{exercice}[Application : Construire un argumentaire pour un Plan Communal de Sauvegarde (PCS)]
\textbf{Contexte :} Vous êtes conseiller municipal dans une commune rurale du Département du Mayo-Danay (Extrême-Nord), traversée par le fleuve Logone, exposée aux inondations saisonnières (juillet-octobre) et aux vents violents (harmattan, tornades). La commune dispose d'un collège, d'un centre de santé, de 45 villages, pas de sirène, radio communautaire locale.
\textbf{Consignes :}
\begin{enumerate}
    \item En vous appuyant sur la distinction \textbf{Prévision / Prédiction / Prévention}, proposez \textbf{une action concrète par pilier} adaptée à ce contexte (ex: prévision = carte aléa inondation basée sur 30 ans de données hydrométriques Logone à Bongor/Kousseri).
    \item Identifiez \textbf{les 3 acteurs locaux prioritaires} à associer dans la commission locale de sécurité civile (hors mairie) et justifiez leur rôle.
    \item Rédigez le \textbf{message d'alerte type (30 secondes max)} à diffuser sur la radio communautaire en cas de crue imminente (prédiction), en français et en langue locale (ex: \og \emph{Laawol Logone doggol...} \fg en fulfulde), respectant la structure : \textbf{Nature du danger -- Localisation -- Comportement à adopter -- Point de rassemblement}.
\end{enumerate}
\end{exercice}

\begin{corrige}[Exercice : PCS Mayo-Danay]
\begin{enumerate}
    \item \textbf{Prévision (Long terme) :} Réalisation d'une \textbf{carte d'aléa inondation} à l'échelle 1/10 000ème par croisement des données hydrométriques (niveaux Logone 30 ans, courbes de décrue/crue) et topographiques (LiDAR ou drone). Zonage : Zone Rouge (inondable 1 an/2 - interdiction construction), Zone Orange (inondable 1 an/10 - construction sur pilotis obligatoire), Zone Verte (hors aléa). Intégration au Plan Local d'Urbanisme (PLU).
    \item \textbf{Prédiction (Court terme) :} Convention avec la \textbf{DMN (Direction Météorologie Nationale)} et l'\textbf{IRGM} pour réception des bulletins de vigilance crues (modèles hydro-météorologiques) 48h/24h avant. Mise en place d'un \textbf{observateur local} (agent communal / chef village) formé à la lecture d'échelle limnimétrique sur le Logone (seuil alerte orange = 4,50 m à Bongor ; rouge = 5,20 m). Transmission radio/SMS vers mairie.
    \item \textbf{Prévention (Amont) :}
        \begin{itemize}
            \item \textit{Structurelle :} Reboisement des berges du Logone (espèces locales : \emph{Acacia nilotica}, \emph{Mitragyna inermis}) sur 50 m de large pour freiner l'érosion et l'envasement ; construction de \textbf{buttes de refuge} surélevées (1,50 m au-dessus cote crue centennale) dans chaque village (abri bétail + population).
            \item \textit{Non structurelle :} Exercices d'évacuation annuels (mois de juin, avant hivernage) dans les 45 villages ; kits d'urgence familiaux subventionnés (bidon 20L, savon, moustiquaire, sifflet) ; formation 2 secouristes PSC1 par village (Croix-Rouge).
        \end{itemize}
    \item \textbf{Acteurs locaux prioritaires :}
        \begin{enumerate}
            \item \textbf{Chefs traditionnels (Lamibe, Chefs de village) :} Autorité coutumière, connaissance parfaite du terrain et des populations vulnérables, relais de l'alerte porte-à-porte (confiance), mobilisation main-d'œuvre pour digues/buttes.
            \item \textbf{Directeur du Collège / Chef Centre de Santé :} Enjeux majeurs (enfants, malades). Responsables de la mise en sécurité de leurs usagers (plan d'évacuation interne, stock eau/médicaments). Le centre de santé est poste de secours avancé.
            \item \textbf{Animateur Radio Communautaire :} Seul média temps réel couvrant toute la commune. Diffusion alerte, consignes, informations retour (langues locales). Lien permanent avec mairie.
        \end{enumerate}
    \item \textbf{Message d'alerte type (Radio Communautaire) :}
    \begin{quote}
    \textbf{Français :} \og \textbf{ALERTE CRUE ROUGE LOGONE.} Niveau fleuve 5,30 m à Bongor (seuil 5,20 m). Montée rapide. Villages rive gauche (Kousseri, Blangoua, Darak) : \textbf{ÉVACUEZ IMMÉDIATEMENT} vers buttes de refuge. Coupez feu/électricité. Prenez kit urgence. Aidez voisins âgés/handicapés. \textbf{NE TRAVERSEZ PAS L'EAU.} Point rassemblement : Butte refuge village / École surélevée. Écoutez cette radio pour suite. \fg
    \medskip
    \textbf{Fulfulde (exemple) :} \og \textbf{LAAWOL LOGONE DOGGOL.} Logone 5,30 m Bongor. Yola doggol. Suuduji fowru (Kousseri, Blangoua, Darak): \textbf{NANNAA HANNO} suudu nguuli. Tawi wuta/lektrisiti. Dabbere kit urgence. Enndira mawni/horo. \textbf{WADDA LOGONE.} Sukko: Suudu nguuli / Kolaare nder. Dembi radio. \fg
    \end{quote}
    \textit{Critères de réussite :} Court ($< 30$ s), clair, impératif, localisé, comportement ordonné (action verb), point de rassemblement nommé, rappel écoute radio.
\end{enumerate}
\end{corrige}

\coursec{Les gestes qui sauvent et la conduite à tenir}

Dans la section précédente, nous avons vu que la gestion des catastrophes repose sur trois piliers : la \cle{prévision}, la \cle{prédiction} et la \cle{prévention}. Mais lorsque, malgré tout, la catastrophe survient — un séisme qui secoue une ville, une inondation qui envahit un quartier de Yaoundé, un incendie dans un marché —, la survie des personnes dépend alors des \cle{réflexes} et des \cle{gestes} accomplis dans les premières minutes, avant même l'arrivée des secours professionnels. C'est l'objet de cette dernière section : apprendre la \cle{conduite à tenir} face aux principales catastrophes et maîtriser les \cle{gestes qui sauvent}.

\courssub{Les principes généraux du secourisme : Protéger, Alerter, Secourir}

Face à une victime ou à une situation d'urgence, l'angoisse pousse souvent à agir dans le désordre : on se précipite vers le blessé, on oublie d'appeler les secours, ou l'on devient soi-même une victime. Le secourisme repose sur une séquence simple, dans un \textbf{ordre impératif} : \cle{Protéger}, puis \cle{Alerter}, puis \cle{Secourir}.

\begin{methode}[La séquence Protéger -- Alerter -- Secourir]
\begin{enumerate}
\item \textbf{Protéger} : supprimer le danger si possible (couper le courant électrique, éteindre le feu naissant, baliser un accident de la route) ou, à défaut, soustraire la victime au danger \emph{sans se mettre soi-même en danger}. Un secouriste victime à son tour aggrave le bilan et prive les autres d'aide.
\item \textbf{Alerter} : prévenir les secours le plus tôt possible, même si l'on doit ensuite revenir vers la victime. Au Cameroun : le \textbf{118} (sapeurs-pompiers) et le \textbf{119} (SAMU / urgences médicales). Le message doit comporter : la \emph{nature} de l'accident, le \emph{nombre} de victimes, le \emph{lieu exact}, les \emph{risques persistants}. Ne jamais raccrocher le premier : c'est le régulateur qui met fin à l'appel.
\item \textbf{Secourir} : n'intervenir qu'après avoir sécurisé la zone et alerté, en appliquant les gestes de premiers secours appris (libération des voies aériennes, position latérale de sécurité, compression d'une hémorragie).
\end{enumerate}
\end{methode}

\begin{attention}[Ne jamais s'exposer soi-même]
La première règle du secouriste est de \textbf{ne pas devenir une victime supplémentaire}. On ne pénètre pas dans un bâtiment qui s'effondre, on ne touche pas une personne électrisée avant d'avoir coupé le courant, on ne s'approche pas d'une nappe d'hydrocarbure en feu. Protéger passe avant secourir.
\end{attention}

\courssub{Conduite à tenir selon le type de catastrophe}

\textbf{En cas de séisme.} À l'\textbf{intérieur} d'un bâtiment, le danger principal vient de la chute d'objets, des vitres brisées et de l'effondrement des structures légères. La conduite correcte est de se réfugier \textbf{sous un meuble solide} (table, bureau) en protégeant sa tête, ou contre un \textbf{mur porteur}, loin des fenêtres, des placards et des luminaires. On n'utilise jamais les ascenseurs et l'on ne se précipite pas dans les escaliers pendant les secousses. À l'\textbf{extérieur}, il faut s'éloigner des bâtiments (chutes de façades, de tuiles, de verre), des lignes électriques et des arbres, puis se placer à découvert. Après les secousses, on coupe le gaz et l'électricité pour prévenir les incendies, et l'on reste attentif aux \cle{répliques}.

\textbf{En cas d'inondation.} L'eau rend les installations électriques dangereuses : on \textbf{coupe l'électricité} dès que l'eau monte. On gagne ensuite les \textbf{points hauts} (étages, collines) en emportant si possible eau potable, lampe et téléphone. La règle d'or : \textbf{ne jamais traverser à pied ou en voiture une eau courante} — \SI{30}{cm} d'eau en mouvement suffisent à emporter un piéton, \SI{60}{cm} une voiture. Chaque année, au Cameroun, des noyades surviennent lors de la traversée de ravins en crue pendant la saison des pluies.

\textbf{En cas d'incendie.} Les fumées chaudes et toxiques montent vers le plafond : c'est l'intoxication, plus que les flammes, qui tue. On se déplace donc en \textbf{rampant sous les fumées}, idéalement avec un linge humide sur le nez et la bouche. On \textbf{ferme les portes derrière soi} pour ralentir la propagation du feu, on n'utilise \textbf{jamais les ascenseurs} (coupure de courant, effet cheminée), et l'on teste la chaleur d'une porte avec le dos de la main avant de l'ouvrir. Un \textbf{feu naissant} peut être étouffé : couverture jetée sur les flammes, extincteur, sable ; jamais d'eau sur un feu d'huile ou d'origine électrique.

\textbf{En cas d'éruption volcanique ou de nuage de gaz toxique.} L'expérience tragique du \textbf{lac Nyos} (21 août 1986, région du Nord-Ouest) l'a montré : le \ce{CO2}, plus dense que l'air, s'écoule et stagne dans les fonds de vallée. La conduite à tenir est donc de \textbf{fuir vers les hauteurs}, jamais vers les points bas, et de \textbf{se couvrir les voies respiratoires avec un linge humide} pour filtrer partiellement les gaz et poussières. En cas de coulée de lave ou de nuée ardente (Mont Cameroun), on s'écarte perpendiculairement à l'axe d'écoulement en gagnant les reliefs.

\begin{popfigure}
\begin{center}
\begin{tabular}{|c|c|c|}
\hline
\rowcolor{popblueL} \textbf{Séisme} & \textbf{Inondation} & \textbf{Incendie} \\
\hline
\begin{tikzpicture}[scale=0.8]
\draw[thick,popdark] (0,0) rectangle (2.6,1.9);
\draw[thick,popblue] (0.5,0) -- (0.5,0.9) -- (2.1,0.9) -- (2.1,0);
\draw[thick,popblue] (0.5,0.9) -- (0.35,1.05) -- (2.25,1.05) -- (2.1,0.9);
\fill[poporange] (1.3,0.35) circle (0.13);
\draw[thick,poporange] (1.3,0.48) -- (1.3,0.85);
\draw[thick,poporange] (1.3,0.6) -- (0.95,0.75);
\draw[thick,poporange] (1.3,0.6) -- (1.65,0.75);
\node[popdark,font=\scriptsize] at (1.3,-0.35) {Sous une table solide};
\end{tikzpicture}
&
\begin{tikzpicture}[scale=0.8]
\fill[popblue!50] (0,0) rectangle (2.6,0.7);
\draw[popblue,thick] (0,0.7) sin (0.65,0.82) cos (1.3,0.7) sin (1.95,0.58) cos (2.6,0.7);
\draw[thick,popdark] (0.3,0.7) -- (0.9,1.7) -- (1.5,0.7);
\draw[thick,popdark] (1.4,0.7) -- (2.0,1.7) -- (2.6,0.7);
\fill[poporange] (1.0,1.55) circle (0.11);
\draw[thick,poporange] (1.0,1.66) -- (1.0,1.95);
\draw[thick,poporange] (1.0,1.8) -- (0.75,1.7);
\draw[thick,poporange] (1.0,1.8) -- (1.25,1.7);
\node[popdark,font=\scriptsize] at (1.3,-0.35) {Gagner les points hauts};
\end{tikzpicture}
&
\begin{tikzpicture}[scale=0.8]
\fill[popmuted!40] (0,1.2) rectangle (2.6,1.9);
\node[popmuted,font=\scriptsize] at (1.3,1.55) {fumées};
\fill[poporange] (0.7,0.35) circle (0.12);
\draw[thick,poporange] (0.82,0.35) -- (1.5,0.35);
\draw[thick,poporange] (1.0,0.35) -- (1.15,0.6);
\draw[thick,poporange] (1.5,0.35) -- (1.75,0.15);
\draw[thick,poporange] (1.5,0.35) -- (1.75,0.55);
\draw[thick,popred,->] (1.9,0.35) -- (2.5,0.35);
\node[popdark,font=\scriptsize] at (1.3,-0.35) {Ramper sous les fumées};
\end{tikzpicture}
\\
\hline
\end{tabular}
\end{center}
\legende{Pictogrammes des conduites à tenir : en cas de séisme, se réfugier sous un meuble solide ; en cas d'inondation, gagner les points hauts ; en cas d'incendie, ramper sous la couche de fumées chaudes.}
\end{popfigure}

\courssub{Les gestes de premiers secours}

\textbf{Libération des voies aériennes.} Une victime inconsciente allongée sur le dos risque l'obstruction des voies aériennes par la chute de la langue en arrière. Le premier geste consiste à \textbf{basculer doucement la tête en arrière} (une main sur le front, deux doigts sous le menton) pour dégager le passage de l'air, puis à vérifier la respiration : regarder le thorax, écouter, sentir l'air expiré pendant \SI{10}{\second} au maximum.

\textbf{La position latérale de sécurité (PLS).} Si la victime est \textbf{inconsciente mais respire}, on la place sur le côté : c'est la \cle{position latérale de sécurité}.

\begin{popfigure}
\begin{tikzpicture}[scale=0.85]
% Étape 1
\begin{scope}[shift={(0,0)}]
\fill[poporange] (0.2,0.5) circle (0.14);
\draw[thick,poporange] (0.34,0.5) -- (1.5,0.5);
\draw[thick,poporange] (0.6,0.5) -- (0.85,0.85);
\draw[thick,poporange] (0.6,0.5) -- (0.85,0.2);
\draw[thick,poporange] (1.5,0.5) -- (2.0,0.5);
\draw[thick,poporange] (1.5,0.5) -- (2.0,0.75);
\node[popdark,font=\scriptsize,align=center] at (1.1,-0.5) {\textbf{1.} Victime sur le dos,\\ bras en croix};
\end{scope}
% Étape 2
\begin{scope}[shift={(3.2,0)}]
\fill[poporange] (0.2,0.5) circle (0.14);
\draw[thick,poporange] (0.34,0.5) -- (1.5,0.5);
\draw[thick,poporange] (0.6,0.5) -- (0.85,0.85);
\draw[thick,poporange] (0.6,0.5) -- (0.85,0.2);
\draw[thick,poporange] (1.5,0.5) -- (1.85,0.95);
\draw[thick,poporange] (1.5,0.5) -- (2.0,0.5);
\node[popdark,font=\scriptsize,align=center] at (1.1,-0.5) {\textbf{2.} Plier la jambe\\ opposée au secouriste};
\end{scope}
% Étape 3
\begin{scope}[shift={(6.4,0)}]
\fill[poporange] (0.2,0.55) circle (0.14);
\draw[thick,poporange] (0.34,0.55) -- (1.5,0.75);
\draw[thick,poporange] (0.6,0.58) -- (0.9,0.95);
\draw[thick,poporange] (1.5,0.75) -- (1.9,1.15);
\draw[thick,poporange] (1.5,0.75) -- (2.0,0.65);
\draw[thick,popblue,->] (2.15,1.3) arc (120:40:0.55);
\node[popdark,font=\scriptsize,align=center] at (1.1,-0.5) {\textbf{3.} Rouler la victime\\ vers soi};
\end{scope}
% Étape 4
\begin{scope}[shift={(9.6,0)}]
\fill[poporange] (0.2,0.75) circle (0.14);
\draw[thick,poporange] (0.34,0.72) -- (1.5,0.55);
\draw[thick,poporange] (0.55,0.7) -- (0.3,1.1);
\draw[thick,poporange] (1.5,0.55) -- (1.9,0.95);
\draw[thick,poporange] (1.5,0.55) -- (2.0,0.4);
\draw[thick,popblue,->] (0.2,0.95) -- (0.2,1.25);
\node[popdark,font=\scriptsize,align=center] at (1.1,-0.5) {\textbf{4.} Bouche vers le sol,\\ tête basculée en arrière};
\end{scope}
\end{tikzpicture}
\legende{Les quatre étapes de la position latérale de sécurité (PLS) : 1. placer la victime sur le dos, bras en croix ; 2. plier la jambe opposée au secouriste ; 3. saisir l'épaule et le genou et rouler la victime vers soi ; 4. ajuster la position, bouche orientée vers le sol et tête basculée en arrière pour maintenir les voies aériennes libres.}
\end{popfigure}

\begin{savaistu}[Pourquoi la PLS sauve des vies]
Chez une victime inconsciente allongée sur le dos, les muscles sont relâchés : la \textbf{langue bascule en arrière} et obstrue le pharynx, et d'éventuels \textbf{vomissements} peuvent être inhalés dans les poumons (inhalation mortelle). En position latérale de sécurité, la gravité maintient la langue contre la paroi latérale et fait s'écouler les liquides vers l'extérieur : les voies aériennes restent libres. Ce simple changement de position, sans aucun matériel, prévient la mort par asphyxie.
\end{savaistu}

\textbf{Compression d'une hémorragie.} Une hémorragie externe se combat par la \textbf{pression} : on comprime fortement la plaie avec les mains protégées (gants, tissu propre, sachet plastique) ou on fait comprimer la victime elle-même. Si la compression directe est impossible ou insuffisante, on peut comprimer l'\textbf{artère en amont} contre un plan osseux, aux \cle{points de compression} artérielle.

\begin{popfigure}
\begin{tikzpicture}[scale=1.0]
% Silhouette
\fill[popblue!60] (0,2.6) circle (0.35);
\draw[thick,popblue] (0,2.25) -- (0,0.9);
\draw[thick,popblue] (0,2.0) -- (-0.9,1.3);
\draw[thick,popblue] (0,2.0) -- (0.9,1.3);
\draw[thick,popblue] (0,0.9) -- (-0.55,-0.4);
\draw[thick,popblue] (0,0.9) -- (0.55,-0.4);
% Points de compression
\fill[popred] (-0.28,2.35) circle (0.09);
\draw[popred,->] (-0.28,2.35) -- (-2.1,2.6) node[left,font=\scriptsize,align=right,popdark]{1. Artère carotide\\ (cou, plaies du cou)};
\fill[popred] (-0.62,1.62) circle (0.09);
\draw[popred,->] (-0.62,1.62) -- (-2.1,1.5) node[left,font=\scriptsize,align=right,popdark]{2. Artère brachiale\\ (face interne du bras)};
\fill[popred] (0.45,0.95) circle (0.09);
\draw[popred,->] (0.45,0.95) -- (2.3,1.3) node[right,font=\scriptsize,align=left,popdark]{3. Artère fémorale\\ (pli de l'aine)};
\fill[popred] (0.52,0.1) circle (0.09);
\draw[popred,->] (0.52,0.1) -- (2.3,0.0) node[right,font=\scriptsize,align=left,popdark]{4. Artère poplitée\\ (arrière du genou)};
\end{tikzpicture}
\legende{Principaux points de compression artérielle : l'artère est comprimée contre l'os sous-jacent, en amont de la plaie, pour réduire l'afflux sanguin vers l'hémorragie. La compression directe de la plaie reste le geste de première intention.}
\end{popfigure}

\begin{attention}[Le garrot : un geste de dernier recours]
Le \cle{garrot} interrompt totalement la circulation en aval : il expose à une \textbf{nécrose du membre} et à des complications graves au moment du desserrage. Il est donc réservé aux \textbf{hémorragies massives non contrôlables} par la compression (membre arraché, plaie délabrante), et ne doit jamais être posé pour une plaie banale. S'il est posé : ne jamais le desserrer soi-même, le laisser apparent, et \textbf{noter l'heure de pose} (inscrite sur la peau ou communiquée aux secours), car la durée d'ischémie conditionne les décisions médicales.
\end{attention}

\begin{exemplebox}[Une question de minutes]
L'organisme humain adulte contient environ \SI{5}{L} de sang. Une hémorragie artérielle (sang rouge vif, pulsatile, qui jaillit par jets synchrones du cœur) non comprimée peut entraîner la mort en \textbf{3 à 5 minutes}. Or le délai d'arrivée des secours dépasse presque toujours ce laps de temps. Une \textbf{compression immédiate}, effectuée par n'importe quel témoin, maintient le sang dans l'organisme et \textbf{sauve la vie} de la victime : c'est le geste dont l'efficacité par rapport à sa simplicité est la plus extraordinaire.
\end{exemplebox}

\textbf{Protection et surveillance de la victime.} En attendant les secours, on \textbf{couvre la victime} (couverture, vêtement) pour lutter contre le refroidissement, on la \textbf{rassure} en lui parlant, et l'on surveille sa respiration. On ne donne jamais à boire ni à manger à une victime inconsciente ou devant être opérée.

\begin{attention}[Ne pas déplacer un blessé sans danger imminent]
L'erreur fréquente consiste à soulever ou transporter un blessé de la route « pour le mettre à l'abri » ou « l'emmener à l'hôpital ». Or toute mobilisation d'une victime suspectée de \textbf{fracture}, en particulier de la \textbf{colonne vertébrale}, peut transformer une fracture stable en \textbf{lésion de la moelle épinière} et provoquer une paralysie définitive. On ne déplace une victime que si un \textbf{danger imminent et réel} la menace (incendie, explosion, effondrement, noyade). Sinon : on la protège sur place, on la couvre et l'on attend les secours.
\end{attention}

\courssub{Exercice résolu et importance de la formation}

\begin{exoresolu}[Analyser une situation d'urgence]
Énoncé. Sur la route Yaoundé--Douala, tu assistes à un accident : une moto a percuté une voiture. Le motard gît sur la chaussée, inconscient ; il respire et présente une plaie de la cuisse qui saigne abondamment, par jets. Décris, dans l'ordre, les actions à mener en justifiant chaque étape.
\tcblower
Solution détaillée.
\begin{enumerate}
\item \textbf{Protéger} : stationner son véhicule en amont, allumer les feux de détresse, baliser la zone (triangle, branchages) pour éviter une suraccident. On ne touche pas encore la victime.
\item \textbf{Alerter} : appeler le 118 ou le 119 ; préciser le lieu (route Yaoundé--Douala, kilométrage ou repère), la nature (accident de la circulation), le nombre de victimes (une, inconsciente, hémorragie de la cuisse). Ne pas raccrocher le premier.
\item \textbf{Secourir} : le saignement par jets indique une \textbf{hémorragie artérielle} (artère fémorale touchée) : urgence absolue, mort possible en 3 à 5 minutes. On comprime immédiatement et fortement la plaie avec un tissu propre ou les mains protégées ; si la compression directe est insuffisante, on comprime l'artère fémorale au pli de l'aine.
\item La victime étant inconsciente \emph{mais respirant}, on la place en \textbf{position latérale de sécurité} dès que l'hémorragie est contrôlée, afin de protéger ses voies aériennes.
\item On \textbf{ne déplace pas} la victime (risque de fracture du rachis), on la couvre, on la rassure si elle reprend conscience, et l'on surveille sa respiration jusqu'à l'arrivée des secours. Le garrot n'est envisagé qu'en dernier recours si la compression échoue, en notant l'heure de pose.
\end{enumerate}
\end{exoresolu}

\textbf{Se former pour sauver.} Ces gestes ne s'improvisent pas : ils s'apprennent et s'entretiennent. Les \textbf{exercices pratiques} (simulations d'évacuation, alertes-incendie, mises en situation) et les \textbf{formations aux premiers secours} organisées dans les établissements scolaires transforment les élèves en citoyens capables d'agir efficacement. Un lycéen formé, c'est une famille et un quartier mieux protégés : la culture du risque et du secourisme est l'un des piliers concrets de l'éducation à l'environnement et au développement durable, car une communauté résiliente est une communauté préparée.

\begin{aretenir}[L'essentiel de la section]
\begin{itemize}
\item La séquence du secourisme, dans l'ordre impératif : \textbf{Protéger} (sans se mettre en danger), \textbf{Alerter} (118 pompiers, 119 SAMU au Cameroun), \textbf{Secourir}.
\item Séisme : sous une table solide à l'intérieur, à découvert loin des bâtiments à l'extérieur. Inondation : couper l'électricité, gagner les points hauts, ne jamais traverser une eau courante. Incendie : ramper sous les fumées, fermer les portes, pas d'ascenseur. Nuage de gaz : fuir vers les hauteurs, linge humide sur le visage.
\item Victime inconsciente qui respire : \textbf{position latérale de sécurité}. Hémorragie : \textbf{compression} immédiate ; le garrot est un dernier recours dont on note l'heure de pose.
\item On ne déplace jamais un blessé suspecté de fracture, sauf danger imminent.
\end{itemize}
\end{aretenir}

\newpage

\coursec{Activité d'intégration}

\courssub{Objectifs de l'activité}
\begin{itemize}
    \item Mobiliser les connaissances sur la \cle{classification des catastrophes} (origine naturelle, humaine, mixte) pour qualifier l'événement de Bafoussam 2019.
    \item Exploiter des \cle{données pluviométriques} (courbe mensuelle) et des \cle{documents cartographiques} (pentes urbanisées) pour identifier les \cle{facteurs de vulnérabilité} et la \cle{saison à haut risque}.
    \item Proposer des \cle{mesures de prévention} structurelles et non structurelles adaptées au contexte local.
    \item Concevoir un \cle{outil de sensibilisation} (affiche ou dépliant) intégrant la classification, les \cle{signes avant-coureurs} d'un glissement de terrain et les \cle{gestes qui sauvent}.
    \item Développer des compétences de \cle{communication scientifique}, de \cle{travail collaboratif} et d'\cle{esprit critique} face à la gestion des risques.
\end{itemize}

\courssub{Situation}
\textbf{Contexte réel :} Dans la nuit du \textbf{28 au 29 octobre 2019}, après plusieurs jours de pluies diluviennes, un important \cle{glissement de terrain} (coulée de boue et de débris) s'est produit dans le quartier \textbf{Ngouache} à \textbf{Bafoussam}, chef-lieu de la région de l'Ouest du Cameroun. Le bilan officiel fait état de \textbf{43 décès}, de nombreux blessés, de \textbf{plus de 20 maisons ensevelies} et de centaines de sinistrés. Ce drame a profondément marqué la population et les autorités.

\textbf{Dossier documentaire mis à votre disposition (extraits) :}

\textbf{Document 1 : Article de presse « Bafoussam : le lourd tribut de la pluie » (Cameroon Tribune, 30 octobre 2019).}
\begin{quote}
« ... Les secours ont travaillé sans relâche pour déblayer les gravats. Les témoins racontent avoir entendu un grondement sourd vers 03 h du matin, suivi d'un nuage de poussière. Le quartier Ngouache est situé en contrebas d'une colline à forte pente, où l'urbanisation anarchique a conduit à la suppression de la couverture végétale et à l'imperméabilisation des sols. Les caniveaux d'évacuation des eaux pluviales, là où ils existent, sont obstrués par les déchets solides... »
\end{quote}

\textbf{Document 2 : Carte simplifiée des pentes urbanisées du quartier Ngouache (extrait Plan Local d'Urbanisme - PLU).}
\begin{center}
\begin{tabularx}{\linewidth}{|l|X|}\hline
\textbf{Légende} & \textbf{Description} \\ \hline
\cellcolor{popred!30} Zone rouge & Pentes $> 30\%$ (\SI{> 17}{\degree}), habitat dense spontané, absence de réseau d'assainissement \\ \hline
\cellcolor{poporange!30} Zone orange & Pentes $15\%$ à $30\%$ (\SIrange{8.5}{17}{\degree}), habitat mixte, caniveaux saturés \\ \hline
\cellcolor{popgreen!30} Zone verte & Pentes $< 15\%$ (\SI{< 8.5}{\degree}), zone relativement stable, infrastructures publiques \\ \hline
\end{tabularx}
\end{center}
\textit{Note :} Le point de rupture initial du glissement (couronne de décrochement) se situe en limite Zone rouge / Zone orange, à une altitude d'environ \SI{1450}{m}. La coulée a parcouru \SI{\approx 300}{m} horizontalement pour un dénivelé de \SI{\approx 120}{m}.

\textbf{Document 3 : Courbe mensuelle des précipitations moyennes à Bafoussam (période 1981--2010, source ASECNA / Météo Cameroun).}
\begin{popfigure}
\begin{tikzpicture}
\begin{axis}[
    width=\linewidth, height=7cm,
    xlabel={Mois}, ylabel={Précipitations moyennes (mm)},
    xtick={1,2,3,4,5,6,7,8,9,10,11,12},
    xticklabels={J,F,M,A,M,J,J,A,S,O,N,D},
    ymin=0, ymax=300,
    ymajorgrids=true, grid style=dashed,
    popcurve, thick, mark=*, mark size=2pt
]
\addplot coordinates {
(1,15) (2,25) (3,85) (4,155) (5,210) (6,205)
(7,180) (8,240) (9,285) (10,260) (11,95) (12,20)
};
\addplot[poporange, dashed, thick, mark=square*, mark size=2pt] coordinates {
(1,10) (2,30) (3,110) (4,180) (5,230) (6,220)
(7,190) (8,260) (9,310) (10,290) (11,110) (12,15)
};
\legend{Normales 1981-2010, Année 2019 (estimée)}
\end{axis}
\end{tikzpicture}
\legende{Évolution mensuelle des précipitations à Bafoussam. La courbe en pointillés représente l'année 2019, caractérisée par un excédent pluviométrique marqué en saison des pluies (juin--octobre).}
\end{popfigure}

\textbf{Document 4 : Photographies de terrain (prises en novembre 2019).}
\begin{itemize}
    \item Photo A : Vue aérienne du cône de déjection recouvrant les habitations ; on distingue la cicatrice de décrochement en amont (sol latéritique nu).
    \item Photo B : Détail d'une maison en briques de terre crue (banco) dont le mur de soutènement a cédé ; fondations affleurantes.
    \item Photo C : Caniveau principal obstrué par des déchets plastiques et des sédiments ; eau stagnante en surface.
\end{itemize}

\courssub{Consignes}
Travaillant en groupe de 4 à 5 élèves, vous disposez de \textbf{2 heures} pour réaliser les tâches suivantes et préparer une présentation orale de \textbf{5 minutes} :

\begin{enumerate}
    \item \textbf{Analyse documentaire :} À partir des documents 1 à 4, identifiez l'\cle{origine} de la catastrophe (naturelle, humaine, mixte) et justifiez votre réponse. Listez \textbf{au moins 4 facteurs de vulnérabilité} (physiques, anthropiques, organisationnels) qui ont aggravé le bilan.
    \item \textbf{Interprétation climatique :} En exploitant le Document 3, déterminez la \textbf{période de l'année à haut risque} de glissements de terrain à Bafoussam. Calculez le \textbf{surplus pluviométrique} (en \%) du mois de septembre 2019 par rapport à la normale 1981--2010. Expliquez le rôle de la pluviométrie dans le déclenchement (mécanisme hydrogéologique).
    \item \textbf{Propositions de prévention :} Formulez \textbf{trois mesures de prévention} concrètes et hiérarchisées (une \emph{structurelle}, une \emph{non structurelle / réglementaire}, une \emph{basée sur la sensibilisation / éducation}), adaptées au contexte socio-économique de Bafoussam. Pour chaque mesure, précisez l'acteur principal responsable de sa mise en œuvre (Mairie, MINHDU, Population, ONG, etc.).
    \item \textbf{Production de l'outil de sensibilisation :} Concevez une \textbf{affiche} (format A3 paysage) ou un \textbf{dépliant} (3 volets A4) intitulé \emph{« Bafoussam : plus jamais ça »}. Votre production doit impérativement comporter :
    \begin{itemize}
        \item La \cle{classification} de la catastrophe de 2019 (type, origine, ampleur).
        \item Les \textbf{5 signes avant-coureurs} visuels d'un glissement de terrain imminent (ex. : fissures au sol, arbres penchés, eau trouble aux sources, etc.).
        \item Les \textbf{gestes qui sauvent} \textbf{AVANT} (évacuation préventive), \textbf{PENDANT} (mise à l'abri verticale, ne pas traverser le courant) et \textbf{APRÈS} (alerte secours, ne pas rentrer chez soi sans avis d'expert).
        \item Un message clair d'appel à l'action (numéros d'urgence : \textbf{112} / \textbf{119}).
    \end{itemize}
    \item \textbf{Évaluation par les pairs :} Chaque groupe présente sa production. La classe utilise la \textbf{grille critériée} ci-après pour noter sur 20 points.
\end{enumerate}

\courssub{Ressources à mobiliser}
\begin{propriete}[Classification des catastrophes (Rappel de cours)]
Une catastrophe est un événement brutal, imprévisible ou non, causant des dommages graves (humains, matériels, environnementaux) et dépassant les capacités de réponse locales.
\begin{itemize}
    \item \textbf{Origine naturelle :} Séisme, éruption volcanique, cyclone, inondation \emph{strictement} climatique, glissement de terrain \emph{sur pente naturelle vierge}.
    \item \textbf{Origine humaine (technologique/anthropique) :} Accident industriel (usine, barrage), pollution majeure, incendie urbain, effondrement d'ouvrage d'art par défaut d'entretien.
    \item \textbf{Origine mixte (naturelle aggravée par l'Homme) :} Phénomène naturel déclencheur (pluie, séisme) dont la gravité est amplifiée par l'action humaine (déforestation, urbanisation en zone à risque, défaut d'aménagement, changement climatique). \textbf{C'est le cas le plus fréquent au Cameroun.}
\end{itemize}
\end{propriete}

\begin{propriete}[Facteurs de vulnérabilité]
La vulnérabilité ($V$) caractérise la propension d'un élément (population, bâtiment, infrastructure) à subir des dommages face à un aléa ($A$). Le risque $R = A \times V \times E$ (Enjeu). Facteurs clés :
\begin{itemize}
    \item \textbf{Physiques :} Pente forte, nature du sol (argile, latérite), absence de végétation, sismicité.
    \item \textbf{Anthropiques :} Habitat précaire (banco, tôle), occupation de zones non constructibles, imperméabilisation, déchets obstruant l'écoulement.
    \item \textbf{Organisationnels :} Absence de Plan de Prévention des Risques (PPR), manque de système d'alerte précoce, méconnaissance des gestes qui sauvent.
\end{itemize}
\end{propriete}

\begin{methode}[Mécanisme de déclenchement d'un glissement de terrain pluvial]
\begin{enumerate}
    \item Infiltration massive de l'eau de pluie $\rightarrow$ saturation du sol (pores remplis d'eau).
    \item Augmentation de la \textbf{pression interstitielle} (pression de l'eau dans les pores) $\rightarrow$ diminution de la \textbf{contrainte efficace} ($\sigma' = \sigma - u$) qui assure la cohésion du sol.
    \item Perte de la résistance au cisaillement du sol (critère de Mohr-Coulomb simplifié : $\tau = c' + \sigma' \tan \phi'$).
    \item Rupture le long d'une surface de glissement (souvent circulaire ou plane) $\rightarrow$ mouvement de la masse (coulée, glissement, éboulement).
\end{enumerate}
\end{methode}

\begin{aretenir}[Gestes qui sauvent : Glissement de terrain / Coulée de boue]
\begin{itemize}
    \item \textbf{AVANT (Signes avant-coureurs) :} Fissures récentes sur murs/sol, portes/fenêtres qui coincent, arbres/poteaux penchés, eau jaillissant anormalement sur pente, bruit sourd (craquement roches). \textbf{Alerter (112/119) et évacuer vers zone haute stable.}
    \item \textbf{PENDANT :} Se déplacer \textbf{perpendiculairement} à l'axe de la coulée (vers les côtés), gagner un point haut solide (rocher, étage supérieur bâtiment solide). \textbf{Ne jamais courir dans l'axe de la coulée ni traverser le dépôt.}
    \item \textbf{APRÈS :} Ne pas regagner son domicile tant que les experts (Génie civil, Protection civile) n'ont pas déclaré la zone stable. Risque de reprise, gaz toxiques, réseaux électriques coupés. Signaler les personnes disparues.
\end{itemize}
\end{aretenir}

\courssub{Démarche et corrigé commenté}
\begin{exoresolu}[Activité d'intégration : Campagne « Bafoussam : plus jamais ça »]
\textbf{Étape 1 : Analyse documentaire et classification (Consigne 1)}

\textbf{Origine de la catastrophe :} \textbf{Mixte (Naturelle aggravée par l'Homme).}
\textbf{Justification :}
\begin{itemize}
    \item \textbf{Facteur naturel déclencheur (Aléa) :} Pluies exceptionnelles (Doc 3 : \SI{310}{mm} en sept. 2019 vs \SI{285}{mm} normale ; excédent sur juin--oct). Saturation hydrique des sols latéritiques (horizon argilique kaolinitique/halloysitique) $\rightarrow$ augmentation de la pression interstitielle et perte de cohésion (Méthode ci-dessus).
    \item \textbf{Facteurs humains aggravants (Vulnérabilité) :}
    \begin{enumerate}
        \item \textbf{Urbanisation anarchique en zone à risque :} Occupation de pentes $> 30\%$ (Zone rouge Doc 2) interdites à la construction (Code de l'urbanisme, Loi n°2004/003).
        \item \textbf{Déforestation / Dénudation des sols :} Suppression de la couverture végétale (racines fixatrices) $\rightarrow$ érosion, infiltration rapide, perte de cohésion superficielle.
        \item \textbf{Imperméabilisation \& Défaut d'assainissement :} Toitures, sols en dur, caniveaux obstrués (Doc 4, Photo C) $\rightarrow$ ruissellement concentré, surcharge hydrique locale, affouillement au pied des pentes.
        \item \textbf{Habitat précaire :} Maisons en banco (terre crue) sans fondations profondes ni chaînages, sensibles à l'eau (Doc 4, Photo B). Absence de murs de soutènement dimensionnés.
    \end{enumerate}
\end{itemize}

\textbf{Étape 2 : Interprétation climatique (Consigne 2)}

\begin{itemize}
    \item \textbf{Période à haut risque :} \textbf{Juillet à Octobre} (cœur de la grande saison des pluies). Pic en \textbf{Septembre} (\SI{285}{mm} normale, \SI{310}{mm} en 2019). C'est la période de saturation maximale des sols (antécédents pluviométriques juin-août).
    \item \textbf{Calcul du surplus pluviométrique de Septembre 2019 :}
    \[
    \text{Surplus } (\%) = \frac{P_{2019} - P_{\text{normale}}}{P_{\text{normale}}} \times 100
    = \frac{310 - 285}{285} \times 100
    = \frac{25}{285} \times 100
    \approx \textbf{8,8\%}
    \]
    \textit{Note :} Ce surplus modéré en valeur relative masque un excédent cumulé sur la saison (juin-oct 2019 $\approx$ \SI{1270}{mm} vs \SI{1170}{mm} normale, soit $+8,5\%$), suffisant pour saturer la zone non saturée sur une épaisseur critique.
    \item \textbf{Rôle hydrogéologique :} L'eau de pluie infiltre la latérite (perméable en surface) jusqu'à l'horizon argilique moins perméable. La nappe perchée monte, la pression interstitielle $u$ augmente, la contrainte efficace $\sigma'$ chute, le facteur de sécurité $F = \frac{\text{Résistance}}{\text{Poussée}}$ passe sous 1.
\end{itemize}

\textbf{Étape 3 : Propositions de prévention hiérarchisées (Consigne 3)}

\begin{center}
\begin{tabularx}{\linewidth}{|l|X|X|X|}\hline
\rowcolor{popblueL}
\textbf{Type de mesure} & \textbf{Mesure concrète} & \textbf{Acteur principal} & \textbf{Justification / Faisabilité} \\ \hline
\textbf{Structurelle} (Investissement lourd, effet durable) & \textbf{Stabilisation physique des pentes critiques (Zone rouge) :} Construction de \textbf{murs de soutènement en béton armé} (berlinois ou en L) au pied des couronnes de décrochement + \textbf{drainage profond} (drains horizontaux, puits de décompression) pour abaisser la nappe perchée. & \textbf{MINHDU} (Ministère de l'Habitat et du Développement Urbain) + \textbf{Mairie de Bafoussam} (Maîtrise d'ouvrage) / Financement \textbf{FEICOM} / \textbf{Banque Mondiale}. & Traite la cause physique (poussée des terres, pression d'eau). Coût élevé mais indispensable pour zones déjà urbanisées qu'on ne peut reloger immédiatement. \\ \hline
\textbf{Non structurelle / Réglementaire} (Cadre légal, planification) & \textbf{Application stricte du Plan Local d'Urbanisme (PLU) :} \textbf{Interdiction formelle de tout nouveau permis de construire} en Zone rouge (pente $>30\%$) et Zone orange non assainie. \textbf{Opération de relogement progressif} des populations les plus exposées (Zone rouge) vers sites sécurisés (ex. site de \emph{Kouoptamo} ou extension urbaine plane), avec indemnisation juste. & \textbf{Mairie} (Police administrative du maire, Art. 25 Loi 2004/017) + \textbf{MINDCAF} (Domaine) + \textbf{MINATD} (Administration territoriale). & S'attaque à la cause racine : l'enjeu humain en zone d'aléa fort. Nécessite volonté politique, foncier disponible, budget relogement. \\ \hline
\textbf{Sensibilisation / Éducation} (Faible coût, effet comportemental) & \textbf{Programme « Quartier Vigilant » :} Formation de \textbf{Comités de Veille Communautaires (CVC)} par quartier (10--15 bénévoles formés par Protection Civile / Croix-Rouge). Missions : surveillance signes avant-coureurs (fissures, eau), nettoyage hebdomadaire caniveaux (journées citoyennes), diffusion alertes SMS/WhatsApp, organisation exercices d'évacuation annuels. & \textbf{Protection Civile} (MINATD) + \textbf{Croix-Rouge Camerounaise} + \textbf{ONG locales} + \textbf{Chefs de quartiers/Blocs}. & Renforce la résilience locale, coût faible, appropriation communautaire. Complète les mesures structurelles lentes à mettre en œuvre. \\ \hline
\end{tabularx}
\end{center}

\textbf{Étape 4 : Production de l'outil de sensibilisation (Consigne 4) — Guide de réalisation pour l'affiche « Bafoussam : plus jamais ça »}

\textbf{Structure recommandée (A3 Paysage - 3 colonnes logiques) :}

\begin{center}
\begin{tabularx}{\linewidth}{|X|X|X|}\hline
\rowcolor{popblueL}
\textbf{COLONNE 1 : COMPRENDRE (Classification \& Aléa)} & \textbf{COLONNE 2 : SURVEILLER (5 Signes Avant-Coureurs)} & \textbf{COLONNE 3 : AGIR (Gestes Qui Sauvent)} \\ \hline
\vspace{0.2cm}
\textbf{TITRE :} BAFOUSSAM : PLUS JAMAIS ÇA ! \\
\vspace{0.2cm}
\textbf{Sous-titre :} Glissement de terrain de Ngouache (Oct. 2019) : 43 vies brisées. \\
\vspace{0.3cm}
\textbf{CLASSIFICATION :} \\
\cle{Catastrophe d'origine mixte} (Naturelle aggravée par l'Homme). \\
\cle{Aléa :} Pluies extrêmes (Saison des pluies Juil--Oct). \\
\cle{Vulnérabilité :} Habitat sur pentes $>30\%$, déforestation, caniveaux bouchés, maisons en banco. \\
\cle{Ampleur :} Locale (quartier), mortelle (43 morts), destructrice (20+ maisons). \\
\vspace{0.3cm}
\textbf{POURQUOI ICI ?} Schéma simplifié (TikZ) : Pente forte + Pluie + Pas d'arbres + Caniveaux bouchés = Sol saturé = GLISSEMENT. \\
\vspace{0.2cm}
\textbf{SAISON À RISQUE :} \textbf{JUILLET -- OCTOBRE} (Pic Septembre). &
\vspace{0.2cm}
\textbf{⚠ 5 SIGNES QUI DOIVENT VOUS ALERTER :} \\
\vspace{0.1cm}
\textbf{1. FISSURES :} Apparition soudaine de fissures diagonales sur les murs, au sol, sur la route. \\
\textbf{2. ARBRES / POTEAUX PENCHÉS :} Trunks courbés vers le bas de pente, poteaux électriques inclinés. \\
\textbf{3. EAU ANORMALE :} Sources nouvelles, eau trouble/jaillissante sur la pente, ruissellement intense hors caniveaux. \\
\textbf{4. BRUITS SOURDS :} Craquements de roches, grondement sourd (comme un camion) venant de la colline, surtout la nuit. \\
\textbf{5. PORTES / FENÊTRES COINCENT :} Déformation des cadres due au mouvement lent du sol (fluage). \\
\vspace{0.2cm}
\textbf{SI VOUS OBSERVEZ 1 SIGNE :} \\
$\rightarrow$ \textbf{ALERTEZ IMMÉDIATEMENT :} \textbf{112} (Urgences) / \textbf{119} (Protection Civile) / Chef de Quartier. \\
$\rightarrow$ \textbf{ÉVACUEZ VERS LE HAUT} (zone stable, rocher, bâtiment R+2 solide) \textbf{SANS ATTENDRE}. &
\vspace{0.2cm}
\textbf{🛡 LES GESTES QUI SAUVENT :} \\
\vspace{0.1cm}
\textbf{AVANT (Prévention / Alerte) :} \\
$\bullet$ Connaître les zones à risque (carte Mairie). \\
$\bullet$ Nettoyer \textbf{chaque semaine} les caniveaux (Journée citoyenne). \\
$\bullet$ Planter des arbres (vetiver, eucalyptus) en pied de pente. \\
$\bullet$ Repérer son \textbf{point de rassemblement haut} (école, église, rocher). \\
$\bullet$ Sac d'urgence prêt (papiers, eau, lampe, radio, médicaments). \\
\vspace{0.2cm}
\textbf{PENDANT (Survie immédiate) :} \\
$\bullet$ \textbf{FUYEZ PERPENDICULAIREMENT} à la coulée (latéralement), \textbf{JAMAIS DANS L'AXE}. \\
$\bullet$ Montez \textbf{VITE} vers le point haut le plus proche. \\
$\bullet$ Si coincé : \textbf{TOIT TERRASSE} ou \textbf{ÉTAGE SUPÉRIEUR} bâtiment solide. \\
$\bullet$ Protégez la tête (casque, sac, bras). \\
$\bullet$ N'essayez \textbf{JAMAIS} de traverser la boue (piège mortel). \\
\vspace{0.2cm}
\textbf{APRÈS (Sécurité / Solidarité) :} \\
$\bullet$ \textbf{NE RENTREZ PAS} chez vous sans feu vert \textbf{Génie Civil / Protection Civile}. \\
$\bullet$ Signalez disparus au \textbf{112 / 119}. \\
$\bullet$ Méfiez-vous : gaz (fosse septique), électricité, reprise de glissement, maladies (eau souillée). \\
$\bullet$ Entraide : hébergement, don sang (Centre Régional Transfusion Sanguine Bafoussam). \\
\vspace{0.3cm}
\textbf{NUMÉROS D'URGENCE :} \textbf{112} (SAMU/Pompiers) --- \textbf{119} (Protection Civile) --- Mairie Bafoussam : \textbf{2 33 42 XX XX} \\
\textbf{Parrainage :} Mairie de Bafoussam --- MINATD --- Croix-Rouge --- OnBuch+ SVT Terminale C/TI
\end{tabularx}
\end{center}

\textbf{Étape 5 : Grille d'évaluation par les pairs (Consigne 5) — Sur 20 points}

\begin{center}
\begin{tabularx}{\linewidth}{|l|X|c|}\hline
\rowcolor{popblueL}
\textbf{Critère} & \textbf{Indicateurs de réussite} & \textbf{Points} \\ \hline
\textbf{Exactitude scientifique (6 pts)} & Classification correcte (Mixte), 4 facteurs vulnérabilité justifiés, calcul surplus correct ($8,8\%$), mécanisme hydrogéologique exact (pression interstitielle), 5 signes avant-coureurs pertinents, gestes AVANT/PENDANT/APRÈS conformes aux consignes officielles. & /6 \\ \hline
\textbf{Pertinence des mesures prévention (4 pts)} & 3 mesures distinctes (Structurelle, Réglementaire, Éducative), adaptées contexte Bafoussam, acteurs identifiés, hiérarchisation logique. & /4 \\ \hline
\textbf{Qualité de l'outil de sensibilisation (6 pts)} & Lisibilité (polices, contraste), exhaustivité (tous éléments demandés présents), clarté messages (pictogrammes, codes couleur danger/sécurité), adaptation public cible (langue française simple, visuels parlants), format respecté (A3 ou 3 volets). & /6 \\ \hline
\textbf{Présentation orale \& Posture (4 pts)} & Respect temps (5 min), participation équitable groupe, aisance expression, réponse pertinente aux questions classe, argumentaire solide. & /4 \\ \hline
\textbf{TOTAL} & & \textbf{/20} \\ \hline
\end{tabularx}
\end{center}
\tcblower
\textbf{Commentaire pédagogique du corrigé :}
Cette activité d'intégration est conçue pour être \textbf{situationnelle} (ancrée dans le drame réel de Bafoussam 2019) et \textbf{complexe} (mobilisant classification, analyse de données chiffrées, raisonnement physico-chimique, réglementation urbaine, communication de crise). Le corrigé type fourni ici sert de \textbf{référence pour l'enseignant} lors de l'évaluation des productions élèves. On attend des élèves qu'ils fassent le lien explicite entre la \textbf{courbe pluviométrique} (Doc 3 : saturation saisonnière) et la \textbf{carte des pentes} (Doc 2 : localisation aléa/vulnérabilité) pour expliquer la \textbf{localisation précise} du drame (Ngouache). La distinction entre \emph{prévision} (probabilité statistique, ex. « septembre est risqué »), \emph{prédiction} (alerte courte échéance, ex. « pluie diluvienne annoncée pour cette nuit, risque imminent ») et \emph{prévention} (mesures structurelles/non structurelles) doit être claire dans l'oral. L'affiche doit être \textbf{opérationnelle} : un habitant de Ngouache doit pouvoir, en 30 secondes de lecture, savoir \emph{quoi regarder} (signes), \emph{qui appeler} (112/119) et \emph{où courir} (haut, latéralement).
\end{exoresolu}

\courssub{Bilan de l'activité}
Cette activité d'intégration a permis de mettre en évidence, par une démarche d'investigation complète, que :
\begin{enumerate}
    \item La \textbf{plupart des « catastrophes naturelles » au Cameroun sont en réalité des catastrophes d'origine mixte}, où l'aléa naturel (pluies intenses, séismes, éruptions) rencontre une vulnérabilité anthropique majeure (urbanisation non maîtrisée, pauvreté du bâti, déforestation, gestion des déchets défaillante). Le cas de Bafoussam 2019 en est l'archétype.
    \item La \cle{prévention} ne se résume pas à des travaux d'ingénierie coûteux (murs, drains) ; elle repose sur un \textbf{trio indissociable} : \textbf{Aménagement du territoire} (réglementation, PLU, relogement), \textbf{Ingénierie de protection} (drainage, stabilisation) et \textbf{Culture du risque} (sensibilisation, veille communautaire, gestes qui sauvent). L'échec d'un seul maillon (ex. caniveaux bouchés malgré murs de soutènement) peut entraîner la catastrophe.
    \item L'\cle{interprétation de données chiffrées} (courbes pluviométriques, calculs de surplus, pentes en \%) est indispensable pour objectiver le risque, définir les seuils d'alerte et prioriser les zones d'intervention (zonage aléa).
    \item La \cle{sensibilisation} efficace doit être \textbf{localisée} (langue, contexte, numéros d'urgence locaux), \textbf{visuelle} (pictogrammes pour signes avant-coureurs), \textbf{actionnable} (gestes simples, point de rassemblement connu) et \textbf{régulièrement réactivée} (exercices, journées nettoyage), non limitée à une campagne ponctuelle post-catastrophe.
    \item En tant que futurs citoyens, techniciens, ingénieurs ou décideurs, vous avez la responsabilité de \textbf{refuser l'installation en zone à risque}, de \textbf{signaler les anomalies} (fissures, caniveaux bouchés) et de \textbf{relayer les consignes de survie}. La résilience d'un territoire se construit \textbf{avant} la catastrophe.
\end{enumerate}
\textbf{Compétences évaluées (Programme MINESEC) :}
\begin{itemize}
    \item \textbf{Savoir :} Classer une catastrophe, expliquer le mécanisme d'un glissement, définir prévision/prédiction/prévention.
    \item \textbf{Savoir-faire :} Analyser un dossier documentaire complexe, exploiter une courbe climatique, élaborer un outil de communication à visée préventive, travailler en équipe.
    \item \textbf{Savoir-être :} Responsabilité citoyenne, solidarité, rigueur scientifique, esprit critique face à l'aménagement du territoire.
\end{itemize}

\newpage

\coursec{Méthodes et exercices résolus}

\coursec{Les catastrophes : classification, gestion et prévention}

\courssub{Définition et caractéristiques d'une catastrophe}

\begin{definition}[Catastrophe]
Une \cle{catastrophe} est un événement brutal ou progressif, d'origine naturelle ou humaine, qui provoque des dommages considérables — pertes en vies humaines, destructions matérielles, dégradation durable de l'environnement — et dont l'ampleur dépasse les capacités de réaction et de réparation de la communauté touchée, nécessitant souvent une aide extérieure.
\end{definition}

\begin{propriete}[Les trois composantes du risque]
Le risque de catastrophe résulte de l'interaction de trois facteurs :
\[
\text{Risque} = \text{Aléa} \times \text{Vulnérabilité} \times \text{Enjeu}
\]
\begin{itemize}
\item \textbf{Aléa} : probabilité d'occurrence, en un lieu et une période donnés, d'un phénomène potentiellement dangereux (séisme, inondation, éruption, accident industriel) d'une intensité donnée.
\item \textbf{Vulnérabilité} : degré de fragilité des populations, des constructions, des activités et de l'environnement exposés à l'aléa (habitat précaire, absence de normes parasismiques, déforestation, insalubrité).
\item \textbf{Enjeu} : valeur des biens, du nombre de vies humaines et des activités économiques présents dans la zone exposée.
\end{itemize}
Une catastrophe survient lorsque l'aléa rencontre une zone à forte vulnérabilité et à fort enjeu.
\end{propriete}

\begin{exemplebox}[Le drame du lac Nyos, 21 août 1986]
Le lac Nyos (région du Nord-Ouest, Cameroun) a libéré brutalement un nuage de \ce{CO2} d'origine magmatique (éruption limnique). Bilan : \num{1746} morts, plus de \num{3000} têtes de bétail tuées, villages dévastés sur \SI{20}{km}. L'aléa (dégazage) était naturel ; la vulnérabilité (habitat en zone basse sans surveillance) et l'enjeu (population dense) ont fait la catastrophe.
\end{exemplebox}

\courssub{Classification des catastrophes}

\begin{methode}[Classer une catastrophe : naturelle, humaine ou mixte]
Pour classer rigoureusement une catastrophe, suis ces quatre étapes :
\begin{enumerate}
\item \textbf{Identifier l'agent causal} : quel phénomène ou action est à l'origine directe (séisme, inondation, explosion, déversement chimique, épidémie, conflit) ?
\item \textbf{Déterminer l'origine de l'agent} :
      \begin{itemize}
      \item Processus naturels (tectonique, volcanisme, climat, microorganismes) $\Rightarrow$ \cle{catastrophe naturelle}.
      \item Activité, négligence ou décision humaine (industrie, transport, guerre, déforestation, urbanisation anarchique) $\Rightarrow$ \cle{catastrophe d'origine humaine} (technologique, sanitaire, sociale, environnementale).
      \end{itemize}
\item \textbf{Vérifier les critères de catastrophe} : dommages importants (pertes humaines, destructions, dégradation durable) dépassant les capacités locales de réponse. Sinon, c'est un simple accident ou aléa.
\item \textbf{Nuancer : les catastrophes mixtes}. Un aléa naturel (pluies diluviennes) devient catastrophe par la vulnérabilité humaine (habitat en zone inondable, déforestation). Au Bac, précise toujours cette nuance.
\end{enumerate}
\textbf{Réflexe rédactionnel} : « L'événement X est une catastrophe d'origine ... car l'agent causal est ... ; il a entraîné ... (bilan chiffré) ».
\end{methode}

\courssub{Les catastrophes d'origine naturelle}

\begin{definition}[Catastrophe naturelle]
Catastrophe dont l'agent causal relève exclusivement des processus géologiques, climatiques ou biologiques internes à la Terre ou au vivant, sans action humaine déclenchante.
\end{definition}

\begin{propriete}[Principaux types de catastrophes naturelles]
\begin{center}
\begin{tabularx}{\linewidth}{|l|X|X|}\hline
\rowcolor{popblueL} \textbf{Type} & \textbf{Aléa principal} & \textbf{Exemples au Cameroun} \\ \hline
\cle{Géologiques} & Séismes, éruptions volcaniques, éruptions limniques, glissements de terrain & Éruptions du Mont Cameroun (1999, 2000) ; Éruption limnique du lac Nyos (1986) ; Glissements de Bafoussam (oct. 2019, $>40$ morts) \\ \hline
\cle{Climatiques / Météorologiques} & Inondations, sécheresses, tempêtes, canicules & Inondations récurrentes de Douala (quartiers marécageux) ; Sécheresses dans l'Extrême-Nord \\ \hline
\cle{Biologiques} & Épidémies zoonotiques, invasions acridiennes, pullulations algales & Épidémies de choléra (liées à l'insalubrité + pluviométrie) ; Fièvre Ebola (réservoir animal) \\ \hline
\end{tabularx}
\end{center}
\end{propriete}

\begin{experience}[Modélisation d'une éruption limnique (démonstration)]
\textbf{Matériel} : grand bécher, eau, colorant alimentaire (eau profonde), huile végétale (couche superficielle moins dense), comprimé effervescent (source de \ce{CO2}) ou tuyau à bulles.
\textbf{Protocole} :
\begin{enumerate}
\item Remplir le bécher d'eau colorée (eau profonde riche en \ce{CO2} dissous).
\item Verser délicatement l'huile (eau de surface, moins dense, stratifiée).
\item Injecter du \ce{CO2} (bulles) au fond : le gaz se dissout sous pression.
\item Provoquer un brassage brutal (secousse ou retrait de l'huile) : libération soudaine du gaz.
\end{enumerate}
\textbf{Observation} : formation d'un nuage gazeux dense qui s'écoule par-dessus le bord du bécher (vers les zones basses).
\textbf{Interprétation} : la stratification empêche le dégazage progressif ; un déclencheur (glissement, séisme, refroidissement) brise la stratification $\Rightarrow$ dégazage catastrophique. Le \ce{CO2} (densité $\approx 1,5 \times$ air) s'accumule dans les vallées.
\legende{Modélisation de la stratification et du dégazage brutal type lac Nyos.}
\end{experience}

\begin{exoresolu}[Le lac Nyos : classification rigoureuse]
\begin{enumerate}
\item Définir « catastrophe ».
\item Classer l'éruption limnique du 21 août 1986 en justifiant.
\item Citer une autre catastrophe naturelle camerounaise.
\end{enumerate}
\tcblower
\textbf{1.} Voir définition ci-dessus.
\textbf{2.} \emph{Agent causal} : libération de \SIrange{100000}{300000}{tonnes} de \ce{CO2}. \emph{Origine} : magmatisme de la ligne volcanique du Cameroun (processus géologique naturel). \emph{Bilan} : \num{1746} morts, \num{3000} bovins, capacités locales dépassées. \textbf{Conclusion} : catastrophe naturelle de type \emph{géologique} (volcanisme).
\textbf{3.} Éruptions du Mont Cameroun (1999, 2000) ; Glissements de terrain de Bafoussam (2019) ; Inondations de Douala.
\end{exoresolu}

\courssub{Les catastrophes d'origine humaine}

\begin{definition}[Catastrophe d'origine humaine (technologique, sanitaire, environnementale)]
Catastrophe dont l'agent causal est une activité, une négligence ou une décision humaine : accident industriel, pollution massive, accident de transport de matières dangereuses, déforestation massive entraînant glissements/sécheresses, conflits armés, épidémies liées à l'insalubrité chronique.
\end{definition}

\begin{propriete}[Exemples de catastrophes d'origine humaine au Cameroun et dans le monde]
\begin{itemize}
\item \textbf{Industrielles / Technologiques} : Explosion du dépôt d'hydrocarbures de Douala-Bassa (1998, $>200$ morts) ; Explosions de camions-citernes (ex. Mbankolo 2016, >50 morts) ; Pollution industrielle du Wouri (métaux lourds, hydrocarbures).
\item \textbf{Environnementales (anthropiques)} : Déforestation massive (sud, est) $\Rightarrow$ érosion, perte de biodiversité, amplification des crues ; Gestion des déchets urbains (décharge de Nkolfoulou, Yaoundé ; HYSACAM) $\Rightarrow$ pollution sols/eaux, gaz à effet de serre.
\item \textbf{Sanitaires / Sociales} : Épidémies de choléra récurrentes (accès à l'eau/assainissement insuffisant) ; Accidents de la route (première cause de mortalité traumatique, >1000 morts/an).
\end{itemize}
\end{propriete}

\begin{attention}[Catastrophes mixtes : la règle au Bac]
La plupart des « catastrophes naturelles » sont \emph{mixtes}.
\begin{itemize}
\item \textbf{Inondations de Douala} : aléa climatique (pluies) + urbanisation anarchique en zone marécageuse + colmatage caniveaux (déchets plastiques) + déforestation amont.
\item \textbf{Glissements de Bafoussam (2019)} : aléa pluviométrique + habitat sur versants raides déboisés + absence de drainage.
\end{itemize}
\textbf{Consigne} : ne jamais écrire « catastrophe naturelle » sans préciser la part de vulnérabilité humaine si le sujet porte sur la gestion/prévention.
\end{attention}

\courssub{La gestion des catastrophes : prévision, prédiction, prévention}

\begin{methode}[Ne jamais confondre prévision, prédiction, prévention]
\begin{enumerate}
\item \textbf{Prévision} : estimation de la \emph{probabilité} qu'un phénomène se produise dans une zone et une période données, à partir de la surveillance (bulletins météo, surveillance sismique Mont Cameroun par l'IRGM). Question : « \emph{Où et avec quelle probabilité ?} ».
\item \textbf{Prédiction} : annonce, avec précision visée sur la \emph{date}, le \emph{lieu} et l'\emph{ampleur}, d'un événement \emph{imminent} fondée sur des signes précurseurs (essaims sismiques, déformation du sol, dégazage anormal). Question : « \emph{Quand exactement ?} ». Difficile pour les séismes.
\item \textbf{Prévention} : mesures prises \emph{en amont} pour réduire vulnérabilité et enjeux : normes de construction, zonage (non-construction lit majeur Sanaga), dégazage artificiel lac Nyos (tuyau 2001), sensibilisation, plans d'évacuation, reboisement. Question : « \emph{Comment limiter les dégâts ?} ».
\end{enumerate}
\textbf{Réflexe Bac} : associe chaque terme à un exemple camerounais concret.
\end{methode}

\begin{exoresolu}[Attribuer les actions : prévision, prédiction ou prévention]
Actions au Cameroun :
\begin{enumerate}
\item[a.] Tuyau de dégazage lac Nyos (2001).
\item[b.] Surveillance sismique Mont Cameroun (réseau stations).
\item[c.] Évacuation flancs Mont Cameroun (1999, éruption imminente annoncée).
\item[d.] Interdiction construction lit majeur Sanaga.
\item[e.] Bulletins alerte fortes pluies Douala (saison des pluies).
\item[f.] Formation élèves gestes qui sauvent / exercices évacuation.
\end{enumerate}
\tcblower
\textbf{a. Prévention} : réduit durablement l'aléa (stock de \ce{CO2}).
\textbf{b. Prévision} : estime probabilité éruption sans date précise.
\textbf{c. Prédiction (application)} : évacuation suite à annonce imminente (signes précurseurs).
\textbf{d. Prévention} : zonage réduit enjeu/vulnérabilité avant crue.
\textbf{e. Prévision} : modèles météo estiment probabilité pluies.
\textbf{f. Prévention} : éducation réduit vulnérabilité (volet éducatif).
\textbf{Bilan} : Prévention (a,d,f) ; Prévision (b,e) ; Prédiction (c). La prévention est le levier le plus efficace.
\end{exoresolu}

\courssub{Analyser un document sur une catastrophe (compétence Bac)}

\begin{methode}[Exploiter un tableau, une courbe, une carte de risques]
\begin{enumerate}
\item \textbf{Lire méthodiquement} : titre, légendes, unités, échelles, dates. Cite les données \emph{avec leurs unités}.
\item \textbf{Dégager les faits} : comparer grandeurs (morts, coût, fréquence), repérer tendances.
\item \textbf{Interpréter} : relier aux notions (aléa, vulnérabilité, enjeu). Rappel : $\text{Risque} = \text{Aléa} \times \text{Vulnérabilité} \times \text{Enjeu}$.
\item \textbf{Conclure} : répondre exactement à la question, mobiliser au moins une donnée chiffrée.
\end{enumerate}
\end{methode}

\begin{exoresolu}[Comparaison séismes Haïti (2010) vs Japon (2011)]
\begin{center}
\begin{tabularx}{\linewidth}{|l|X|X|}\hline
\rowcolor{popblueL} \textbf{Paramètre} & \textbf{Séisme A (Haïti, 2010)} & \textbf{Séisme B (Japon, 2011)} \\ \hline
Magnitude & 7,0 & 9,0 \\ \hline
Morts (séisme seul) & $\approx 220\,000$ & $\approx 16\,000$ (avec tsunami) \\ \hline
Normes parasismiques & Quasi absentes & Strictement appliquées \\ \hline
Habitat précaire & Très dense & Faible \\ \hline
\end{tabularx}
\end{center}
\begin{enumerate}
\item Montrer que l'ampleur ne dépend pas seulement de l'aléa.
\item Définir aléa, vulnérabilité ; expliquer l'écart.
\end{enumerate}
\tcblower
\textbf{1.} Magnitude B (9,0) $>$ A (7,0). Énergie $\propto 10^{1,5 \times \Delta M} = 10^{3} = 1000$ fois plus grande pour B. Pourtant morts A ($\approx 220\,000$) $\gg$ morts B ($\approx 16\,000$). \textbf{L'ampleur ne dépend pas que de l'aléa.}
\textbf{2.} \emph{Aléa} : probabilité phénomène dangereux. \emph{Vulnérabilité} : fragilité populations/bâtiments. Japon : normes strictes, habitat adapté $\Rightarrow$ faible vulnérabilité $\Rightarrow$ bâtiments résistent à aléa énorme. Haïti : absence normes, habitat précaire dense $\Rightarrow$ vulnérabilité extrême $\Rightarrow$ aléa modéré $\rightarrow$ catastrophe majeure. Relation $\text{Risque} = \text{Aléa} \times \text{Vulnérabilité} \times \text{Enjeu}$ confirmée. La prévention (construction, éducation) est clé.
\end{exoresolu}

\courssub{Les gestes qui sauvent et la conduite à tenir}

\begin{methode}[Conduite à tenir : Protéger — Alerter — Secourir]
Face à une victime/catastrophe, applique \textbf{dans l'ordre} :
\begin{enumerate}
\item \textbf{Protéger} : supprimer danger ou mettre victime/soi à l'abri (couper courant, éloigner bâtiment fissuré). \emph{Jamais de secouriste en danger.}
\item \textbf{Alerter} : appeler secours (Cameroun : \textbf{112} urgence européenne, \textbf{118} pompiers, \textbf{119} SAMU). Donner : nature, lieu précis, nombre victimes. \emph{Ne jamais raccrocher le premier.}
\item \textbf{Secourir} :
      \begin{itemize}
      \item Vérifier conscience puis respiration.
      \item Inconscient \textbf{qui respire} $\Rightarrow$ \cle{Position Latérale de Sécurité (PLS)} : victime sur le côté, tête en arrière (dégage voies respiratoires, évite étouffement langue/vomissements). Surveiller respiration.
      \item Inconscient \textbf{ne respirant pas} $\Rightarrow$ \cle{Massage cardiaque} : compressions centre poitrine, fréquence \SIrange{100}{120}{min^{-1}}, profondeur $\approx \SI{5}{cm}$.
      \item Hémorragie $\Rightarrow$ compression manuelle ou point de compression.
      \end{itemize}
\end{enumerate}
\textbf{Conduites spécifiques} : Séisme $\Rightarrow$ sous table solide, loin fenêtres, évacuer sans ascenseur ; Inondation $\Rightarrow$ couper eau/gaz/électricité, rejoindre point haut ; Incendie $\Rightarrow$ ramper sous fumées, portes fermées.
\end{methode}

\begin{exoresolu}[Accident axe Yaoundé–Douala]
Témoin : moto-taxi renversé, conducteur inconscient, poitrine se soulève (respire). Passant veut le redresser et donner de l'eau.
\begin{enumerate}
\item Expliquer pourquoi empêcher le passant.
\item Décrire les 3 actions dans l'ordre, justifiées.
\end{enumerate}
\tcblower
\textbf{1.} Victime inconsciente : boire $\Rightarrow$ \emph{fausse route} (asphyxie). Redresser brutalement $\Rightarrow$ risque aggraver \emph{lésion colonne vertébrale} (traumatisme cervical fréquent motocycliste) $\Rightarrow$ paralysie. \textbf{Règle} : ne jamais donner à boire à un inconscient, ne pas déplacer sauf danger imminent (incendie).
\textbf{2.}
\begin{itemize}
\item \emph{Protéger} : signaler (triangle, témoins amont), baliser zone. \emph{Justification} : éviter suraccident, sécuriser secouriste/victime.
\item \emph{Alerter} : 118/112. Préciser : accident circulation, axe Yaoundé–Douala, repère kilométrique, 1 victime inconsciente respirant. \emph{Justification} : prise en charge rapide conditionne survie.
\item \emph{Secourir} : Inconscient \textbf{respire} $\Rightarrow$ \cle{PLS}. Victime sur côté, tête en arrière. Surveiller respiration jusqu'aux secours. \emph{Justification} : maintient voies aériennes libres, évite étouffement.
\end{itemize}
\end{exoresolu}

\courssub{Élaborer un outil de sensibilisation (compétence savoir-faire)}

\begin{methode}[Concevoir affiche, sketch, dépliant efficace]
\begin{enumerate}
\item \textbf{Cible + Risque précis} : ex. riverains lac Nyos / éruption limnique ; élèves / séisme ; habitants Douala / inondation.
\item \textbf{UN message central} : une idée clé (ex. « Séisme : sous la table ! »).
\item \textbf{Contenu scientifique structuré} : nature risque, signes annonciateurs, consignes (avant/pendant/après), notions cours (aléa, vulnérabilité, prévention).
\item \textbf{Forme soignée} : langue simple (+ locale), phrases courtes, impératifs numérotés, pictogrammes, couleurs alerte (rouge/orange), \textbf{numéros d'urgence visibles}.
\item \textbf{Tester} : compris en $< \SI{30}{s}$ ? Consignes réalisables ? Test échantillon avant diffusion.
\end{enumerate}
\end{methode}

\begin{exoresolu}[Affiche riverains lac Nyos]
\begin{enumerate}
\item Cible et risque.
\item Message central + 3 consignes justifiées scientifiquement.
\end{enumerate}
\tcblower
\textbf{1.} Cible : populations réinstallées lac Nyos (Nord-Ouest), dont non-francophones (pictogrammes, langues locales). Risque : \cle{éruption limnique} (libération brutale \ce{CO2}, incolore, inodore, densité $\approx 1,5 \times$ air $\Rightarrow$ s'accumule vallées/points bas $\Rightarrow$ asphyxie).
\textbf{2.}
\begin{itemize}
\item \emph{Message} : « Grondement lac / odeur suspecte / animaux meurent : \textbf{FUIS VERS LES HAUTEURS !} »
\item \emph{Consigne 1} : Reconnaître signes (oiseaux/petits animaux morts au sol, eaux bouillonnantes/décolorées) $\rightarrow$ alerter autorités. \emph{Justification} : \ce{CO2} dense $\Rightarrow$ tue d'abord au ras du sol = indicateurs précoces.
\item \emph{Consigne 2} : Gagner \textbf{immédiatement les points hauts}, JAMAIS vallées/grottes. \emph{Justification} : nuage \ce{CO2} s'écoule/stagne en bas ; air respirable en altitude.
\item \emph{Consigne 3} : Respecter zones interdites habitat (zonage) ; ne pas approcher berges après séisme/glissement. \emph{Justification} : brassage eaux peut déclencher dégazage ; zonage réduit enjeu (prévention).
\end{itemize}
\textbf{Conclusion} : message unique impératif, pictogrammes (flèche colline, alerte rouge), consignes fondées physique \ce{CO2} = outil prévention efficace, adapté cible/risque local.
\end{exoresolu}

\courssub{Construire une argumentation de type Bac}

\begin{methode}[Rédiger un paragraphe argumenté (sujet ouvert)]
\begin{enumerate}
\item \textbf{Thèse} en 1 phrase (ex. « Catastrophes dites naturelles : part responsabilité humaine »).
\item \textbf{Définir termes-clés} (catastrophe, aléa, vulnérabilité, prévention) : 1ère attente correcteur.
\item \textbf{2--3 arguments} étayés par \emph{exemples précis et chiffrés} (Nyos 1986 : 1746 morts ; Inondations Douala ; Choléra/insalubrité).
\item \textbf{Conclusion} : réponse explicite + ouverture gestion (prévention, éducation).
\end{enumerate}
\textbf{Formules} : $\text{Risque} = \text{Aléa} \times \text{Vulnérabilité} \times \text{Enjeu}$ ; distinction naturel/humain/mixte ; triptyque prévision/prédiction/prévention.
\end{methode}

\begin{exoresolu}[Sujet synthèse : « Les catastrophes naturelles ne sont jamais totalement naturelles »]
Discuter en $\approx 15$ lignes avec 2 exemples argumentés.
\tcblower
\textbf{Thèse + Définitions.} Catastrophe = dommages majeurs $>$ capacités locales. Naturelle = agent causal processus Terre/vivant. Or $\text{Risque} = \text{Aléa} \times \text{Vulnérabilité} \times \text{Enjeu}$ : aléa naturel, mais vulnérabilité/enjeu façonnés par l'homme.

\textbf{Argument 1 (Nyos).} Aléa purement naturel (dégazage \ce{CO2} magmatique). Mais bilan 1746 morts par installation villages dans vallées traversées par nuage : absence zonage/surveillance $\Rightarrow$ vulnérabilité humaine transforme phénomène géologique en catastrophe. Dégazage artificiel (tuyau 2001) prouve action humaine réduit risque.

\textbf{Argument 2 (Douala / Bafoussam).} Inondations Douala / Glissements Bafoussam (2019, $>40$ morts) : aléa climatique (pluies) aggravé facteurs humains : urbanisation anarchique zones marécageuses/versants raides, déforestation (ruissellement), colmatage caniveaux (déchets plastiques). Catastrophe \emph{mixte}.

\textbf{Conclusion.} Affirmation fondée : aléa naturel, mais vulnérabilité choix humains (habitat, déforestation, absence prévention) détermine ampleur. D'où importance capitale prévention + éducation aux risques.
\end{exoresolu}

\begin{attention}[Les 5 erreurs qui coûtent des points au Bac]
\begin{enumerate}
\item \textbf{Confondre prévision, prédiction, prévention.} Prévision = probabilité ; Prédiction = annonce imminente (date/lieu/ampleur) ; Prévention = agir en amont (réduire vulnérabilité/enjeux). Attribuer tuyau Nyos à « prédiction » = erreur grave (c'est \emph{prévention}).
\item \textbf{Classer sans justifier l'agent causal.} « Nyos = naturelle » sans expliquer \ce{CO2} magmatique = moitié points perdus. Toute classification doit être argumentée.
\item \textbf{Oublier unités/chiffres / physique.} « 1746 morts Nyos » sans source/cohérence ; « \ce{CO2} plus lourd que l'air » sans déduire « fuir vers \emph{hauteurs} (pas vallées) » = copie affaiblie.
\item \textbf{Inverser chaîne secours.} Secourir avant Protéger/Alerter = faute grave. Ordre obligatoire : \textbf{Protéger $\rightarrow$ Alerter $\rightarrow$ Secourir}. Donner à boire à inconscient / déplacer blessé sans danger imminent = erreur conduite sanctionnée.
\item \textbf{Hors sujet synthèse.} Empiler définitions sans exemple chiffré, ou raconter catastrophe sans relier à $\text{Risque} = \text{Aléa} \times \text{Vulnérabilité} \times \text{Enjeu}$ et à la question. Correcteur attend : thèse, arguments illustrés, conclusion explicite.
\end{enumerate}
\end{attention}

\begin{savaistu}[Le lac Nyos : de la catastrophe à la prévention active]
Depuis 2001, un tuyau (puis deux autres en 2011) installé par une équipe internationale (IRGM, université de Savoie, PNUD) permet un \textbf{dégazage artificiel continu} : l'eau profonde riche en \ce{CO2} remonte par le tuyau, la pression diminue, le gaz se dégage en surface de façon contrôlée (geyser permanent). Cela a réduit la concentration de \ce{CO2} au fond de \SI{>80}{\%}, écartant le risque d'éruption limnique majeure. C'est un exemple unique au monde de \textbf{prévention technique d'une catastrophe géologique}. Le site est désormais surveillé en temps réel (capteurs pression, température, concentration \ce{CO2}) et un plan d'évacuation des populations est opérationnel.
\end{savaistu}

\begin{exercice}[Entraînement Bac - Classification et Gestion]
\begin{enumerate}
\item Une usine chimique rejette accidentellement un nuage toxique (chlore) sur un quartier d'habitation précaire : 15 morts, 200 blessés. Classer cette catastrophe et justifier en 3 lignes.
\item Lors d'une forte pluie, un quartier de Douala built sur un ancien marécage est inondé : 3 morts, 500 sinistrés. Préciser si c'est une catastrophe naturelle, humaine ou mixte. Justifier par les composantes du risque.
\item Citer les 3 numéros d'urgence au Cameroun et l'ordre strict de la chaîne de secours.
\item Une affiche de prévention contre le choléra doit cibler les élèves. Proposer un message central et 2 consignes impératives avec pictogrammes suggérés.
\end{enumerate}
\end{exercice}

\begin{corrige}[Exercice n°1]
\begin{enumerate}
\item \textbf{Catastrophe d'origine humaine (technologique/industrielle).} Agent causal : fuite de chlore (produit industriel manipulé par l'homme). Origine : défaillance technique/erreur humaine dans l'usine. Bilan lourd (15 morts) $>$ capacités locales $\Rightarrow$ catastrophe.
\item \textbf{Catastrophe mixte.} Aléa : pluie intense (naturel). Vulnérabilité : habitat précaire en zone marécageuse (humain), absence drainage (humain). Enjeu : population dense. L'aléa seul n'aurait pas causé ce bilan sans la vulnérabilité anthropique.
\item Numéros : \textbf{112} (urgence généraliste), \textbf{118} (pompiers), \textbf{119} (SAMU). Ordre : \textbf{1. Protéger} $\rightarrow$ \textbf{2. Alerter} $\rightarrow$ \textbf{3. Secourir}.
\item \emph{Message} : « Eau sale = Danger ! Lave-toi les mains, bois de l'eau traitée. » \emph{Consigne 1} : « Lave-toi les mains au savon après les toilettes et avant de manger » (pictogramme mains + savon + eau). \emph{Consigne 2} : « Bois uniquement de l'eau bouillie ou traitée (pastilles/javel) » (pictogramme bidon + feu ou pastille). Numéros urgence visibles.
\end{enumerate}
\end{corrige}

\newpage

\coursec{Exercices}

\courssub{Je vérifie mes connaissances}

\begin{exercice}[QCM : les notions clés]
Pour chaque question, choisir la (ou les) bonne(s) réponse(s).
\begin{enumerate}
\item Une \cle{catastrophe} est :
\begin{enumerate}
\item tout événement brutal provoquant des dégâts ;
\item un événement soudain causant des pertes humaines et matérielles importantes, dépassant les capacités locales de réaction ;
\item uniquement un phénomène naturel violent ;
\item toujours d'origine humaine.
\end{enumerate}
\item L'\cle{aléa} désigne :
\begin{enumerate}
\item la probabilité d'occurrence d'un phénomène dangereux ;
\item la vulnérabilité des populations ;
\item le nombre de victimes d'une catastrophe ;
\item le coût des dégâts matériels.
\end{enumerate}
\item La \cle{prévision} d'un cyclone consiste à :
\begin{enumerate}
\item annoncer la date exacte du phénomène ;
\item estimer, à partir de données météorologiques, la probabilité, la zone et l'intensité du phénomène à échéance courte ;
\item construire des digues ;
\item évacuer les populations.
\end{enumerate}
\item La catastrophe du lac Nyos (1986) est une catastrophe :
\begin{enumerate}
\item d'origine humaine ;
\item d'origine naturelle ;
\item mixte ;
\item nucléaire.
\end{enumerate}
\end{enumerate}
\end{exercice}

\begin{exercice}[Vrai ou faux, à justifier]
Répondre par VRAI ou FAUX, puis justifier chaque réponse en une ou deux phrases.
\begin{enumerate}
\item Un accident et une catastrophe se distinguent uniquement par leur origine.
\item La prédiction d'un séisme consiste à annoncer la date, l'heure et le lieu exacts du tremblement de terre.
\item La prévention d'une catastrophe passe par la réduction de la vulnérabilité des populations.
\item Une épidémie de choléra dans un quartier sans eau potable est une catastrophe dont l'origine peut être à la fois naturelle et humaine.
\item Les gestes qui sauvent s'organisent selon la séquence : Alerter -- Protéger -- Secourir.
\end{enumerate}
\end{exercice}

\begin{exercice}[Définitions]
Définir précisément, en une phrase chacune, les termes suivants : \cle{catastrophe} ; \cle{risque} ; \cle{vulnérabilité} ; \cle{prévention}. Puis, à l'aide d'un exemple camerounais de votre choix, montrer comment l'aléa, la vulnérabilité et les enjeux se combinent pour produire un risque.
\end{exercice}

\begin{exercice}[Quelques ordres de grandeur]
Le 21 août 1986, le lac Nyos a libéré brutalement un volume estimé à environ \SI{1,6}{\kilo\metre\cubed} de \ce{CO2}. Le nuage gazeux, plus dense que l'air (densité du \ce{CO2} : environ \num{1,5} fois celle de l'air), a dévalé les vallées voisines et causé la mort d'environ \num{1746} personnes et \num{3500} têtes de bétail dans un rayon d'environ \SI{25}{km}.
\begin{enumerate}
\item Rappeler la formule chimique du dioxyde de carbone et expliquer pourquoi le nuage est resté plaqué au sol.
\item Sachant que la masse molaire du \ce{CO2} est $M = \SI{44}{g.mol^{-1}}$ et que le volume molaire dans les conditions du lac est voisin de \SI{24}{L.mol^{-1}}, estimer la masse de \ce{CO2} dégagée (on prendra une densité du gaz libéré de l'ordre de \SI{2}{kg.m^{-3}} pour simplifier). Exprimer le résultat en tonnes.
\item Citer deux conséquences sanitaires directes de l'inhalation d'un air très riche en \ce{CO2}.
\end{enumerate}
\end{exercice}

\courssub{Je m'entraîne}

\begin{exercice}[Classer des catastrophes]
Voici une liste de dix catastrophes :
\begin{enumerate}
\item émission limnique du lac Nyos (Cameroun, 1986) ;
\item accident nucléaire de Fukushima (Japon, 2011) ;
\item épidémie de choléra dans un quartier de Yaoundé ;
\item séisme de Haïti (2010) ;
\item marée noire provoquée par le naufrage d'un pétrolier ;
\item feu de brousse dans la région de l'Adamaoua ;
\item conflit armé provoquant des déplacements massifs de populations ;
\item inondations récurrentes de certains quartiers de Douala en saison des pluies ;
\item éruption du Mont Cameroun (1999) ;
\item accident ferroviaire d'Eséka (Cameroun, 2016).
\end{enumerate}
\begin{enumerate}
\item Classer chaque catastrophe dans l'une des trois catégories : d'origine \cle{naturelle}, d'origine \cle{humaine} (technologique ou sociale), ou \cle{mixte}.
\item Justifier chaque classement en une phrase, en précisant la part du phénomène naturel et celle de l'activité humaine.
\item Pour la catastrophe de Fukushima, expliquer pourquoi elle est qualifiée de catastrophe « en cascade » ou « en chaîne ».
\end{enumerate}
\end{exercice}

\begin{exercice}[Aléa, vulnérabilité, risque : le cas des inondations de Douala]
Chaque année, les pluies diluviennes provoquent des inondations dans plusieurs quartiers de Douala construits dans des zones basses et marécageuses, où les caniveaux sont souvent obstrués par les déchets.
\begin{enumerate}
\item Identifier dans cette situation : l'aléa, les facteurs de vulnérabilité, et les enjeux exposés.
\item Proposer trois mesures de prévention concrètes, en précisant pour chacune si elle agit sur l'aléa ou sur la vulnérabilité.
\item Expliquer pourquoi deux quartiers soumis à la même pluie peuvent subir des dégâts très différents.
\end{enumerate}
\end{exercice}

\begin{exercice}[Prévision, prédiction, prévention : à chacun sa définition]
Trois affirmations sont rapportées dans la presse :
\begin{itemize}
\item Affirmation A : « Les sismologues affirment qu'un grand séisme frappera Douala le 15 mars prochain à 14 h 30. »
\item Affirmation B : « Météo Cameroun annonce de fortes pluies sur le littoral dans les prochaines 48 heures, avec un risque d'inondation. »
\item Affirmation C : « La mairie impose de nouvelles normes de construction parasismique et organise un exercice d'évacuation dans les écoles. »
\end{itemize}
\begin{enumerate}
\item Associer chaque affirmation à l'un des trois termes : prévision, prédiction, prévention. Justifier.
\item Expliquer pourquoi l'affirmation A doit être accueillie avec la plus grande méfiance, en rappelant l'état actuel des connaissances scientifiques sur les séismes.
\item Citer un exemple de catastrophe pour laquelle la prévision à court terme est fiable, et un exemple pour laquelle seule la prévention est possible.
\end{enumerate}
\end{exercice}

\begin{exercice}[Rédiger une consigne de sécurité]
Tu es délégué de classe en Terminale C au lycée de ta localité. Le proviseur te charge de rédiger l'affiche de sécurité à apposer dans chaque salle, intitulée « Conduite à tenir en cas de séisme ou d'incendie ».
\begin{enumerate}
\item Rédiger cette affiche en respectant la séquence des \cle{gestes qui sauvent} : Protéger -- Alerter -- Secourir (au moins deux consignes par étape, rédigées à l'impératif).
\item Ajouter deux consignes spécifiques au séisme (comportement pendant la secousse) et deux consignes spécifiques à l'incendie.
\item Expliquer pourquoi il est interdit d'utiliser les ascenseurs et de revenir dans le bâtiment avant l'autorisation des secours.
\end{enumerate}
\end{exercice}

\courssub{Je me prépare au Bac}

\begin{exercice}[Type Bac : l'explosion limnique du lac Nyos et sa prévention]
Le lac Nyos (Nord-Ouest du Cameroun) occupe un ancien cratère volcanique. Le 21 août 1986, il a brutalement dégagé un énorme volume de \ce{CO2}, tuant environ \num{1746} personnes. Le document ci-dessous présente la concentration en \ce{CO2} dissous dans l'eau du lac, mesurée à différentes profondeurs, en 1986 (avant toute intervention) puis après l'installation en 2001 d'un tuyau de dégazage reliant les eaux profondes à la surface.

\begin{center}
\begin{tabularx}{\linewidth}{|c|X|X|}\hline
\rowcolor{popblueL} Profondeur (m) & \ce{CO2} dissous en 1986 (mmol.L$^{-1}$) & \ce{CO2} dissous après dégazage (mmol.L$^{-1}$) \\ \hline
0 (surface) & 5 & 4 \\ \hline
$-50$ & 30 & 20 \\ \hline
$-100$ & 120 & 45 \\ \hline
$-150$ & 300 & 80 \\ \hline
$-200$ (fond) & 550 & 110 \\ \hline
\end{tabularx}
\end{center}

\begin{enumerate}
\item Décrire l'évolution de la concentration en \ce{CO2} avec la profondeur en 1986. Calculer le rapport entre la concentration au fond et celle en surface.
\item Sachant que le \ce{CO2} provient d'une source magmatique sous-jacente et s'accumule dans les eaux profondes froides et sous pression, proposer une hypothèse expliquant le dégazage brutal de 1986 (on pourra raisonner par analogie avec une bouteille d'eau gazeuse que l'on débouche).
\item Comparer les deux profils de concentration et montrer que le tuyau de dégazage constitue une mesure de \cle{prévention}. Préciser si cette mesure agit sur l'aléa ou sur la vulnérabilité.
\item Le dégazage contrôlé ne supprime pas tout danger. Proposer deux mesures complémentaires de gestion du risque pour les populations riveraines (surveillance, aménagement du territoire, éducation des populations).
\item Rédiger une conclusion de quatre à six lignes montrant comment l'étude scientifique du lac Nyos illustre la démarche de gestion d'une catastrophe naturelle : compréhension du mécanisme, surveillance, prévention.
\end{enumerate}
\end{exercice}

\begin{exercice}[Type Bac : évolution mondiale du nombre de catastrophes]
Le tableau ci-dessous donne le nombre moyen de catastrophes (naturelles et technologiques) recensées chaque année dans le monde, par décennie.

\begin{center}
\begin{tabularx}{\linewidth}{|l|X|}\hline
\rowcolor{popblueL} Décennie & Nombre moyen de catastrophes recensées par an \\ \hline
1970--1979 & environ 80 \\ \hline
1980--1989 & environ 160 \\ \hline
1990--1999 & environ 260 \\ \hline
2000--2009 & environ 350 \\ \hline
2010--2019 & environ 330 \\ \hline
\end{tabularx}
\end{center}

\begin{enumerate}
\item Tracer, sur une feuille de papier millimétré, la courbe du nombre de catastrophes recensées en fonction du temps (échelles : \SI{2}{cm} pour 10 ans en abscisse ; \SI{2}{cm} pour 100 catastrophes en ordonnée). Dégager la tendance générale.
\item Calculer le pourcentage d'augmentation du nombre de catastrophes recensées entre la décennie 1970--1979 et la décennie 2000--2009.
\item Trois hypothèses sont proposées pour expliquer cette augmentation : (a) l'amélioration des moyens d'observation et de recensement ; (b) l'urbanisation croissante des zones exposées (littoraux, bassins fluviaux) ; (c) le changement climatique qui accroît la fréquence de certains événements extrêmes. Pour chacune, indiquer si elle porte sur l'aléa, sur la vulnérabilité, ou sur la perception du phénomène, et justifier.
\item Expliquer pourquoi la légère baisse observée pour la décennie 2010--2019 ne doit pas conduire à relâcher les efforts de prévention.
\item En t'appuyant sur deux exemples camerounais, dégager l'importance de la prévention dans la réduction des effets des catastrophes.
\end{enumerate}
\end{exercice}

\begin{exercice}[Type Bac : question de synthèse]
Sujet de synthèse : « Montrer, à partir d'exemples camerounais et mondiaux, que la gravité d'une catastrophe dépend autant de la vulnérabilité des populations que de l'aléa lui-même. »

Pour traiter ce sujet en une vingtaine de lignes, tu pourras t'appuyer sur les éléments suivants :
\begin{itemize}
\item le séisme de Haïti (2010, magnitude 7,0 : plus de \num{200000} morts) comparé à des séismes de magnitude comparable au Japon (bilan humain très inférieur) ;
\item l'émission limnique du lac Nyos (1986) : un aléa rare mais des villages situés dans les vallées de déversement du nuage ;
\item les inondations de Douala : même pluie, dégâts inégaux selon les quartiers ;
\item les éruptions du Mont Cameroun : des coulées de lave fréquentes mais peu meurtrières grâce à la faible densité de population sur les zones de coulées et à la surveillance.
\end{itemize}

\begin{enumerate}
\item Définir les notions d'aléa, de vulnérabilité et de risque, et rappeler la relation qui les lie.
\item Construire ta démonstration en comparant au moins deux situations où un aléa semblable produit des bilans très différents.
\item Conclure en dégageant les implications pratiques pour la gestion des catastrophes au Cameroun (aménagement du territoire, éducation des populations, plans d'urgence).
\end{enumerate}
\end{exercice}

\newpage

\coursec{Corrigés des exercices}

\begin{corrige}[Exercice 1 — QCM : les notions clés]
\textbf{Énoncé :} Se reporter à l'énoncé original de l'exercice.
\tcblower
\begin{enumerate}
\item \textbf{Bonne réponse : b.} Une catastrophe est un événement soudain et brutal causant des pertes humaines et matérielles importantes, qui dépasse les capacités locales de réaction et d'adaptation. La réponse a est insuffisante (elle ne distingue pas l'accident de la catastrophe : c'est l'ampleur des dégâts et le dépassement des capacités locales qui comptent) ; les réponses c et d sont fausses car une catastrophe peut être d'origine naturelle, humaine ou mixte.
\item \textbf{Bonne réponse : a.} L'aléa est la probabilité d'occurrence d'un phénomène dangereux d'une intensité donnée, en un lieu et pendant une période donnés. La vulnérabilité (b) est une notion distincte ; le nombre de victimes (c) et le coût des dégâts (d) sont des conséquences, non l'aléa.
\item \textbf{Bonne réponse : b.} La prévision est une estimation probabiliste (probabilité, zone, intensité, échéance courte) fondée sur des données scientifiques (satellites, modèles météorologiques). Annoncer une date exacte (a) relèverait de la prédiction, scientifiquement impossible ; construire des digues (c) et évacuer (d) relèvent de la prévention et de la gestion de crise.
\item \textbf{Bonne réponse : b.} L'émission limnique du lac Nyos est une catastrophe d'origine naturelle : le \ce{CO2} est d'origine magmatique et le dégazage brutal est un phénomène limnologique naturel, sans cause humaine directe. La vulnérabilité des villages riverains est humaine, mais l'origine du phénomène est naturelle.
\end{enumerate}
\end{corrige}

\begin{corrige}[Exercice 2 — Vrai ou faux, à justifier]
\textbf{Énoncé :} Se reporter à l'énoncé original de l'exercice.
\tcblower
\begin{enumerate}
\item \textbf{FAUX.} Un accident et une catastrophe se distinguent par l'ampleur des dégâts : la catastrophe provoque des pertes massives qui dépassent les capacités locales de réaction, alors que l'accident reste gérable localement. L'origine (naturelle ou humaine) ne permet pas de les distinguer.
\item \textbf{FAUX.} La prédiction (date, heure, lieu exacts) d'un séisme est scientifiquement impossible à l'heure actuelle : aucune méthode ne permet de déterminer précisément le déclenchement d'un tremblement de terre. Seule la prévision probabiliste (zones à risque, probabilités à long terme) et la prévention sont réalisables.
\item \textbf{VRAI.} Le risque résulte de la combinaison de l'aléa, de la vulnérabilité et des enjeux ; comme on ne peut le plus souvent pas supprimer l'aléa naturel, la prévention consiste essentiellement à réduire la vulnérabilité (normes de construction, éducation, plans d'évacuation) et à limiter l'exposition des enjeux.
\item \textbf{VRAI.} Le choléra est causé par un agent biologique (\emph{Vibrio cholerae}), donc d'origine naturelle ; mais sa diffusion épidémique est favorisée par des facteurs humains : absence d'eau potable, insuffisance d'assainissement, densité de population. C'est donc une catastrophe mixte.
\item \textbf{FAUX.} La séquence correcte des gestes qui sauvent est : \cle{Protéger} (se mettre à l'abri, éliminer le danger), puis \cle{Alerter} (prévenir les secours), puis \cle{Secourir} (premiers gestes). Alerter avant de se protéger exposerait le sauveteur au danger.
\end{enumerate}
\end{corrige}

\begin{corrige}[Exercice 3 — Définitions]
\textbf{Énoncé :} Se reporter à l'énoncé original de l'exercice.
\tcblower
\begin{itemize}
\item \textbf{Catastrophe} : événement soudain et brutal, d'origine naturelle ou humaine, provoquant des pertes humaines, matérielles et environnementales importantes qui dépassent les capacités locales de réaction.
\item \textbf{Risque} : possibilité de dommages résultant de la combinaison d'un aléa, de la vulnérabilité des populations et des enjeux exposés ; on écrit souvent : Risque $=$ Aléa $\times$ Vulnérabilité $\times$ Enjeux.
\item \textbf{Vulnérabilité} : degré de fragilité d'une population, d'une société ou d'un territoire face à un aléa, lié à la qualité des constructions, au niveau de préparation, à la pauvreté, à l'organisation des secours.
\item \textbf{Prévention} : ensemble des mesures prises avant la catastrophe pour réduire le risque, en agissant sur l'aléa quand c'est possible (dégazage d'un lac) et surtout sur la vulnérabilité et l'exposition des enjeux (normes, aménagement, éducation).
\end{itemize}

\textbf{Exemple camerounais.} À Douala, les pluies diluviennes de saison constituent l'\cle{aléa} (probabilité élevée de fortes précipitations) ; les constructions en zones basses et marécageuses, les caniveaux obstrués par les déchets et l'absence de plans d'évacuation constituent la \cle{vulnérabilité} ; les habitants, leurs habitations et les infrastructures sont les \cle{enjeux} exposés. La combinaison des trois produit le \cle{risque} d'inondations meurtrières et destructrices, qui se réalise chaque année dans certains quartiers.
\end{corrige}

\begin{corrige}[Exercice 4 — Quelques ordres de grandeur]
\textbf{Énoncé :} Se reporter à l'énoncé original de l'exercice.
\tcblower
\begin{enumerate}
\item La formule chimique du dioxyde de carbone est \ce{CO2}. Sa densité est environ \num{1,5} fois celle de l'air : le gaz libéré, plus lourd que l'air, ne s'élève pas mais s'écoule au ras du sol en suivant les vallées, comme un liquide. C'est pourquoi le nuage est resté plaqué au sol et a asphyxié les populations des vallées voisines.
\item Conversion du volume : $V = \SI{1,6}{km^3} = \num{1,6e9}~\si{m^3}$. Avec une densité du gaz $d = \SI{2}{kg.m^{-3}}$ :
\[ m = d \times V = 2 \times \num{1,6e9} = \num{3,2e9}~\si{kg} = \num{3,2e6}~\si{t}. \]
La masse de \ce{CO2} dégagée est donc de l'ordre de \textbf{3,2 millions de tonnes}. (Vérification par les moles : $n = \dfrac{V}{V_m} = \dfrac{\num{1,6e12}~\si{L}}{24~\si{L.mol^{-1}}} \approx \num{6,7e10}~\si{mol}$, soit $m = n \times M \approx \num{6,7e10} \times 44 \approx \num{2,9e12}~\si{g} \approx \num{2,9e6}~\si{t}$ : même ordre de grandeur, le calcul est cohérent.)
\item Deux conséquences sanitaires directes de l'inhalation d'un air très riche en \ce{CO2} :
\begin{itemize}
\item l'\textbf{asphyxie} par privation d'oxygène (le \ce{CO2} chasse l'\ce{O2} de l'air respirable), entraînant la perte de conscience puis la mort en quelques minutes ;
\item l'\textbf{acidose respiratoire} : le \ce{CO2} dissous dans le sang abaisse le pH sanguin, perturbant le fonctionnement du cœur et du système nerveux.
\end{itemize}
\end{enumerate}
\end{corrige}

\begin{corrige}[Exercice 5 — Classer des catastrophes]
\textbf{Énoncé :} Se reporter à l'énoncé original de l'exercice.
\tcblower
\begin{enumerate}
\item \textbf{Classement et justifications :}
\begin{enumerate}
\item \emph{Lac Nyos (1986)} : \textbf{naturelle} — dégazage limnique d'origine magmatique, sans cause humaine directe.
\item \emph{Fukushima (2011)} : \textbf{mixte} — le séisme et le tsunami sont naturels, mais la catastrophe nucléaire résulte d'une installation technologique humaine vulnérable.
\item \emph{Choléra à Yaoundé} : \textbf{mixte} — agent biologique naturel, diffusion favorisée par le défaut d'eau potable et d'assainissement (facteurs humains).
\item \emph{Séisme de Haïti (2010)} : \textbf{naturelle} — rupture de faille tectonique ; l'ampleur du bilan tient à la vulnérabilité humaine, non à l'origine du phénomène.
\item \emph{Marée noire} : \textbf{humaine (technologique)} — naufrage d'un pétrolier, accident lié à une activité industrielle et au transport.
\item \emph{Feu de brousse (Adamaoua)} : \textbf{humaine} — les feux de brousse sont le plus souvent allumés par l'homme (défrichage, culture sur brûlis), même si la sécheresse les favorise.
\item \emph{Conflit armé} : \textbf{humaine (sociale)} — catastrophe provoquée entièrement par des décisions et des affrontements humains.
\item \emph{Inondations de Douala} : \textbf{mixte} — aléa pluvieux naturel, aggravé par l'urbanisation des zones marécageuses et l'obstruction des caniveaux par les déchets.
\item \emph{Éruption du Mont Cameroun (1999)} : \textbf{naturelle} — phénomène volcanique indépendant de l'activité humaine.
\item \emph{Accident ferroviaire d'Eséka (2016)} : \textbf{humaine (technologique)} — déraillement lié au transport et à la gestion des infrastructures.
\end{enumerate}
\item \textbf{Fukushima, catastrophe « en cascade » :} un premier événement naturel (séisme de magnitude 9,0) déclenche un second phénomène naturel (tsunami), qui provoque à son tour une catastrophe technologique (fusion des cœurs des réacteurs, rejets radioactifs). Chaque événement en engendre un autre : c'est l'enchaînement des causes, de l'aléa naturel vers la défaillance technologique, qui justifie le terme de catastrophe « en chaîne ».
\end{enumerate}
\end{corrige}

\begin{corrige}[Exercice 6 — Aléa, vulnérabilité, risque : les inondations de Douala]
\textbf{Énoncé :} Se reporter à l'énoncé original de l'exercice.
\tcblower
\begin{enumerate}
\item \textbf{Aléa} : les pluies diluviennes annuelles (probabilité élevée de précipitations intenses sur le littoral). \textbf{Facteurs de vulnérabilité} : constructions en zones basses et marécageuses, caniveaux obstrués par les déchets (écoulement impossible), absence de plans d'évacuation, pauvreté des habitants. \textbf{Enjeux exposés} : les habitants (vies humaines), les habitations, les biens, les routes et les activités économiques.
\item Trois mesures de prévention :
\begin{itemize}
\item curage régulier des caniveaux et gestion des déchets : agit sur la \textbf{vulnérabilité} (elle ne supprime pas la pluie, mais permet l'écoulement) ;
\item interdiction de construire dans les zones inondables et relogement : agit sur la \textbf{vulnérabilité} et l'exposition des enjeux ;
\item reboisement des versants et aménagement de bassins de rétention : agit partiellement sur l'\textbf{aléa} (réduction du ruissellement) — c'est l'une des rares façons d'agir sur l'aléa hydrologique.
\end{itemize}
\item Le risque dépend de la combinaison aléa $\times$ vulnérabilité $\times$ enjeux. À aléa identique (même pluie), un quartier bien drainé, construit sur un terrain élevé avec des habitants sensibilisés subira peu de dégâts, tandis qu'un quartier bas, aux caniveaux bouchés et densément peuplé sera dévasté. C'est la vulnérabilité, et non le seul aléa, qui détermine l'ampleur de la catastrophe.
\end{enumerate}
\end{corrige}

\begin{corrige}[Exercice 7 — Prévision, prédiction, prévention]
\textbf{Énoncé :} Se reporter à l'énoncé original de l'exercice.
\tcblower
\begin{enumerate}
\item \textbf{Affirmation A = prédiction} : elle annonce une date, une heure et un lieu exacts. \textbf{Affirmation B = prévision} : estimation probabiliste à échéance courte (48 h) fondée sur des données météorologiques, avec mention d'un risque. \textbf{Affirmation C = prévention} : mesures prises en amont (normes parasismiques, exercices d'évacuation) pour réduire la vulnérabilité.
\item L'affirmation A doit être accueillie avec la plus grande méfiance car la prédiction précise des séismes est \textbf{scientifiquement impossible} : on ne connaît aucun précurseur fiable permettant d'annoncer la date, l'heure et le lieu exacts d'un tremblement de terre. La sismologie ne fournit que des prévisions probabilistes (cartes d'aléa, probabilités à long terme). Une telle annonce relève de la rumeur et peut provoquer des mouvements de panique dangereux.
\item Prévision à court terme fiable : les \textbf{cyclones} et les \textbf{fortes pluies} (satellites météorologiques, trajectoires prévisibles plusieurs jours à l'avance). Cas où seule la prévention est possible : les \textbf{séismes} (normes parasismiques, éducation, plans d'urgence), car aucune prévision à court terme n'existe.
\end{enumerate}
\end{corrige}

\begin{corrige}[Exercice 8 — Rédiger une consigne de sécurité]
\textbf{Énoncé :} Se reporter à l'énoncé original de l'exercice.
\tcblower
\begin{enumerate}
\item \textbf{Affiche « Conduite à tenir en cas de séisme ou d'incendie » :}
\begin{itemize}
\item \textbf{PROTÉGER} : Garde ton calme, ne cours pas. Mets-toi à l'abri sous une table ou contre un mur porteur, loin des fenêtres. Coupe, si possible, le gaz et l'électricité.
\item \textbf{ALERTER} : Dès que tu es en sécurité, alerte les secours (pompiers : 118) et signale précisément le lieu et la nature du sinistre. Préviens le proviseur ou le surveillant.
\item \textbf{SECOURIR} : Évacue en silence par l'escalier vers le point de rassemblement. Aide les blessés sans les déplacer s'ils sont gravement atteints, et obéis aux consignes des secours.
\end{itemize}
\item \textbf{Spécifique séisme} : pendant la secousse, reste à l'abri sous un meuble solide et protège ta tête avec tes bras ; ne t'approche jamais des fenêtres ni des façades après la secousse (chutes d'objets, répliques possibles). \textbf{Spécifique incendie} : baisse-toi sous les fumées (l'air respirable est près du sol) et couvre ton nez avec un tissu humide ; n'ouvre pas une porte chaude au toucher et n'utilise jamais d'eau sur un feu d'origine électrique.
\item Les ascenseurs sont interdits car une coupure d'électricité ou une déformation du bâtiment peut y bloquer les occupants, et la cage d'ascenseur agit comme une cheminée qui propage les fumées. Il est interdit de revenir dans le bâtiment car des répliques (séisme) ou des structures fragilisées (incendie) peuvent provoquer l'effondrement : seuls les secours, après expertise, autorisent le retour.
\end{enumerate}
\end{corrige}

\begin{corrige}[Exercice 9 — Type Bac : l'explosion limnique du lac Nyos et sa prévention]
\textbf{Énoncé :} Se reporter à l'énoncé original de l'exercice.
\tcblower
\begin{enumerate}
\item En 1986, la concentration en \ce{CO2} dissous \textbf{augmente fortement avec la profondeur} : de \SI{5}{mmol.L^{-1}} en surface à \SI{550}{mmol.L^{-1}} au fond. Rapport fond/surface :
\[ \frac{550}{5} = 110. \]
L'eau du fond est donc \textbf{110 fois plus concentrée} en \ce{CO2} que l'eau de surface : le lac est chargé en gaz dans ses couches profondes.
\item \textbf{Hypothèse :} le \ce{CO2} magmatique s'accumule dans les eaux profondes, froides et soumises à une forte pression, qui le maintiennent dissous (comme le gaz dans une bouteille d'eau gazeuse fermée). Si un événement déclencheur (glissement de terrain, brassement des couches d'eau, refroidissement de surface) fait remonter les eaux profondes, la pression diminue brutalement : le gaz se dégage massivement (comme à l'ouverture de la bouteille), ce qui entraîne encore plus d'eau vers la surface — le dégazage s'auto-entretient et devient explosif.
\item Après dégazage, les concentrations profondes ont fortement chuté : à $-200$ m, on passe de 550 à \SI{110}{mmol.L^{-1}} (division par 5) ; à $-100$ m, de 120 à \SI{45}{mmol.L^{-1}}. Le tuyau ramène en continu les eaux profondes chargées en gaz vers la surface, où le \ce{CO2} se dégage de façon contrôlée : le lac ne peut plus accumuler une charge explosive. C'est une mesure de \textbf{prévention} qui agit sur l'\textbf{aléa} (elle réduit la probabilité et l'intensité du phénomène dangereux lui-même), ce qui est rare et remarquable.
\item Deux mesures complémentaires : (a) une \textbf{surveillance permanente} du lac (capteurs de \ce{CO2}, station de mesure, système d'alerte des populations) ; (b) l'\textbf{aménagement du territoire} : interdiction d'habiter les vallées de déversement du nuage, relogement des villages, et \textbf{éducation des populations} (exercices d'évacuation, information sur les signes avant-coureurs).
\item \textbf{Conclusion (modèle) :} La catastrophe du lac Nyos illustre la démarche complète de gestion d'une catastrophe naturelle. L'étude scientifique a d'abord permis de comprendre le mécanisme : accumulation de \ce{CO2} magmatique dans les eaux profondes sous pression, puis dégazage brutal. Cette compréhension a rendu possible une prévention agissant directement sur l'aléa : le dégazage contrôlé par tuyau, qui a divisé par cinq la charge en gaz du fond. Complétée par la surveillance du lac et l'éducation des populations, elle montre que la science transforme une menace mortelle en risque maîtrisé.
\end{enumerate}
\textbf{Barème indicatif (sur 20) :} question 1 : 4 pts (description 2, calcul du rapport 2) ; question 2 : 4 pts ; question 3 : 4 pts ; question 4 : 4 pts ; question 5 : 4 pts.

\textbf{Points de vigilance :} exiger le calcul explicite du rapport (110) avec les unités ; l'analogie avec la bouteille gazeuse doit citer le rôle de la pression ; bien distinguer aléa et vulnérabilité à la question 3 (le tuyau agit sur l'aléa, non sur la vulnérabilité) ; la conclusion doit mobiliser les trois verbes : comprendre, surveiller, prévenir.
\end{corrige}

\begin{corrige}[Exercice 10 — Type Bac : évolution mondiale du nombre de catastrophes]
\textbf{Énoncé :} Se reporter à l'énoncé original de l'exercice.
\tcblower
\begin{enumerate}
\item \textbf{Construction de la courbe :} on place en abscisse les décennies (2 cm pour 10 ans) et en ordonnée le nombre de catastrophes (2 cm pour 100). Points à placer : (1975 ; 80), (1985 ; 160), (1995 ; 260), (2005 ; 350), (2015 ; 330). La courbe monte fortement de 1970 à 2009 (multiplication par plus de 4), puis marque une légère baisse sur 2010--2019. \textbf{Tendance générale : forte augmentation} du nombre de catastrophes recensées sur cinquante ans.
\item Pourcentage d'augmentation entre 1970--1979 (80) et 2000--2009 (350) :
\[ \frac{350 - 80}{80} \times 100 = \frac{270}{80} \times 100 = 337{,}5~\%. \]
Le nombre de catastrophes recensées a augmenté de \textbf{337,5 \%} (il a été multiplié par $350/80 = 4{,}375$).
\item Analyse des hypothèses :
\begin{itemize}
\item (a) Amélioration de l'observation et du recensement : porte sur la \textbf{perception} du phénomène — on déclare aujourd'hui des catastrophes qui passaient inaperçues (satellites, médias, registres), sans que le risque réel ait forcément changé.
\item (b) Urbanisation des zones exposées : porte sur la \textbf{vulnérabilité} et les enjeux — plus de populations et de biens dans les zones dangereuses, donc plus de catastrophes pour un même aléa.
\item (c) Changement climatique : porte sur l'\textbf{aléa} — il accroît la fréquence et l'intensité de certains événements extrêmes (sécheresses, inondations, cyclones).
\end{itemize}
\item La baisse de 2010--2019 (de 350 à 330, soit environ $-6~\%$) est faible et peut n'être qu'une fluctuation ; elle ne remet pas en cause la tendance lourde à la hausse. De plus, l'urbanisation des zones exposées et le changement climatique se poursuivent : relâcher la prévention reviendrait à accroître la vulnérabilité face à un aléa qui reste élevé.
\item Exemples camerounais : (a) au \textbf{lac Nyos}, le dégazage artificiel et la surveillance ont transformé un piège mortel en risque surveillé ; (b) autour du \textbf{Mont Cameroun}, la surveillance volcanique et la faible occupation des zones de coulées font que les éruptions de 1999 et 2000 n'ont fait aucune victime. La prévention réduit donc concrètement les effets des catastrophes, même quand l'aléa demeure.
\end{enumerate}
\textbf{Barème indicatif (sur 20) :} question 1 : 5 pts (axes et échelles 2, points 2, tendance 1) ; question 2 : 3 pts ; question 3 : 6 pts (2 par hypothèse) ; question 4 : 3 pts ; question 5 : 3 pts.

\textbf{Points de vigilance :} le pourcentage se calcule par rapport à la valeur initiale (80), non à la valeur finale ; ne pas confondre « nombre de catastrophes recensées » et « nombre de catastrophes réelles » (biais de recensement) ; à la question 3, chaque hypothèse doit être explicitement rattachée à aléa, vulnérabilité ou perception.
\end{corrige}

\begin{corrige}[Exercice 11 — Type Bac : question de synthèse]
\textbf{Énoncé :} Se reporter à l'énoncé original de l'exercice.
\tcblower
\textbf{Corrigé type (environ vingt lignes) :}

\begin{enumerate}
\item \textbf{Définitions.} L'\cle{aléa} est la probabilité d'occurrence d'un phénomène dangereux d'intensité donnée en un lieu et une période donnés. La \cle{vulnérabilité} est le degré de fragilité des populations et des sociétés face à cet aléa (qualité des constructions, pauvreté, préparation, organisation des secours). Le \cle{risque} résulte de leur combinaison avec les enjeux exposés : Risque $=$ Aléa $\times$ Vulnérabilité $\times$ Enjeux. La gravité d'une catastrophe n'est donc pas inscrite dans le seul phénomène naturel.
\item \textbf{Démonstration.} Le séisme de Haïti (2010, magnitude 7,0) a tué plus de \num{200000} personnes, alors que des séismes de magnitude comparable au Japon font des bilans humains très inférieurs : à aléa semblable, c'est la vulnérabilité — bâtiments non parasismiques, pauvreté, secours insuffisants — qui a transformé le séisme haïtien en hécatombe. De même, au lac Nyos (1986), l'émission limnique, aléa rare, n'a été meurtrière (\num{1746} morts) que parce que des villages se trouvaient dans les vallées de déversement du nuage de \ce{CO2}. À Douala, la même pluie diluvienne dévaste les quartiers bas aux caniveaux obstrués et épargne les quartiers drainés : l'aléa est identique, les dégâts diffèrent selon la vulnérabilité. À l'inverse, les éruptions fréquentes du Mont Cameroun (1999, 2000) n'ont fait aucune victime grâce à la faible densité de population sur les zones de coulées et à la surveillance volcanique : un aléa réel peut ne produire aucune catastrophe si la vulnérabilité est faible.
\item \textbf{Conclusion.} Puisqu'on ne peut supprimer la plupart des aléas naturels, la gestion des catastrophes au Cameroun doit porter sur la vulnérabilité : aménagement du territoire (interdire les constructions dans les vallées du Nyos et les zones inondables de Douala), normes de construction, éducation des populations aux gestes qui sauvent, plans d'urgence et surveillance scientifique. Réduire la vulnérabilité, c'est réduire la catastrophe avant même qu'elle ne survienne.
\end{enumerate}
\textbf{Barème indicatif (sur 20) :} définitions et relation aléa--vulnérabilité--risque : 4 pts ; comparaison argumentée d'au moins deux situations à aléa semblable et bilans différents : 8 pts ; mobilisation d'exemples camerounais : 4 pts ; conclusion dégageant les implications pratiques : 4 pts.

\textbf{Points de vigilance :} une question de synthèse exige une démonstration construite (introduction, définitions, comparaisons chiffrées, conclusion), et non une liste d'exemples ; les comparaisons doivent être explicites (même aléa, bilans différents) ; les chiffres (magnitude 7,0 ; \num{200000} morts ; \num{1746} victimes) doivent être exacts ; la conclusion doit retomber sur des mesures concrètes applicables au Cameroun.
\end{corrige}

\newpage

\coursec{Fiche bilan}

\begin{popfigure}
\centering
\begin{tikzpicture}[mindmap, grow cyclic,
  every node/.style={concept, font=\small\bfseries},
  root concept/.append style={concept color=popblue, font=\large\bfseries, minimum size=3.2cm},
  level 1 concept/.append style={level distance=4.2cm, sibling angle=60, font=\footnotesize\bfseries, minimum size=2.2cm},
  level 2 concept/.append style={level distance=2.6cm, font=\scriptsize, minimum size=1.6cm}]
\node[root concept]{Les catastrophes}
  child[concept color=poporange]{ node {Définition}
    child { node {Événement soudain}
    }
    child { node {Bilan : morts, dégâts}
    }
  }
  child[concept color=popgreen]{ node {Classification}
    child { node {Origine naturelle}
    }
    child { node {Origine humaine}
    }
  }
  child[concept color=popblue]{ node {Catastrophes naturelles}
    child { node {Séismes, volcans}
    }
    child { node {Inondations, sécheresses}
    }
    child { node {Épidémies}
    }
  }
  child[concept color=poppurple]{ node {Catastrophes humaines}
    child { node {Accidents, incendies}
    }
    child { node {Guerres, pollutions}
    }
  }
  child[concept color=popgold]{ node {Gestion}
    child { node {Prévision}
    }
    child { node {Prédiction}
    }
    child { node {Prévention}
    }
  }
  child[concept color=poppink]{ node {Gestes qui sauvent}
    child { node {Alerter, protéger}
    }
    child { node {Secourir}
    }
  };
\end{tikzpicture}
\legende{Carte mentale de la leçon : les six branches essentielles à maîtriser — définition, classification, catastrophes naturelles, catastrophes humaines, gestion (prévision, prédiction, prévention) et gestes qui sauvent.}
\end{popfigure}

\begin{bilan}
\textbf{Définitions clés}
\begin{itemize}
  \item \cle{Catastrophe} : événement brutal, d'origine naturelle ou humaine, provoquant de nombreuses pertes en vies humaines et des dégâts matériels et environnementaux importants, dépassant les capacités locales de réaction.
  \item \cle{Aléa} : phénomène dangereux susceptible de se produire (séisme, inondation\ldots).
  \item \cle{Vulnérabilité} : degré de fragilité des populations et des biens face à l'aléa. \textbf{Risque = aléa $\times$ vulnérabilité}.
  \item \cle{Prévision} : annoncer qu'un phénomène peut se produire dans une zone donnée, sans date précise.
  \item \cle{Prédiction} : annoncer le lieu, la date approximative et l'intensité du phénomène.
  \item \cle{Prévention} : ensemble des mesures prises \emph{avant} la catastrophe pour en limiter les effets (éducation, normes de construction, plans d'évacuation).
\end{itemize}

\textbf{Classification à connaître par c\oe ur}
\begin{center}
\begin{tabularx}{\linewidth}{|l|X|X|}
\hline
\rowcolor{popblueL} \textbf{Origine} & \textbf{Exemples} & \textbf{Exemples camerounais} \\
\hline
Naturelle & Séismes, éruptions volcaniques, tsunamis, inondations, sécheresses, glissements de terrain, épidémies & Éruption limnique du lac Nyos (1986, environ \num{1700} morts), inondations de l'Extrême-Nord, choléra \\
\hline
Humaine (technologique ou sociale) & Accidents industriels et routiers, incendies, pollutions, guerres, déversements d'hydrocarbures & Accidents de la circulation sur l'axe Douala--Yaoundé, feux de brousse, pollution urbaine \\
\hline
\end{tabularx}
\end{center}

\textbf{Méthodes express}
\begin{itemize}
  \item \textbf{Classer une catastrophe} : identifier la cause première $\rightarrow$ si elle relève des processus de la nature (géologiques, climatiques, biologiques) : naturelle ; si elle résulte d'une activité ou d'une négligence humaine : humaine (certaines sont \emph{mixtes} : inondation aggravée par la déforestation).
  \item \textbf{Gérer une catastrophe} : retenir la trilogie \cle{prévision -- prédiction -- prévention}, puis les phases : avant (préparation), pendant (conduite à tenir, gestes qui sauvent), après (secours, reconstruction).
  \item \textbf{Gestes qui sauvent} : protéger -- alerter (SAMU : 119, sapeurs-pompiers : 118) -- secourir (position latérale de sécurité, compression d'une hémorragie, massage cardiaque).
\end{itemize}
\end{bilan}

\begin{aretenir}[Checklist avant le Bac]
\begin{itemize}
  \item $\square$ Je sais définir une catastrophe et citer ses caractéristiques (soudaineté, ampleur, pertes).
  \item $\square$ Je distingue aléa, vulnérabilité et risque.
  \item $\square$ Je sais classer une catastrophe selon son origine : naturelle ou humaine.
  \item $\square$ Je cite au moins trois catastrophes naturelles avec un exemple camerounais (lac Nyos, 1986).
  \item $\square$ Je cite au moins trois catastrophes d'origine humaine.
  \item $\square$ Je sais que certaines catastrophes sont mixtes (cause naturelle aggravée par l'activité humaine).
  \item $\square$ Je distingue clairement prévision, prédiction et prévention.
  \item $\square$ Je connais les trois phases de la gestion d'une catastrophe : avant, pendant, après.
  \item $\square$ Je connais la conduite à tenir en cas de séisme, d'inondation et d'incendie.
  \item $\square$ Je maîtrise la chaîne « protéger -- alerter -- secourir » et les numéros d'urgence (118, 119).
  \item $\square$ Je sais décrire la position latérale de sécurité et l'arrêt d'une hémorragie.
  \item $\square$ Je sais élaborer un message de sensibilisation clair sur la gestion des catastrophes.
\end{itemize}
\end{aretenir}

\findecours

\end{document}
